% concatenated from VROG_icml2025_camera_ready.tex, math_commands.tex, main_paper.tex, main_refs.tex, supp_docs.tex, supp_refs.tex
%%%%%%%% ICML 2025 EXAMPLE LATEX SUBMISSION FILE %%%%%%%%%%%%%%%%%
\documentclass{article}

% Recommended, but optional, packages for figures and better typesetting:
\usepackage{microtype}
\usepackage{graphicx}
\usepackage{subfigure}
\usepackage{booktabs} % for professional tables

% hyperref makes hyperlinks in the resulting PDF.
% If your build breaks (sometimes temporarily if a hyperlink spans a page)
% please comment out the following usepackage line and replace
% \usepackage{icml2025} with \usepackage[nohyperref]{icml2025} above.
\usepackage{hyperref}
\usepackage{amssymb}
\usepackage{pifont}

% Attempt to make hyperref and algorithmic work together better:
\newcommand{\theHalgorithm}{\arabic{algorithm}}

% Use the following line for the initial blind version submitted for review:
%\usepackage{icml2025}

% If accepted, instead use the following line for the camera-ready submission:
\usepackage[accepted]{icml2025}

% For theorems and such
\usepackage{amsmath}
\usepackage{amssymb}
\usepackage{mathtools}
\usepackage{amsthm}
\usepackage{makecell}

% if you use cleveref..
\usepackage[capitalize,noabbrev]{cleveref}

%%%%% NEW MATH DEFINITIONS %%%%%

\usepackage{amsmath,amsfonts,bm}
\usepackage{amsthm}
\usepackage{amssymb}
\usepackage{paralist}
\usepackage{graphicx}
\usepackage{algorithm}
\usepackage{algorithmic}
\usepackage{caption}

% Mark sections of captions for referring to divisions of figures
\newcommand{\figleft}{{\em (Left)}}
\newcommand{\figcenter}{{\em (Center)}}
\newcommand{\figright}{{\em (Right)}}
\newcommand{\figtop}{{\em (Top)}}
\newcommand{\figbottom}{{\em (Bottom)}}
\newcommand{\captiona}{{\em (a)}}
\newcommand{\captionb}{{\em (b)}}
\newcommand{\captionc}{{\em (c)}}
\newcommand{\captiond}{{\em (d)}}

% Highlight a newly defined term
\newcommand{\newterm}[1]{{\bf #1}}


% Figure reference, lower-case.
\def\figref#1{figure~\ref{#1}}
% Figure reference, capital. For start of sentence
\def\Figref#1{Figure~\ref{#1}}
\def\twofigref#1#2{figures \ref{#1} and \ref{#2}}
\def\quadfigref#1#2#3#4{figures \ref{#1}, \ref{#2}, \ref{#3} and \ref{#4}}
% Section reference, lower-case.
\def\secref#1{section~\ref{#1}}
% Section reference, capital.
\def\Secref#1{Section~\ref{#1}}
% Reference to two sections.
\def\twosecrefs#1#2{sections \ref{#1} and \ref{#2}}
% Reference to three sections.
\def\secrefs#1#2#3{sections \ref{#1}, \ref{#2} and \ref{#3}}
% Reference to an equation, lower-case.
%\def\eqref#1{equation~\ref{#1}}
% Reference to an equation, upper case
\def\Eqref#1{Equation~\ref{#1}}
% A raw reference to an equation---avoid using if possible
\def\plaineqref#1{\ref{#1}}
% Reference to a chapter, lower-case.
\def\chapref#1{chapter~\ref{#1}}
% Reference to an equation, upper case.
\def\Chapref#1{Chapter~\ref{#1}}
% Reference to a range of chapters
\def\rangechapref#1#2{chapters\ref{#1}--\ref{#2}}
% Reference to an algorithm, lower-case.
\def\algref#1{algorithm~\ref{#1}}
% Reference to an algorithm, upper case.
\def\Algref#1{Algorithm~\ref{#1}}
\def\twoalgref#1#2{algorithms \ref{#1} and \ref{#2}}
\def\Twoalgref#1#2{Algorithms \ref{#1} and \ref{#2}}
% Reference to a part, lower case
\def\partref#1{part~\ref{#1}}
% Reference to a part, upper case
\def\Partref#1{Part~\ref{#1}}
\def\twopartref#1#2{parts \ref{#1} and \ref{#2}}

\def\ceil#1{\lceil #1 \rceil}
\def\floor#1{\lfloor #1 \rfloor}
\def\1{\bm{1}}
\newcommand{\train}{\mathcal{D}}
\newcommand{\valid}{\mathcal{D_{\mathrm{valid}}}}
\newcommand{\test}{\mathcal{D_{\mathrm{test}}}}

\def\eps{{\epsilon}}


% Random variables
\def\reta{{\textnormal{$\eta$}}}
\def\ra{{\textnormal{a}}}
\def\rb{{\textnormal{b}}}
\def\rc{{\textnormal{c}}}
\def\rd{{\textnormal{d}}}
\def\re{{\textnormal{e}}}
\def\rf{{\textnormal{f}}}
\def\rg{{\textnormal{g}}}
\def\rh{{\textnormal{h}}}
\def\ri{{\textnormal{i}}}
\def\rj{{\textnormal{j}}}
\def\rk{{\textnormal{k}}}
\def\rl{{\textnormal{l}}}
% rm is already a command, just don't name any random variables m
\def\rn{{\textnormal{n}}}
\def\ro{{\textnormal{o}}}
\def\rp{{\textnormal{p}}}
\def\rq{{\textnormal{q}}}
\def\rr{{\textnormal{r}}}
\def\rs{{\textnormal{s}}}
\def\rt{{\textnormal{t}}}
\def\ru{{\textnormal{u}}}
\def\rv{{\textnormal{v}}}
\def\rw{{\textnormal{w}}}
\def\rx{{\textnormal{x}}}
\def\ry{{\textnormal{y}}}
\def\rz{{\textnormal{z}}}

% Random vectors
\def\rvepsilon{{\mathbf{\epsilon}}}
\def\rvtheta{{\mathbf{\theta}}}
\def\rva{{\mathbf{a}}}
\def\rvb{{\mathbf{b}}}
\def\rvc{{\mathbf{c}}}
\def\rvd{{\mathbf{d}}}
\def\rve{{\mathbf{e}}}
\def\rvf{{\mathbf{f}}}
\def\rvg{{\mathbf{g}}}
\def\rvh{{\mathbf{h}}}
\def\rvu{{\mathbf{i}}}
\def\rvj{{\mathbf{j}}}
\def\rvk{{\mathbf{k}}}
\def\rvl{{\mathbf{l}}}
\def\rvm{{\mathbf{m}}}
\def\rvn{{\mathbf{n}}}
\def\rvo{{\mathbf{o}}}
\def\rvp{{\mathbf{p}}}
\def\rvq{{\mathbf{q}}}
\def\rvr{{\mathbf{r}}}
\def\rvs{{\mathbf{s}}}
\def\rvt{{\mathbf{t}}}
\def\rvu{{\mathbf{u}}}
\def\rvv{{\mathbf{v}}}
\def\rvw{{\mathbf{w}}}
\def\rvx{{\mathbf{x}}}
\def\rvy{{\mathbf{y}}}
\def\rvz{{\mathbf{z}}}

% Elements of random vectors
\def\erva{{\textnormal{a}}}
\def\ervb{{\textnormal{b}}}
\def\ervc{{\textnormal{c}}}
\def\ervd{{\textnormal{d}}}
\def\erve{{\textnormal{e}}}
\def\ervf{{\textnormal{f}}}
\def\ervg{{\textnormal{g}}}
\def\ervh{{\textnormal{h}}}
\def\ervi{{\textnormal{i}}}
\def\ervj{{\textnormal{j}}}
\def\ervk{{\textnormal{k}}}
\def\ervl{{\textnormal{l}}}
\def\ervm{{\textnormal{m}}}
\def\ervn{{\textnormal{n}}}
\def\ervo{{\textnormal{o}}}
\def\ervp{{\textnormal{p}}}
\def\ervq{{\textnormal{q}}}
\def\ervr{{\textnormal{r}}}
\def\ervs{{\textnormal{s}}}
\def\ervt{{\textnormal{t}}}
\def\ervu{{\textnormal{u}}}
\def\ervv{{\textnormal{v}}}
\def\ervw{{\textnormal{w}}}
\def\ervx{{\textnormal{x}}}
\def\ervy{{\textnormal{y}}}
\def\ervz{{\textnormal{z}}}

% Random matrices
\def\rmA{{\mathbf{A}}}
\def\rmB{{\mathbf{B}}}
\def\rmC{{\mathbf{C}}}
\def\rmD{{\mathbf{D}}}
\def\rmE{{\mathbf{E}}}
\def\rmF{{\mathbf{F}}}
\def\rmG{{\mathbf{G}}}
\def\rmH{{\mathbf{H}}}
\def\rmI{{\mathbf{I}}}
\def\rmJ{{\mathbf{J}}}
\def\rmK{{\mathbf{K}}}
\def\rmL{{\mathbf{L}}}
\def\rmM{{\mathbf{M}}}
\def\rmN{{\mathbf{N}}}
\def\rmO{{\mathbf{O}}}
\def\rmP{{\mathbf{P}}}
\def\rmQ{{\mathbf{Q}}}
\def\rmR{{\mathbf{R}}}
\def\rmS{{\mathbf{S}}}
\def\rmT{{\mathbf{T}}}
\def\rmU{{\mathbf{U}}}
\def\rmV{{\mathbf{V}}}
\def\rmW{{\mathbf{W}}}
\def\rmX{{\mathbf{X}}}
\def\rmY{{\mathbf{Y}}}
\def\rmZ{{\mathbf{Z}}}

% Elements of random matrices
\def\ermA{{\textnormal{A}}}
\def\ermB{{\textnormal{B}}}
\def\ermC{{\textnormal{C}}}
\def\ermD{{\textnormal{D}}}
\def\ermE{{\textnormal{E}}}
\def\ermF{{\textnormal{F}}}
\def\ermG{{\textnormal{G}}}
\def\ermH{{\textnormal{H}}}
\def\ermI{{\textnormal{I}}}
\def\ermJ{{\textnormal{J}}}
\def\ermK{{\textnormal{K}}}
\def\ermL{{\textnormal{L}}}
\def\ermM{{\textnormal{M}}}
\def\ermN{{\textnormal{N}}}
\def\ermO{{\textnormal{O}}}
\def\ermP{{\textnormal{P}}}
\def\ermQ{{\textnormal{Q}}}
\def\ermR{{\textnormal{R}}}
\def\ermS{{\textnormal{S}}}
\def\ermT{{\textnormal{T}}}
\def\ermU{{\textnormal{U}}}
\def\ermV{{\textnormal{V}}}
\def\ermW{{\textnormal{W}}}
\def\ermX{{\textnormal{X}}}
\def\ermY{{\textnormal{Y}}}
\def\ermZ{{\textnormal{Z}}}

% Vectors
\def\vzero{{\bm{0}}}
\def\vone{{\bm{1}}}
\def\vmu{{\bm{\mu}}}
\def\vtheta{{\bm{\theta}}}
\def\va{{\bm{a}}}
\def\vb{{\bm{b}}}
\def\vc{{\bm{c}}}
\def\vd{{\bm{d}}}
\def\ve{{\bm{e}}}
\def\vf{{\bm{f}}}
\def\vg{{\bm{g}}}
\def\vh{{\bm{h}}}
\def\vi{{\bm{i}}}
\def\vj{{\bm{j}}}
\def\vk{{\bm{k}}}
\def\vl{{\bm{l}}}
\def\vm{{\bm{m}}}
\def\vn{{\bm{n}}}
\def\vo{{\bm{o}}}
\def\vp{{\bm{p}}}
\def\vq{{\bm{q}}}
\def\vr{{\bm{r}}}
\def\vs{{\bm{s}}}
\def\vt{{\bm{t}}}
\def\vu{{\bm{u}}}
\def\vv{{\bm{v}}}
\def\vw{{\bm{w}}}
\def\vx{{\bm{x}}}
\def\vy{{\bm{y}}}
\def\vz{{\bm{z}}}

% Elements of vectors
\def\evalpha{{\alpha}}
\def\evbeta{{\beta}}
\def\evepsilon{{\epsilon}}
\def\evlambda{{\lambda}}
\def\evomega{{\omega}}
\def\evmu{{\mu}}
\def\evpsi{{\psi}}
\def\evsigma{{\sigma}}
\def\evtheta{{\theta}}
\def\eva{{a}}
\def\evb{{b}}
\def\evc{{c}}
\def\evd{{d}}
\def\eve{{e}}
\def\evf{{f}}
\def\evg{{g}}
\def\evh{{h}}
\def\evi{{i}}
\def\evj{{j}}
\def\evk{{k}}
\def\evl{{l}}
\def\evm{{m}}
\def\evn{{n}}
\def\evo{{o}}
\def\evp{{p}}
\def\evq{{q}}
\def\evr{{r}}
\def\evs{{s}}
\def\evt{{t}}
\def\evu{{u}}
\def\evv{{v}}
\def\evw{{w}}
\def\evx{{x}}
\def\evy{{y}}
\def\evz{{z}}

% Matrix
\def\mA{{\bm{A}}}
\def\mB{{\bm{B}}}
\def\mC{{\bm{C}}}
\def\mD{{\bm{D}}}
\def\mE{{\bm{E}}}
\def\mF{{\bm{F}}}
\def\mG{{\bm{G}}}
\def\mH{{\bm{H}}}
\def\mI{{\bm{I}}}
\def\mJ{{\bm{J}}}
\def\mK{{\bm{K}}}
\def\mL{{\bm{L}}}
\def\mM{{\bm{M}}}
\def\mN{{\bm{N}}}
\def\mO{{\bm{O}}}
\def\mP{{\bm{P}}}
\def\mQ{{\bm{Q}}}
\def\mR{{\bm{R}}}
\def\mS{{\bm{S}}}
\def\mT{{\bm{T}}}
\def\mU{{\bm{U}}}
\def\mV{{\bm{V}}}
\def\mW{{\bm{W}}}
\def\mX{{\bm{X}}}
\def\mY{{\bm{Y}}}
\def\mZ{{\bm{Z}}}
\def\mBeta{{\bm{\beta}}}
\def\mPhi{{\bm{\Phi}}}
\def\mLambda{{\bm{\Lambda}}}
\def\mSigma{{\bm{\Sigma}}}

% Tensor
\DeclareMathAlphabet{\mathsfit}{\encodingdefault}{\sfdefault}{m}{sl}
\SetMathAlphabet{\mathsfit}{bold}{\encodingdefault}{\sfdefault}{bx}{n}
\newcommand{\tens}[1]{\bm{\mathsfit{#1}}}
\def\tA{{\tens{A}}}
\def\tB{{\tens{B}}}
\def\tC{{\tens{C}}}
\def\tD{{\tens{D}}}
\def\tE{{\tens{E}}}
\def\tF{{\tens{F}}}
\def\tG{{\tens{G}}}
\def\tH{{\tens{H}}}
\def\tI{{\tens{I}}}
\def\tJ{{\tens{J}}}
\def\tK{{\tens{K}}}
\def\tL{{\tens{L}}}
\def\tM{{\tens{M}}}
\def\tN{{\tens{N}}}
\def\tO{{\tens{O}}}
\def\tP{{\tens{P}}}
\def\tQ{{\tens{Q}}}
\def\tR{{\tens{R}}}
\def\tS{{\tens{S}}}
\def\tT{{\tens{T}}}
\def\tU{{\tens{U}}}
\def\tV{{\tens{V}}}
\def\tW{{\tens{W}}}
\def\tX{{\tens{X}}}
\def\tY{{\tens{Y}}}
\def\tZ{{\tens{Z}}}


% Graph
\def\gA{{\mathcal{A}}}
\def\gB{{\mathcal{B}}}
\def\gC{{\mathcal{C}}}
\def\gD{{\mathcal{D}}}
\def\gE{{\mathcal{E}}}
\def\gF{{\mathcal{F}}}
\def\gG{{\mathcal{G}}}
\def\gH{{\mathcal{H}}}
\def\gI{{\mathcal{I}}}
\def\gJ{{\mathcal{J}}}
\def\gK{{\mathcal{K}}}
\def\gL{{\mathcal{L}}}
\def\gM{{\mathcal{M}}}
\def\gN{{\mathcal{N}}}
\def\gO{{\mathcal{O}}}
\def\gP{{\mathcal{P}}}
\def\gQ{{\mathcal{Q}}}
\def\gR{{\mathcal{R}}}
\def\gS{{\mathcal{S}}}
\def\gT{{\mathcal{T}}}
\def\gU{{\mathcal{U}}}
\def\gV{{\mathcal{V}}}
\def\gW{{\mathcal{W}}}
\def\gX{{\mathcal{X}}}
\def\gY{{\mathcal{Y}}}
\def\gZ{{\mathcal{Z}}}

% Sets
\def\sA{{\mathbb{A}}}
\def\sB{{\mathbb{B}}}
\def\sC{{\mathbb{C}}}
\def\sD{{\mathbb{D}}}
% Don't use a set called E, because this would be the same as our symbol
% for expectation.
\def\sF{{\mathbb{F}}}
\def\sG{{\mathbb{G}}}
\def\sH{{\mathbb{H}}}
\def\sI{{\mathbb{I}}}
\def\sJ{{\mathbb{J}}}
\def\sK{{\mathbb{K}}}
\def\sL{{\mathbb{L}}}
\def\sM{{\mathbb{M}}}
\def\sN{{\mathbb{N}}}
\def\sO{{\mathbb{O}}}
\def\sP{{\mathbb{P}}}
\def\sQ{{\mathbb{Q}}}
\def\sR{{\mathbb{R}}}
\def\sS{{\mathbb{S}}}
\def\sT{{\mathbb{T}}}
\def\sU{{\mathbb{U}}}
\def\sV{{\mathbb{V}}}
\def\sW{{\mathbb{W}}}
\def\sX{{\mathbb{X}}}
\def\sY{{\mathbb{Y}}}
\def\sZ{{\mathbb{Z}}}

% Entries of a matrix
\def\emLambda{{\Lambda}}
\def\emA{{A}}
\def\emB{{B}}
\def\emC{{C}}
\def\emD{{D}}
\def\emE{{E}}
\def\emF{{F}}
\def\emG{{G}}
\def\emH{{H}}
\def\emI{{I}}
\def\emJ{{J}}
\def\emK{{K}}
\def\emL{{L}}
\def\emM{{M}}
\def\emN{{N}}
\def\emO{{O}}
\def\emP{{P}}
\def\emQ{{Q}}
\def\emR{{R}}
\def\emS{{S}}
\def\emT{{T}}
\def\emU{{U}}
\def\emV{{V}}
\def\emW{{W}}
\def\emX{{X}}
\def\emY{{Y}}
\def\emZ{{Z}}
\def\emSigma{{\Sigma}}

% entries of a tensor
% Same font as tensor, without \bm wrapper
\newcommand{\etens}[1]{\mathsfit{#1}}
\def\etLambda{{\etens{\Lambda}}}
\def\etA{{\etens{A}}}
\def\etB{{\etens{B}}}
\def\etC{{\etens{C}}}
\def\etD{{\etens{D}}}
\def\etE{{\etens{E}}}
\def\etF{{\etens{F}}}
\def\etG{{\etens{G}}}
\def\etH{{\etens{H}}}
\def\etI{{\etens{I}}}
\def\etJ{{\etens{J}}}
\def\etK{{\etens{K}}}
\def\etL{{\etens{L}}}
\def\etM{{\etens{M}}}
\def\etN{{\etens{N}}}
\def\etO{{\etens{O}}}
\def\etP{{\etens{P}}}
\def\etQ{{\etens{Q}}}
\def\etR{{\etens{R}}}
\def\etS{{\etens{S}}}
\def\etT{{\etens{T}}}
\def\etU{{\etens{U}}}
\def\etV{{\etens{V}}}
\def\etW{{\etens{W}}}
\def\etX{{\etens{X}}}
\def\etY{{\etens{Y}}}
\def\etZ{{\etens{Z}}}

% The true underlying data generating distribution
\newcommand{\pdata}{p_{\rm{data}}}
% The empirical distribution defined by the training set
\newcommand{\ptrain}{\hat{p}_{\rm{data}}}
\newcommand{\Ptrain}{\hat{P}_{\rm{data}}}
% The model distribution
\newcommand{\pmodel}{p_{\rm{model}}}
\newcommand{\Pmodel}{P_{\rm{model}}}
\newcommand{\ptildemodel}{\tilde{p}_{\rm{model}}}
% Stochastic autoencoder distributions
\newcommand{\pencode}{p_{\rm{encoder}}}
\newcommand{\pdecode}{p_{\rm{decoder}}}
\newcommand{\precons}{p_{\rm{reconstruct}}}

\newcommand{\laplace}{\mathrm{Laplace}} % Laplace distribution

\newcommand{\E}{\mathbb{E}}
\newcommand{\Ls}{\mathcal{L}}
\newcommand{\R}{\mathbb{R}}
\newcommand{\emp}{\tilde{p}}
\newcommand{\lr}{\alpha}
\newcommand{\reg}{\lambda}
\newcommand{\rect}{\mathrm{rectifier}}
\newcommand{\softmax}{\mathrm{softmax}}
\newcommand{\sigmoid}{\sigma}
\newcommand{\softplus}{\zeta}
\newcommand{\KL}{D_{\mathrm{KL}}}
\newcommand{\Var}{\mathrm{Var}}
\newcommand{\standarderror}{\mathrm{SE}}
\newcommand{\Cov}{\mathrm{Cov}}
% Wolfram Mathworld says $L^2$ is for function spaces and $\ell^2$ is for vectors
% But then they seem to use $L^2$ for vectors throughout the site, and so does
% wikipedia.
\newcommand{\normlzero}{L^0}
\newcommand{\normlone}{L^1}
\newcommand{\normltwo}{L^2}
\newcommand{\normlp}{L^p}
\newcommand{\normmax}{L^\infty}

\newcommand{\parents}{Pa} % See usage in notation.tex. Chosen to match Daphne's book.

\DeclareMathOperator*{\argmax}{arg\,max}
\DeclareMathOperator*{\argmin}{arg\,min}

\DeclareMathOperator{\sign}{sign}
\DeclareMathOperator{\Tr}{Tr}
\let\ab\allowbreak


%%%%%%%%%%%%%%%%%%%%%%%%%%%%%%%%%%%%%%%%%%%%%%%%%%%%%%%%
%%%%%% USER DEFINED COMMANDS
%%%%%%%%%%%%%%%%%%%%%%%%%%%%%%%%%%%%%%%%%%%%%%%%%%%%%%%%
%+ Renewcommand symbols
\newcommand{\Id}{\mathbb{I}}
\newcommand{\Hil}{\mathcal{H}}
\newcommand{\Rext}{\R\cup\{+\infty\}}
\newcommand{\abs}[1]{\left\vert#1\right\vert}
\newcommand{\set}[1]{\left\{#1\right\}}
\newcommand{\sets}[1]{\{#1\}}
\newcommand{\norm}[1]{\left\Vert#1\right\Vert}
\newcommand{\norms}[1]{\Vert#1\Vert}
\newcommand{\iprod}[1]{\left\langle #1\right\rangle}
\newcommand{\iprods}[1]{\langle #1\rangle}

\newcommand{\epi}[1]{\mathrm{epi}\left(#1\right)}
\newcommand{\Eproof}{\hfill $\square$}
\newcommand{\prox}{\mathrm{prox}}
\newcommand{\proj}{\mathrm{proj}}
\newcommand{\dom}[1]{\mathrm{dom}(#1)}
\newcommand{\zero}[1]{{\boldsymbol{0}}}
\newcommand{\Exps}[2]{\mathbb{E}_{#1}\left[#2\right]}
\newcommand{\Expsk}[2]{\mathbb{E}_{#1}\big[#2\big]}
\newcommand{\Expk}[1]{\mathbb{E}\big[#1\big]}
\newcommand{\Expsb}[2]{\mathbb{E}_{#1}\big[#2\big]}
\newcommand{\Exp}[1]{\mathbb{E}\left[#1\right]}
\newcommand{\cExp}[2]{\mathbb{E}\left[#1\mid #2\right]}
\newcommand{\Prob}[1]{\mathbf{Prob}\left(#1\right)}
\newcommand{\cProb}[2]{\mathbf{Prob}\left(#1 \mid #2 \right)}
\newcommand{\zer}[1]{\mathrm{zer}(#1)}
\newcommand{\gra}[1]{\mathrm{gra}(#1)}
\newcommand{\range}[1]{\mathrm{range}(#1)}
\newcommand{\mcal}[1]{\mathcal{#1}}

\newcommand{\xopt}{x^{\star}}
\newcommand{\xbar}{\bar{x}}
\newcommand{\xast}{x^{\ast}}
\newcommand{\yast}{y^{\ast}}
%\newcommand{\xb}{x}
%\newcommand{\yb}{y}
%\newcommand{\eb}{\mathbf{e}}
%\newcommand{\zb}{z}
%\newcommand{\wb}{w}
%\newcommand{\ub}{u}
%\newcommand{\vb}{v}
%\newcommand{\lb}{\mathbf{l}}
%\renewcommand{\sb}{s}
%\newcommand{\rb}{r}
%\newcommand{\qb}{q}
%\newcommand{\ab}{a}
%\newcommand{\bb}{b}
%\newcommand{\cb}{c}
%\newcommand{\db}{d}
%\newcommand{\hb}{h}
%\newcommand{\nb}{n}
%\newcommand{\pb}{p}
\newcommand{\Ab}{\mathbf{A}}
\newcommand{\Gb}{\mathbf{G}}
\newcommand{\Hb}{\mathbf{H}}
\newcommand{\Kb}{\mathbf{K}}
\newcommand{\Ac}{\mathcal{A}}
\newcommand{\Hc}{\mathcal{H}}
\newcommand{\Bb}{\mathbf{B}}
\newcommand{\Db}{\mathbf{D}}
\newcommand{\Pb}{\mathbf{P}}
\newcommand{\Pbb}{\mathbb{P}}
\newcommand{\Qb}{\mathbf{Q}}
\newcommand{\Bc}{\mathcal{B}}
\newcommand{\Cb}{\mathbf{C}}
\newcommand{\Ub}{\mathbf{U}}
\newcommand{\Vb}{\mathbf{V}}
\newcommand{\Xb}{\mathbf{X}}
\newcommand{\Yb}{\mathbf{Y}}
\newcommand{\Zb}{\mathbf{Z}}
\newcommand{\Xc}{\mathcal{X}}
\newcommand{\Yc}{\mathcal{Y}}
\newcommand{\Zc}{\mathcal{Z}}
\newcommand{\Sc}{\mathcal{S}}
\newcommand{\Sdp}{\mathbb{S}}
\newcommand{\Dc}{\mathcal{D}}
\newcommand{\Lc}{\mathcal{L}}
\newcommand{\Qc}{\mathcal{Q}}
\newcommand{\Gc}{\mathcal{G}}
\newcommand{\Mc}{\mathcal{M}}
\newcommand{\Tc}{\mathcal{T}}
\newcommand{\Lb}{\mathbf{L}}
\newcommand{\Fb}{\mathbf{F}}
\newcommand{\Ib}{\mathcal{I}}
\newcommand{\Fc}{\mathcal{F}}
\newcommand{\Uc}{\mathcal{U}}
\newcommand{\Vc}{\mathcal{V}}
\newcommand{\Cc}{\mathcal{C}}
\newcommand{\Ec}{\mathcal{E}}
\newcommand{\Nc}{\mathcal{N}}
\newcommand{\Pc}{\mathcal{P}}
\newcommand{\Kc}{\mathcal{K}}
\newcommand{\Wc}{\mathcal{W}}
\newcommand{\Xopt}{\mathcal{X}^{\star}}
\newcommand{\Yopt}{\mathcal{Y}^{\star}}
\newcommand{\Lag}{\widetilde{\mathcal{L}}}

\newcommand{\BigO}[1]{\mathcal{O}\left(#1\right)}
\newcommand{\BigOs}[1]{\mathcal{O}\big(#1\big)}
\newcommand{\SmallOs}[1]{o\big(#1\big)}
\newcommand{\SmallO}[1]{o\left(#1\right)}
\newcommand{\Ball}[2]{\mathbb{B}_{#1}\left(#2\right)}
\newcommand{\BigOc}[1]{{\color{blue}\mathcal{O}\left(#1\right)}}
\newcommand{\SmallOc}[1]{{\color{blue}o\left(#1\right)}}

%% For text colors
\newcommand{\mred}[1]{\textcolor{red}{#1}}
\newcommand{\mblue}[1]{\textcolor{blue}{#1}}

% Editing commands
\newcommand{\mytbi}[1]{\textbf{\textit{#1}}}
\newcommand{\mytb}[1]{\textbf{#1}}
\newcommand{\myti}[1]{\textit{#1}}
\newcommand{\mbf}[1]{\mathbf{#1}}
\newcommand{\mbb}[1]{\mathbb{#1}}
\newcommand{\myeqc}[1]{ {\scriptsize\textcircled{#1}} }

% Gap between section/subsection.
\newcommand{\beforesec}{\vspace{-1.5ex}}
\newcommand{\aftersec}{\vspace{-1.25ex}}
\newcommand{\beforesubsec}{\vspace{-1.5ex}}
\newcommand{\aftersubsec}{\vspace{-1.5ex}}
\newcommand{\beforesubsubsec}{\vspace{-1.5ex}}
\newcommand{\aftersubsubsec}{\vspace{-1.5ex}}
\newcommand{\beforepara}{\vspace{-1.5ex}}
%\newcommand{\beforesec}{\vspace{0ex}}
%\newcommand{\aftersec}{\vspace{0ex}}
%\newcommand{\beforesubsec}{\vspace{0ex}}
%\newcommand{\aftersubsec}{\vspace{0ex}}
%\newcommand{\beforesubsubsec}{\vspace{0ex}}
%\newcommand{\aftersubsubsec}{\vspace{0ex}}
%\newcommand{\beforepara}{\vspace{0ex}}


\arraycolsep=0.2em
\def\arraystretch{1.5}


%\theoremstyle{plain}
%\newtheorem{theorem}{Theorem}[section]
%\newtheorem{corollary}{Corollary}[section]
%\newtheorem{lemma}{Lemma}[section]
%%\newtheorem{proposition}{Proposition}[section]
%\theoremstyle{definition}
%\newtheorem{definition}{Definition}[section]
%\newtheorem{assumption}{Assumption}[section]
%\theoremstyle{remark}
%\newtheorem{remark}{Remark}[section]

\newcommand{\myblue}[1]{{\color{blue}{#1}}}
\newcommand{\myred}[1]{{\color{red}{#1}}}
\newcommand{\mymgt}[1]{{\color{magenta}{#1}}}
\newcommand{\mygreen}[1]{{\color{green}{#1}}}
\newcommand{\qcmt}[1]{{\color{red}\textit{#1}}}
\newcommand{\scmt}[1]{{\color{red}{#1}}}
%\newcommand{\myeqc}[1]{ {\tiny\textcircled{#1}} }
\newcommand{\Expn}[1]{\mathbb{E}\big[#1\big]}

%%%%%%%%%%%%%%%%%%%%%%%%%%%%%%%%
% THEOREMS
%%%%%%%%%%%%%%%%%%%%%%%%%%%%%%%%
\theoremstyle{plain}
\newtheorem{theorem}{Theorem}[section]
\newtheorem{proposition}[theorem]{Proposition}
\newtheorem{lemma}[theorem]{Lemma}
\newtheorem{corollary}[theorem]{Corollary}
\theoremstyle{definition}
\newtheorem{definition}[theorem]{Definition}
\newtheorem{assumption}[theorem]{Assumption}
\theoremstyle{remark}
\newtheorem{remark}[theorem]{Remark}

% Todonotes is useful during development; simply uncomment the next line
%    and comment out the line below the next line to turn off comments
%\usepackage[disable,textsize=tiny]{todonotes}
\usepackage[textsize=tiny]{todonotes}


% The \icmltitle you define below is probably too long as a header.
% Therefore, a short form for the running title is supplied here:
\icmltitlerunning{Variance-Reduced Forward-Reflected-Backward Splitting Methods}

\begin{document}

\twocolumn[
\icmltitle{Variance-Reduced Forward-Reflected-Backward Splitting Methods for \\ Nonmonotone Generalized Equations}

% It is OKAY to include author information, even for blind
% submissions: the style file will automatically remove it for you
% unless you've provided the [accepted] option to the icml2025
% package.

% List of affiliations: The first argument should be a (short)
% identifier you will use later to specify author affiliations
% Academic affiliations should list Department, University, City, Region, Country
% Industry affiliations should list Company, City, Region, Country

% You can specify symbols, otherwise they are numbered in order.
% Ideally, you should not use this facility. Affiliations will be numbered
% in order of appearance and this is the preferred way.
\icmlsetsymbol{equal}{*}

\begin{icmlauthorlist}
\icmlauthor{Quoc Tran-Dinh}{comp}
%\icmlauthor{Firstname2 Lastname2}{equal,yyy,comp}
%\icmlauthor{Firstname3 Lastname3}{comp}
%\icmlauthor{Firstname4 Lastname4}{sch}
%\icmlauthor{Firstname5 Lastname5}{yyy}
%\icmlauthor{Firstname6 Lastname6}{sch,yyy,comp}
%\icmlauthor{Firstname7 Lastname7}{comp}
%%\icmlauthor{}{sch}
%\icmlauthor{Firstname8 Lastname8}{sch}
%\icmlauthor{Firstname8 Lastname8}{yyy,comp}
%\icmlauthor{}{sch}
%\icmlauthor{}{sch}
\end{icmlauthorlist}

\icmlaffiliation{comp}{Department of Statistics and Operations Research, The University of North Carolina at Chapel Hill, Chapel Hill, NC 27599, USA}
%\icmlaffiliation{comp}{Company Name, Location, Country}
%\icmlaffiliation{sch}{School of ZZZ, Institute of WWW, Location, Country}

\icmlcorrespondingauthor{Quoc Tran-Dinh}{quoctd@email.unc.edu}
%\icmlcorrespondingauthor{Firstname2 Lastname2}{first2.last2@www.uk}

% You may provide any keywords that you
% find helpful for describing your paper; these are used to populate
% the "keywords" metadata in the PDF but will not be shown in the document
\icmlkeywords{Machine Learning, ICML}

\vskip 0.3in
]

% this must go after the closing bracket ] following \twocolumn[ ...

% This command actually creates the footnote in the first column
% listing the affiliations and the copyright notice.
% The command takes one argument, which is text to display at the start of the footnote.
% The \icmlEqualContribution command is standard text for equal contribution.
% Remove it (just {}) if you do not need this facility.

\printAffiliationsAndNotice{}  % leave blank if no need to mention equal contribution
%\printAffiliationsAndNotice{\icmlEqualContribution} % otherwise use the standard text.

\begin{abstract}
We develop two novel stochastic variance-reduction methods to approximate solutions of a class of nonmonotone [generalized] equations.
Our algorithms leverage a new combination of ideas from the forward-reflected-backward splitting method and  a class of unbiased variance-reduced estimators. 
We construct two new stochastic estimators within this class, inspired by the well-known SVRG and SAGA estimators. 
These estimators significantly differ from existing approaches used in minimax and variational inequality problems.
By appropriately choosing parameters, both algorithms achieve state-of-the-art oracle complexity of $\mathcal{O}(n + n^{2/3} \epsilon^{-2})$ for obtaining an $\epsilon$-solution in terms of the operator residual norm for a class of nonmonotone problems, where $n$ is the number of  summands and $\epsilon$ signifies the desired accuracy. 
This complexity aligns with the best-known results in SVRG and SAGA methods for stochastic nonconvex optimization. 
We test our algorithms on some numerical examples and compare them with existing methods.
The results demonstrate promising improvements offered by the new methods compared to their competitors.
\end{abstract}


%Tasks to be done
%\begin{itemize}
%%\item On lines 140–142 of the second column, please ensure that the citation of “Diakonikolas et al.” is formatted correctly in accordance with the appropriate referencing style.
%
%%\item The authors ... results. I request that authors verify this by providing a table or a comparison with prior works in terms of complexity and assumptions. This is necessary and would enhance the novelty and clarity.
%%\item[]Thank you! Yes, we will provide a table to summarize and compare our work with existing results. We will also highlight our novelty using this summarized table as suggested.
%%\item Do the complexities of the proposed methods in this paper indeed match the best-known results?
%%\item[] To our best knowledge, our oracle complexity results are the best-known so far for non-monotone problems (i.e., covered by Assumptions 1.3-1.4) using SVRG and SAGA. We will add a table to compare as suggested.
%\item Nevertheless, based on your suggestion, we will increase the sizes of our experiments in the revision.
%\item Example 1 allows  to be non-positive semidefinite. Therefore, it violates the condition  (since ...). This means that  is non-monotone. We will clearly explain this in the revision.
%\item The work by (Cai et al 2023) relies on Halpern's fixed-point iteration, and indeed requires stronger assumptions than ours. We will conduct numerical examples to compare both methods, at least in the co-coercivity case. We are not sure if it will diverge in the non-monotone case, but there is no theoretical result to guarantee convergence for this method under Assumption 1.4.
%
%\item We agree with the reviewer that the paper is dense, and we will try to improve its presentation as suggested. It was significantly challenging for us to conduct a 10-page paper with both strong theories and experiments. Therefore, we chosed to focus on methodology and theoretical results in our submission. However, as we discussed above, additional experiments will be added to the supplementary in the revision.
%
%\item We will rephrase some inaccurate statements related to robust deep learning and generative modeling as suggested.
%
%\item I think this paper is too dense ... applications in ML. For example, what kind of non-convex problem satisfies Assumption 1.4?"
%
%\item 
%\end{itemize}

%%%%%%%%%%%%%%%%%%%%%%%%%%%%%%%
%%%% 1. Introduction.
%%%%%%%%%%%%%%%%%%%%%%%%%%%%%%%%
\vspace{-2ex}
\beforesec
\section{Introduction}\label{sec:intro}
\aftersec
\textbf{[Non]linear equations and inclusions}  are cornerstones of computational mathematics, finding applications in diverse fields like engineering, mechanics, economics, statistics, optimization, and machine learning, see, e.g., \cite{Bauschke2011,reginaset2008,Facchinei2003,phelps2009convex,ryu2022large,ryu2016primer}.
These problems, known as \textit{generalized equations} \citep{Rockafellar1997}, are equivalent to \textit{fixed-point problems}.
The recent revolution in modern machine learning and robust optimization has brought renewed interest to generalized equations and their special case: minimax problem. 
They serve as powerful tools for handling Nash's equilibria and minimax models in generative adversarial nets, adversarial training, and robust learning,  
see \cite{arjovsky2017wasserstein,goodfellow2014generative,madry2018towards,namkoong2016stochastic}.
Notably, most problems arising from these applications are nonmonotone, nonsmooth, and large-scale.
This paper develops new and simple stochastic algorithms with variance reduction for solving this class of problems, equipped with rigorous theoretical guarantees.

\beforesubsec
\subsection{Nonmonotone finite-sum generalized equations}\label{subsec:problem_statements}
\aftersubsec
The central problem we study in this paper is the following [possibly nonmonotone] \textit{generalized equation} (also known as \textit{composite inclusion}) \citep{Rockafellar1997}:
\begin{equation}\label{eq:NI}
\textrm{Find $x^{\star} \in \R^p$ such that:}~ 0 \in  Gx^{\star} + Tx^{\star},
\tag{NI}
\end{equation}
where $G : \R^p \to \R^p$ is a single-valued operator, possibly nonlinear, and $T : \R^p \rightrightarrows 2^{\R^2}$ is a multivalued mapping from $\R^p$ to $2^{\R^p}$ (the set of all subsets of $\R^p$).
In addition,  we assume that $G$ is given in the following \textit{large finite-sum}:{\!\!\!}
\vspace{-1ex}
\begin{equation}\label{eq:finite_sum_form}
Gx := \frac{1}{n}\sum_{i=1}^nG_ix,
\vspace{-0.5ex}
\end{equation}
where $G_i : \R^p\to\R^p$ are given operators for all $i \in [n] := \sets{1, 2, \cdots, n}$ and $n \gg 1$.
This structure often arises from statistical learning, [generative] machine learning, networks, distributed systems, and data science.
For simplicity of notation, we denote $\Psi := G + T$ and $\dom{\Psi} := \dom{G}\cap\dom{T}$, where $\dom{R}$ is the domain of $R$.

We highlight that the methods developed in this paper can be straightforwardly extended to tackle $Gx = \Expsk{\xi\sim\mbb{P}}{\mbf{G}(x, \xi)}$ as the expectation of a stochastic operator $\mbf{G}$ involving a random vector $\xi$ defined on a probability space $(\Omega, \mbb{P}, \Sigma)$.

\beforesubsec
\subsection{Equivalent forms and special cases}\label{subsec:special_cases}
\aftersubsec
The model \eqref{eq:NI} covers many fundamental problems in optimization and related fields, including the following ones.

\textbf{$\mathrm{(a)}$~[Non]linear equation.} If $T = 0$, then \eqref{eq:NI} reduces to the following \textit{[non]linear equation}:
\vspace{-0.5ex}
\begin{equation}\label{eq:NE}
\textrm{Find $x^{\star}\in\dom{G}$ such that:}~ Gx^{\star} = 0.
\tag{NE}
\vspace{-0.5ex}
\end{equation}
Both \eqref{eq:NI} and \eqref{eq:NE} are also called \textit{root-finding problems}.
Clearly, \eqref{eq:NE} is a special case of \eqref{eq:NI}.
However, under appropriate assumptions on $G$ and/or $T$ (e.g., using the resolvent of $T$), one can also transform \eqref{eq:NI} to \eqref{eq:NE}.
Let $\zer{\Psi} := \set{x^{\star} \in \dom{\Psi} : 0 \in \Psi{x^{\star}} }$ and $\zer{G} := \sets{x^{\star} \in \dom{G} : Gx^{\star} = 0}$  be the solution sets of \eqref{eq:NI} and \eqref{eq:NE}, respectively, which are assumed to be nonempty.

\textbf{$\mathrm{(b)}$~Variational inequality problem (VIP).}
If $T = \Nc_{\Xc}$, the normal cone of a nonempty, closed, and convex set $\Xc$ in $\R^p$, then \eqref{eq:NI} reduces to the following VIP:
\vspace{-0.5ex}
\begin{equation}\label{eq:VIP}
\hspace{-1.0ex}
\textrm{Find $x^{\star} \! \in \! \Xc$ such that:} \ \iprods{Gx^{\star}, x - x^{\star}} \geq 0, \forall x\in \Xc.
\tag{VIP}
\hspace{-0ex}
\vspace{-0.5ex}
\end{equation}
If $T = \partial{g}$, the subdifferential of a convex function $g$, then \eqref{eq:NI} reduces to a mixed VIP, denoted by MVIP.
Both VIP and MVIP cover many problems in practice, including minimax problems and Nash's equilibria, see, e.g., \cite{reginaset2008,Facchinei2003,phelps2009convex}.

\textbf{$\mathrm{(c)}$~Minimax optimization.} 
Another important special case of \eqref{eq:NI} (or MVIP) is the following minimax optimization (or saddle-point) problem, which has found various applications in machine learning and robust optimization:
\vspace{-0.5ex}
\begin{equation}\label{eq:minimax}
\min_{ u \in \R^{p_1}} \max_{v \in \R^{p_2}} \Big\{ \Lc(u, v) := \varphi(u) + \Hc(u, v) - \psi(v) \Big\},
\tag{SP}
\vspace{-0.5ex}
\end{equation}
where $\Hc : \R^{p_1}\times\R^{p_2} \to \R$ is a smooth function, and $\varphi$ and $\psi$ are proper, closed, and convex.
Let us define $x := [u, v] \in \R^p$ as the concatenation of $u$ and $v$ with $p := p_1 + p_2$, $Gx := [ \nabla_u{\Hc}(u, v), -\nabla_v{\Hc}(u, v)]$, and $Tx := [\partial{\varphi}(u), \partial{\psi}(v)]$.
Then, the optimality condition of \eqref{eq:minimax} is written in the form of \eqref{eq:NI}.
Since \eqref{eq:VIP}, and in particular, \eqref{eq:minimax} are special cases of \eqref{eq:NI}, our algorithms for \eqref{eq:NI} in the sequel  can be specified to solve these problems.

\textbf{$\mathrm{(d)}$~Fixed-point problem.} Problem \eqref{eq:NE} is equivalent to the following fixed-point problem:
\vspace{-0.75ex}
\begin{equation}\label{eq:FP_prob}
\textrm{Find $x^{\star}\in\dom{F}$ such that:}~ x^{\star} = Fx^{\star},
\tag{FP}
\vspace{-0.75ex}
\end{equation}
where $F := \Id - \lambda G$ with $\Id$ being the identity operator and $\lambda > 0$. 
Since \eqref{eq:FP_prob} is equivalent to \eqref{eq:NE}, our algorithms for \eqref{eq:NE} from this paper can also be applied to solve \eqref{eq:FP_prob}.


\beforesubsec
\subsection{Motivation}\label{subsec:motivation}
\aftersubsec
Our work is mainly motivated by the following aspects.

\vspace{-0.5ex}

%%
\textit{$\mathrm{(i)}$~Recent applications.} Both \eqref{eq:NE} and \eqref{eq:NI} cover the minimax problem \eqref{eq:minimax} as a special case.
This minimax problem, especially in nonconvex-nonconcave settings,  has recently gained its popularity as it provides a powerful tool to model applications in generative adversarial networks \citep{arjovsky2017wasserstein,goodfellow2014generative}, robust and distributionally robust optimization \citep{Ben-Tal2009,Bertsimas2011,levy2020large}, adversarial training \citep{madry2018towards}, online optimization \citep{bhatia2020online}, and reinforcement learning \citep{azar2017minimax,zhang2021multi}. 
Our work is motivated by those applications.

%%
\vspace{-0.5ex}
\textit{$\mathrm{(ii)}$~Optimality certification.} 
Existing stochastic methods often target special cases of \eqref{eq:NI} such as \eqref{eq:NE} and \eqref{eq:VIP}.
In addition, these methods frequently rely on a monotonicity assumption, which excludes many problems of current interest, e.g., \cite{alacaoglu2022beyond,alacaoglu2021stochastic,beznosikov2023stochastic,gorbunov2022stochastic,loizou2021stochastic}.
Furthermore, existing methods analyze convergence based on a [duality] \textbf{gap function} \citep{Facchinei2003} or a \textbf{restricted gap function} \citep{Nesterov2007a}.
As discussed in \cite{cai2023variance,diakonikolas2020halpern}, these metrics have limitations, particularly in nonmonotone settings. 
It is important to note that standard gap functions are not applicable to our settings due to Assumption~\ref{as:A2}.
Regarding oracle complexity, several works, e.g., \cite{alacaoglu2021stochastic,beznosikov2023stochastic,gorbunov2022stochastic,loizou2021stochastic} claim an oracle complexity of $\mathcal{O}(n + \sqrt{n}\epsilon^{-1})$  to attain an $\epsilon$-solution, but this is measured using a restricted gap function. 
Again, as highlighted in \cite{cai2023variance,diakonikolas2020halpern}, this certification does not translate to the operator residual norm and is inapplicable to nonmonotone settings.
Therefore, a direct comparison between our results and these previous works is challenging due to these methodological discrepancies (see Table~\ref{tbl:existing_vs_ours}).

%%
\textit{$\mathrm{(iii)}$~New and simple algorithms.} 
Various existing stochastic methods for solving \eqref{eq:VIP} and \eqref{eq:NI} rely on established techniques. 
These include mirror-prox/averaging and extragradient-type schemes combined with the classic Robbin-Monro stochastic approximation \citep{RM1951} (e.g., \cite{cui2021analysis,iusem2017extragradient,juditsky2011solving,kannan2019optimal,kotsalis2022simple,yousefian2018stochastic}). 
Some approaches utilize increasing mini-batch sizes for variance reduction (e.g., \cite{iusem2017extragradient,pethick2023solving}).
Recent works have explored alternative variance-reduced methods for \eqref{eq:NI} and its special cases (e.g., \cite{alacaoglu2022beyond,alacaoglu2021stochastic,bot2019forward,cai2022stochastic,davis2022variance}). 
However, these methods primarily adapt existing optimization estimators to approximate the operator $G$ without significant differences.
Our approach departs from directly approximating $G$.  
Instead, we construct an intermediate quantity $S_{\gamma}^k := Gx^k - \gamma Gx^{k-1}$ as a linear combination of two consecutive evaluations of $G$ (i.e. $Gx^k$ and $Gx^{k-1}$). 
We then develop stochastic variance-reduced estimators specifically for $S_{\gamma}^k$.
Note that $S^k_{\gamma}$ alone is not new, but our idea of using it in stochastic methods is new.
This idea allows us to design new and simple algorithms with a single loop for both \eqref{eq:NE} and \eqref{eq:NI}  where the state-of-the-art oracle complexity is achieved (\textit{cf.} Sections~\ref{sec:VROG4NE} and \ref{sec:VROG4NI}).


%%%%%%%%%%%%%%%%%%%%%%%%%%%%%%%%%%%%
%%% Basic Assumptions.
%%%%%%%%%%%%%%%%%%%%%%%%%%%%%%%%%%%%
\beforesubsec
\subsection{Basic assumptions}\label{subsec:basic_assumptions}
\aftersubsec
We tackle both \eqref{eq:NE} and \eqref{eq:NI} covered by the following basic assumptions (see \cite{Bauschke2011} for terminologies and concepts used in these assumptions).

%%% Assumption A.0.
\begin{assumption}\label{as:A0}[\textit{Well-definedness}]
$\zer{\Psi}$ of \eqref{eq:NI} and $\zer{G}$ of \eqref{eq:NE} are nonempty.
\end{assumption}

%%% Assumption A.0.b
\begin{assumption}\label{as:A0b}[\textit{Maximal monotonicity of $T$}]
$T$ in \eqref{eq:NI} is maximally monotone on $\dom{T}$.
\end{assumption}

%%% Assumption A.1.
\begin{assumption}\label{as:A1}[\textit{Lipschitz continuity of $G$}]
$G$ in \eqref{eq:finite_sum_form} is $L$-averaged Lipschitz continuous, i.e. $\forall x, y \in \dom{G}$:
\vspace{-0.75ex}
\begin{equation}\label{eq:L_smooth}
\begin{array}{l}
\frac{1}{n} \sum_{i=1}^n \norms{G_ix - G_iy}^2 \leq L^2\norms{x - y}^2.
\end{array}
\vspace{-0.5ex}
\end{equation}
\end{assumption}

%%% Assumption A.2.
\begin{assumption}\label{as:A2}[\textit{Weak-Minty solution}]
There exist a solution $x^{\star} \in \zer{\Psi}$ and  $\kappa \geq 0$ such that $\iprods{Gx + v, x - x^{\star}} \geq - \kappa \norms{Gx + v}^2$ for all $x \in \dom{\Psi}$ and $v \in Tx$.
\end{assumption}

While Assumption~\ref{as:A0} is basic, Assumption~\ref{as:A0b} guarantees the single-valued and well-definiteness of the resolvent $J_{T}$ of $T$. 
In fact, this assumption can be relaxed to some classes of nonmonotone operators $T$, but we omit this extension.
The $L$-averaged Lipschitz continuity \eqref{eq:L_smooth} is standard and has been used in most deterministic, randomized, and stochastic methods.
It is slightly stronger that the $L$-Lipschitz continuity of the sum $G$.
The star-co-hypomonotonicity in Assumption~\ref{as:A2} is significantly different from  the star-strong monotonicity used in, e.g., \cite{kotsalis2022simple}.
It covers a class of nonmonotone operators $G$ (see Supp. Doc. \ref{apdix:subsec:nonmonotone_examples} for a concrete example).
It is also weaker than the co-hypomotonicity, used, e.g., in \cite{cai2023variance}.

%%%%%%%%%%%%%%%%%%%%%%%%%%%%%%%%%%%%
%%% 1.3. Contribution and related work.
%%%%%%%%%%%%%%%%%%%%%%%%%%%%%%%%%%%%
\beforesubsec
\subsection{Contribution and related work}\label{subsec:contribution_related_work}
\aftersubsec
Our primary goal is to develop a class of variance-reduction methods to solve both \eqref{eq:NE} and \eqref{eq:NI},  their special cases such as  \eqref{eq:VIP} and \eqref{eq:minimax}, and equivalent problems, like \eqref{eq:FP_prob}.

%%%%% 
\textbf{Our contribution.} Our main contribution consists of:  
\begin{compactitem}
\item[(a)] We exploit the variable $S_{\gamma}^k$ in \eqref{eq:Sk_op} and introduce a class of unbiased variance-reduced estimators $\widetilde{S}^k_{\gamma}$ for $S^k_{\gamma}$, not for $G$,  covered by Definition~\ref{de:ub_SG_estimator}.

\item[(b)] We construct two instances of $\widetilde{S}^k_{\gamma}$ by leveraging the SVRG \citep{SVRG} and SAGA \citep{Defazio2014} estimators, respectively that fulfill our Definition~\ref{de:ub_SG_estimator}.
These estimators are also of independent interest, and can be applied to develop other methods.

\item[(c)] We develop a variance-reduced forward-reflected-type method \eqref{eq:VROG4NE} to solve \eqref{eq:NE} required $\mathcal{O}( n \! + \! n^{2/3}\epsilon^{-2} )$ evaluations of $G_i$ to obtain an $\epsilon$-solution.  

\item[(d)] We design a novel stochastic variance-reduced forward-reflected-backward splitting method \eqref{eq:VROG4NI} to solve \eqref{eq:NI}, also required $\mathcal{O}( n \! + \! n^{2/3}\epsilon^{-2} \big)$ evaluations of $G_i$.
\end{compactitem}
%%%%%
Table \ref{tbl:existing_vs_ours} below compares our work and some existing single-loop variance-reduction methods, but in either the co-coercive or monotone settings.
Now, let us highlight the following points of our contribution. 
First, our intermediate quantity $S^k_{\gamma}$ can be viewed as a generalization of the forward-reflected-backward splitting (FRBS) operator \citep{malitsky2020forward} or an optimistic gradient operator \citep{daskalakis2018training} used in the literature. 
However, the chosen range $\gamma \in (1/2, 1)$ excludes these classical methods from recovering as special cases of $S_{\gamma}^k$.
Second, since our SVRG and SAGA estimators are designed specifically for $S_{\gamma}^k$, they differ from existing estimators in the literature, including recent works \citep{alacaoglu2022beyond,alacaoglu2021stochastic,bot2019forward}.
Third, both proposed algorithms are single-loop and straightforward to implement.
Fourth, our algorithm for nonlinear inclusions \eqref{eq:NI} significantly differs from existing methods, including deterministic ones, due to the additional term $\gamma^{-1}(2\gamma-1)(y^k-x^k)$.
For a comprehensive survey of deterministic methods, we refer to \cite{tran2023sublinear}.
Fifth, our oracle complexity estimates rely on the metric $\mathbb{E}[\norms{Gx^k}^2]$ or $\mathbb{E}[\norms{Gx^k+v^k}^2]$ for $v^k\in Tx^k$, commonly used in nonmonotone settings. 
Unlike the monotone case, this metric cannot be directly converted to a gap function, see, e.g., \cite{alacaoglu2022beyond,alacaoglu2021stochastic}. 
Our complexity bounds match the best known in stochastic nonconvex optimization using SAGA or SVRG without additional enhancements, e.g., utilizing a nested technique as in \cite{zhou2018stochastic}.

\begin{table*}[ht!]
\vspace{-0ex}
\newcommand{\cell}[1]{{\!\!}{#1}{\!\!}}
\setlength{\tabcolsep}{4pt}      % default is 6pt (horizontal padding)
\renewcommand{\arraystretch}{0.9}
\begin{small}
\begin{center}
\caption{Comparison of recent existing single-loop variance-reduction methods and our algorithms}\label{tbl:existing_vs_ours}
\vspace{-1ex}
\resizebox{0.9\textwidth}{!}{  
\hspace{-1ex}
\begin{tabular}{|c|c|c|c|c|c|c|} \hline
Papers & Problem & \cell{Assumptions}  & Estimators &  \cell{Residual Rates} & Oracle Complexity \\ \hline
\citep{davis2022variance}       &  \eqref{eq:NE} and \eqref{eq:NI}  & co-coercive/SQM  & {\!\!\!} SVRG \& SAGA {\!\!\!}   &  linear  & $\BigOs{(L/\mu)\log(\epsilon^{-1})}$  \\ \hline
\citep{tran2024accelerated} & \eqref{eq:NE} and \eqref{eq:NI} & co-coercive      & a class   & $\BigOs{1/k^2}$  &  $\mcal{O}( n + n^{2/3}\epsilon^{-1} )$  \\ \hline
\citep{cai2023variance} & \eqref{eq:NE} and \eqref{eq:NI} & co-coercive      & SARAH  &   $\BigOs{1/k^2}$ &  $\mcal{O}\big( n + n^{1/2}\log(n)\epsilon^{-1} \big)$ \\ \hline
\citep{alacaoglu2021stochastic} & \eqref{eq:VIP} & monotone          & SVRG & \ding{55}   & \mymgt{$\mcal{O}\big( n + n^{1/2}\epsilon^{-1} \big)$}  \\ \hline
\citep{alacaoglu2022beyond} & \eqref{eq:VIP} & monotone          & SVRG & \ding{55}   & \mymgt{$\mcal{O}\big( n\epsilon^{-1} \big)$} \\ \hline
\myblue{\textbf{Ours}} &   \eqref{eq:NE} and \eqref{eq:NI}  & {\!\!\!} \myblue{weak Minity} {\!\!\!} & \myblue{a class}   &{\!\!\!\!\!}  \makecell{\myblue{ $\BigOs{1/k}$}} {\!\!\!\!\!} &{\!\!\!}  \mblue{$\mcal{O}\big(n + n^{2/3}\epsilon^{-2}\big)$} {\!\!\!}  \\ \hline
\end{tabular}}
\end{center}
\end{small}
%\\\vspace{1ex}
{\footnotesize
\vspace{-0.5ex}
\textbf{Notes:} 
\textbf{SQM} means ``strong quasi-monotonicity";  
\textbf{Residual Rate} is the convergence rate on $\Expn{\norms{Gx^k}^2}$ or $\Expn{\norms{Gx^k + v^k}^2}$ for $v^k \in Tx^k$; 
and \textbf{\myblue{a class}} is a class of variance-reduced estimators.
The complexity of \cite{alacaoglu2021stochastic} and \citep{alacaoglu2022beyond} marked by \mymgt{magenta}   is on a gap function, a different metric than in other works in Table~\ref{tbl:existing_vs_ours}.
Thus, it is unclear how to compare them.
}{\!\!\!}
\vspace{-2ex}
\end{table*}


\textbf{Related work.}
Since both theory and solution methods for solving \eqref{eq:NE} and \eqref{eq:NI} are ubiquitous, see, e.g., \cite{Bauschke2011,reginaset2008,Facchinei2003,phelps2009convex,ryu2022large,ryu2016primer}, especially under the monotonicity, we only highlight the most recent related works (see more in Supp. Doc. \ref{apdix:sec:further_related_work}).

\vspace{-0.75ex}
\textit{$\mathrm{(i)}$~Weak-Minty solution.}
Assumption~\ref{as:A2} is known as a weak-Minty solution of \eqref{eq:NI} (in particular, of \eqref{eq:NE}), which has been widely used in recent works, e.g., \cite{bohm2022solving,diakonikolas2021efficient,lee2021fast,pethick2022escaping,tran2023extragradient} for deterministic methods  and, e.g., \cite{lee2021fast,pethick2023solving,tran2022accelerated} for stochastic methods.
This weak-Minty solution condition is weaker than the co-hypomonotonicity \citep{bauschke2020generalized}, which was used earlier in  proximal-point methods \citep{combettes2004proximal}.
Diakonikolas \textit{et al.} exploited this condition to develop an extragradient variant (called EG+) to solve \eqref{eq:NE}.
Following up works include \cite{bohm2022solving,cai2022baccelerated,luo2022last,pethick2022escaping,tran2023extragradient}. 
A recent survey in \cite{tran2023sublinear} provides several deterministic methods that rely on this condition.
This assumption covers a class of nonmonotone operators $G$ or $G+T$. 

\vspace{-0.75ex}
\textit{$\mathrm{(ii)}$~Stochastic approximation methods.}
Stochastic methods for both \eqref{eq:NE} and \eqref{eq:NI} and their special cases have been extensively developed, see, e.g., \cite{juditsky2011solving,kotsalis2022simple,pethick2023solving}.
Several methods exploited mirror-prox and averaging techniques such as \cite{juditsky2011solving,kotsalis2022simple}, while others relied on projection or extragradient schemes, e.g., \cite{cui2021analysis,iusem2017extragradient,kannan2019optimal,pethick2023solving,yousefian2018stochastic}. 
Many of these algorithms use standard Robbin-Monro stochastic approximation with fixed or increasing batch sizes.
Some other works generalized the analysis to a general class of algorithms such as \citep{beznosikov2023stochastic,gorbunov2022stochastic,loizou2021stochastic} covering both standard stochastic approximation and variance reduction algorithms.

%%% Variance reductions.
\vspace{-0.75ex}
\textit{$\mathrm{(iii)}$~Variance-reduction methods.}	
Variance-reduction techniques have been broadly explored in optimization, where many estimators were proposed, including SAGA \citep{Defazio2014}, SVRG  \citep{SVRG}, SARAH \citep{nguyen2017sarah}, and Hybrid-SGD \citep{Tran-Dinh2019,Tran-Dinh2019a}, and STORM \citep{Cutkosky2019}.
Researchers have adopted these estimators to develop methods for \eqref{eq:NE} and \eqref{eq:NI}. 
For example, \cite{davis2022variance} proposed a SAGA-type methods for \eqref{eq:NE} under a [quasi]-strong monotonicity.
The authors in \cite{alacaoglu2022beyond,alacaoglu2021stochastic} employed SVRG estimators and developed methods for \eqref{eq:VIP}.
Other works can be found in \cite{bot2019forward,carmon2019variance,chavdarova2019reducing,huang2022accelerated,palaniappan2016stochastic,yu2022fast}.
All of these results are different from ours.
Some recent works exploited Halpern's fixed-point iterations and develop corresponding variance-reduced methods, see, e.g., \cite{cai2023variance,cai2022stochastic}.
However, varying parameters or incorporating double-loop/inexact methods must be used to achieve improved theoretical oracle complexity.
We believe that such approaches may be challenging to select parameters and to implement in practice.
Finally, unlike optimization, it has been realized that using biased estimators such as SARAH or Hybrid-SGD/STORM for \eqref{eq:NI} (including \eqref{eq:minimax} and \eqref{eq:VIP}) is challenging due to the lack of an objective function, a key metric to prove convergence, and product terms like $\iprods{e^k, x^{k+1} - x^{\star}}$  in convergence analyses, where $e^k$ is a bias rendered from $\widetilde{S}^k_{\gamma}$ (see Supp. Doc. \ref{apdix:sec:further_related_work}).

\textbf{Notation.}
We use $\Fc_k := \sigma(x^0, x^1, \cdots, x^k)$ to denote the $\sigma$-algebra generated by $x^0,\cdots, x^k$ up to the iteration $k$.
$\Expsk{k}{\cdot} = \Expk{\cdot \mid \Fc_k}$ denotes the conditional expectation w.r.t. $\Fc_k$, and $\Expk{\cdot}$ is the total expectation. 
We also use $\BigO{\cdot}$ to characterize convergence rates and oracle complexity.
For an operator $G$, $\dom{G} := \sets{x : Gx\neq\emptyset}$ denotes its domain, and $J_G$ denotes its resolvent. 

\textbf{Paper organization.}
Section~\ref{sec:VRS_Estimator} introduces  $S^k_{\gamma}$ and defines a class of stochastic estimators for it.
It also constructs two instances: SVRG and SAGA, and proves their key properties.
Section~\ref{sec:VROG4NE} develops an algorithm for solving \eqref{eq:NE} and establishes its oracle complexity.
Section~\ref{sec:VROG4NI} designs a new algorithm for solving \eqref{eq:NI} and proves its oracle complexity.
Section~\ref{sec:NumExps} presents two concrete numerical examples.
Proofs and additional results are deferred to Sup. Docs. \ref{apdix:sec:further_related_work} to \ref{apdx:sec:num_experiments}.

%%%%%%%%%%%%%%%%%%%%%%%%%%%%%%%%%%%%%%%%%%%%%
%%% 2. Variance-Reduced Stochastic Estimators for Operators.
%%%%%%%%%%%%%%%%%%%%%%%%%%%%%%%%%%%%%%%%%%%%%
\beforesec
\section{Forward-Reflected Quantity and Its Stochastic Variance-Reduced Estimators}\label{sec:VRS_Estimator}
\aftersec
We first define our forward-reflected quantity (FRQ) for $G$ in \eqref{eq:NE} and \eqref{eq:NI} using here. 
Next, we propose a class of unbiased variance-reduced estimators for FRQ.
Finally, we construct two instances relying on the two well-known estimators: SVRG from \cite{SVRG} and SAGA from \cite{Defazio2014}.
%However, any other estimator could be used if it satisfies our Definition~\ref{de:ub_SG_estimator} below.

%%%%%%%%%%%%%%%%%%%%%%%%%%%%%%%%%%%%%%%%%%%%%%%%%%%
%%% 2.1. Optimistic gradient operators
%%%%%%%%%%%%%%%%%%%%%%%%%%%%%%%%%%%%%%%%%%%%%%%%%%%
\beforesubsec
\subsection{The forward-reflected quantity}\label{subsec:OG_operator}
\aftersubsec
Our methods for solving \eqref{eq:NE} and \eqref{eq:NI} rely on the following intermediate quantity constructed from $G$ via two consecutive iterates $x^{k-1}$ and $x^k$ controlled by $\gamma \in [0, 1]$:
%\vspace{-0.5ex}
\begin{equation}\label{eq:Sk_op}
S_{\gamma}^k := Gx^k  - \gamma Gx^{k-1}.
\tag{FRQ}
%\vspace{-0.75ex}
\end{equation}
Here, $\gamma \in \big(\frac{1}{2}, 1\big)$ plays a crucial role in our methods in the sequel.
Clearly, if $\gamma = \frac{1}{2}$, then we can write $S_{1/2}^k = \frac{1}{2}Gx^k + \frac{1}{2}(Gx^k - Gx^{k-1}) = \frac{1}{2}[ 2Gx^k - Gx^{k-1}]$ used in both the forward-reflected-backward splitting (FRBS) method \citep{malitsky2020forward} and the optimistic gradient method  \citep{daskalakis2018training}.
In deterministic unconstrained settings (i.e. solving \eqref{eq:NE}), see \citep{tran2023sublinear}, FRBS  is also equivalent to Popov's past-extragradient method \citep{popov1980modification}, reflected-forward-backward splitting algorithm  \citep{cevher2021reflected,malitsky2015projected}, and optimistic gradient scheme \citep{daskalakis2018training}.
In the deterministic constrained case, i.e. solving \eqref{eq:NI}, these methods are different.
Since $\gamma \in \big(\frac{1}{2}, 1\big)$, our methods below exclude these classical schemes.
However, due to a similarity pattern of \eqref{eq:Sk_op} and  FRBS, we still term our quantity  $S_{\gamma}^k$ by the ``\textbf{forward-reflected quantity}'', abbreviated by FRQ.

%%%%%%%%%%%%%%%%%%%%%%%%%%%%%%%%%%%%%%%%%%%%%%%%%%%
%%% 2.2. Stochastic unbiased variance-reduced estimators for the OG operator
%%%%%%%%%%%%%%%%%%%%%%%%%%%%%%%%%%%%%%%%%%%%%%%%%%%
\beforesubsec
\subsection{Unbiased variance-reduced estimators for FRQ}\label{subsec:SG_estimators}
\aftersubsec
Now, let us propose  the following class of stochastic variance-reduced estimators $\widetilde{S}^k_{\gamma}$ of $S_{\gamma}^k$.

%%% Definition 3.1.
\begin{definition}\label{de:ub_SG_estimator}
A stochastic estimator $\widetilde{S}_{\gamma}^k$ is said to be a \textit{stochastic unbiased variance-reduced estimator} of $S_{\gamma}^k$ in \eqref{eq:Sk_op} if there exist constants $\rho \in (0, 1]$, $C \geq 0$ and $\hat{C} \geq 0$, and a nonnegative sequence $\sets{\Delta_k}$  such that:
% the following three conditions hold for $\widetilde{e}_{\gamma}^k := \widetilde{S}_{\gamma}^k - S^k_{\gamma}$:
\vspace{-1.25ex}
\begin{equation}\label{eq:ub_SG_estimator}
\hspace{-0.1ex}
\arraycolsep=0.1em
\left\{ \begin{array}{lcl}
\Expsk{k}{ \widetilde{S}_{\gamma}^k - S^k_{\gamma} }   =   0, \\ %\quad\qquad\qquad \textrm{(unbiasedness)} \\
\Expk{ \norms{ \widetilde{S}_{\gamma}^k - S^k_{\gamma} }^2 }   \leq    \Delta_k, \\ %\qquad\qquad \textrm{(bounded variance)}\\
\Delta_k  \leq   (1 - \rho )\Delta_{k-1} + C \cdot U_k + \hat{C} \cdot U_{k-1},  %\quad \textrm{(variance-reduction)}, %\\
%&& + {~} C U_k + \hat{C} U_{k-1}, %\\
%&&  + {~} \frac{\hat{C}}{n} \cdot \sum_{i=1}^n \Expk{ \norms{G_ix^{k-1} - G_ix^{k-2} }^2 }.
\end{array}\right.
\hspace{-4ex}
\vspace{-1.5ex}
\end{equation}
where $U_k := \frac{1}{n} \sum_{i=1}^n \Expk{ \norms{G_ix^k - G_ix^{k-1} }^2 }$.
\end{definition}
Here, $\Delta_{-1} = 0$, $x^{-2} = x^{-1} = x^0$, and  $\Expsk{k}{\cdot}$ and $\Expk{\cdot}$ are the conditional and total expectations defined earlier, respectively.  
The condition $\rho > 0$ is important to achieve a variance reduction as long as $x^k$ is close to $x^{k-1}$ and $x^{k-1}$ is close to $x^{k-2}$.
Otherwise, $ \widetilde{S}_{\gamma}^k$ may not be a variance-reduced estimator of $S_{\gamma}^k$.
Since $S_{\gamma}^k$ is evaluated at both $x^{k-1}$ and $x^k$, our bound for the estimator $\widetilde{S}_{\gamma}^k$ depends on three consecutive points $x^{k-2}$, $x^{k-1}$, and $x^k$, which is different from previous works, including \cite{alacaoglu2021forward,beznosikov2023stochastic,davis2022variance,driggs2019accelerating}.

Now, we will construct two variance-reduced estimators satisfying Definition~\ref{de:ub_SG_estimator} by exploiting SVRG \citep{SVRG} and SAGA \citep{Defazio2014}. 

%%% 3.1. SVRG Estimator for Root-Finding Probems.
\beforepara
\paragraph{$\mathrm{(a)}$~Loopless-SVRG estimator for $S^k_{\gamma}$.}\label{subsec:SVRG_estimator}
Consider a mini-batch $\Bc_k \subseteq [n] := \sets{1,2,\cdots, n}$ with a fixed batch size $b := \vert \Bc_k\vert$.
Denote $G_{\Bc_k}z := \frac{1}{b}\sum_{i \in \Bc_k}G_iz$ for $z\in\dom{G}$.
We define the following estimator for $S_{\gamma}^k$:
\vspace{-1ex}
\begin{equation}\label{eq:SVRG_estimator}
\begin{array}{lcl}
\widetilde{S}_{\gamma}^k & := &   (1-\gamma )( Gw^k  - G_{\Bc_k}w^k ) \\
&& + {~} G_{\Bc_k}x^k - \gamma G_{\Bc_k}x^{k-1}, 
\end{array}
\tag{L-SVRG}
\vspace{-1ex}
\end{equation}
where the snapshot point $w^k$ is selected randomly as follows:
\vspace{-0.75ex}
\begin{equation}\label{eq:SG_estimator_w}
w^{k+1} := \left\{\begin{array}{ll}
x^{k} &\text{with probability $\mbf{p}$}\\
w^k & \text{with probability $1 - \mbf{p}$}.
\end{array}\right.
\vspace{-0.75ex}
\end{equation}
The probability $\mbf{p} \in (0, 1)$ will appropriately be chosen later by nonuniformly flipping a coin.
This estimator is known as a loopless variant  \citep{kovalev2019don} of the SVRG estimator \citep{SVRG}.
However, it is different from existing ones used in root-finding algorithms, including \cite{davis2022variance} because we define it for $S_{\gamma}^k$, not for $Gx^k$.
In addition, the first term is also damped by a factor $1-\gamma$ to guarantee the unbiasedness of $\widetilde{S}_{\gamma}^k$ to $S_{\gamma}^k$.

The following lemma shows that $\widetilde{S}_{\gamma}^k$ satisfies Definition~\ref{de:ub_SG_estimator}.

%%% Lemma 1.
\begin{lemma}\label{le:SVGR_estimator}
Let $S_{\gamma}^k$ be given by \eqref{eq:Sk_op} and $\widetilde{S}_{\gamma}^k$ be generated by the SVRG estimator \eqref{eq:SVRG_estimator} and
\vspace{-0.75ex}
\begin{equation*}
\begin{array}{lcl}
\Delta_k &:= &  \frac{1}{nb} \sum_{i=1}^n \Expk{ \norms{ G_ix^k - \gamma G_ix^{k-1} - (1 - \gamma) G_iw^k }^2 }.
\end{array}
\vspace{-0.75ex}
\end{equation*}
Then, $\widetilde{S}_{\gamma}^k$ satisfies Definition~\ref{de:ub_SG_estimator} with this $\sets{ \Delta_k }$, $\rho := \frac{\mbf{p} }{2}$, $C := \frac{4 - 6\mbf{p} + 3\mbf{p}^2}{b\mbf{p}}$, and $\hat{C} := \frac{2\gamma^2(2 - 3\mbf{p} + \mbf{p}^2 )}{b\mbf{p}}$.
\end{lemma}

%%%%% SAGA Estimator.
\beforepara
\paragraph{$\mathrm{(b)}$~SAGA estimator for $S^k_{\gamma}$.}
Given $S^k_{\gamma}$ as in \eqref{eq:Sk_op} and a mini-batch estimator $G_{\Bc_k}$ as in \eqref{eq:SVRG_estimator}, we construct the following SAGA estimator for $S_{\gamma}^k$:
\vspace{-0.75ex}
\begin{equation}\label{eq:SAGA_estimator}
\begin{array}{lcl}
\widetilde{S}^k_{\gamma} & := & \big[ G_{\Bc_k}x^k - \gamma G_{\Bc_k}x^{k-1} - (1 - \gamma)\hat{G}_{\Bc_k}^k \big] \\
&& + {~} \frac{1-\gamma}{n}\sum_{i=1}^n\hat{G}_i^k, 
\end{array}
\tag{SAGA}
\vspace{-0.75ex}
\end{equation}
where $\Bc_k$ is a mini-batch of size $b$, and $\hat{G}^k_i$ is updated as
\vspace{-0.75ex}
\begin{equation}\label{eq:SAGA_ref_points}
\hat{G}_i^{k+1} := \left\{\begin{array}{lcl}
G_ix^{k} &\text{if}~ i \in \Bc_k, \\
\hat{G}_i^{k} & \text{otherwise}. %~i \notin \Bc_k.
\end{array}\right.
\vspace{-0.75ex}
\end{equation}
To form $\widetilde{S}^k_{\gamma}$, we need to store $n$ components $\hat{G}^k_i$ computed so far for $i \in [n]$ in a table $\Tc_k := [\hat{G}_1^k, \hat{G}_2^k, \cdots, \hat{G}_n^k]$ initialized at $\hat{G}_i^0 := G_ix^0$ for all $i\in [n]$.
Clearly, the SAGA estimator requires significant memory to store $\Tc_k$ if $n$ and $p$ are both large.
We have the following result.

%%% Lemma 2.2.
\begin{lemma}\label{le:SAGA_estimator}
Let $S^k_{\gamma}$ be defined by \eqref{eq:Sk_op} and $\widetilde{S}_{\gamma}^k$ be generated by the SAGA estimator \eqref{eq:SAGA_estimator}, and  
\vspace{-0.75ex}
\begin{equation*} 
\begin{array}{lcl}
 \Delta_k &:= &  \frac{1}{nb} \sum_{i=1}^n \Expk{ \norms{ G_ix^k - \gamma G_ix^{k-1} - (1-\gamma) \hat{G}_{i}^k }^2 }.
\end{array}
\vspace{-0.75ex}
\end{equation*}
Then,  $\widetilde{S}_{\gamma}^k$  satisfies Definition~\ref{de:ub_SG_estimator} with this $\sets{\Delta_k}$, $\rho := \frac{b}{2n} \in (0, 1]$, $C :=  \frac{[ 2(n - b)(2n+b) + b^2] }{ n b}$, and $\hat{C} := \frac{2(n - b)(2n+b) \gamma^2}{ nb}$.
\end{lemma}

We only provide two instances: \eqref{eq:SVRG_estimator} and \eqref{eq:SAGA_estimator} covered by Definition~\ref{de:ub_SG_estimator}.
However, we believe that similar estimators for $S_{\gamma}^k$ relied on, e.g., JacSketch \citep{gower2021stochastic} or SEGA \citep{hanzely2018sega}, among others can fulfill our Definition~\ref{de:ub_SG_estimator}.

%%%%%%%%%%%%%%%%%%%%%%%%%%%%%%%%%%%%%%%%%%%%%%%%%%%%%%%
%%% 3. Stochastic Fixed-Point Methods.
%%%%%%%%%%%%%%%%%%%%%%%%%%%%%%%%%%%%%%%%%%%%%%%%%%%%%%%
\beforesec
\section{A Variance-Reduced Forward-Reflected Method for [Non]linear Equations}\label{sec:VROG4NE}
\aftersec
\vspace{0.5ex}
We first utilize the class of stochastic estimators in Definition~\ref{de:ub_SG_estimator} to develop a variance-reduced forward-reflected (VFR) method for solving \eqref{eq:NE} under Ass.~\ref{as:A1} and \ref{as:A2}.

%%%%%%%%%%%%%%%%%%%%%%%%%%%%%%%%%%%%%%%%%%%%%%%%%%%
%%% 3.1. The VROG method and its convergence guarantee}
%%%%%%%%%%%%%%%%%%%%%%%%%%%%%%%%%%%%%%%%%%%%%%%%%%%
\beforesubsec
\subsection{The VFR method and its convergence guarantee}\label{subsec:OG_method}
\aftersubsec
\vspace{0.5ex}
\textbf{$\mathrm{(a)}$~VFR Method.}
Our method is described as follows.
\textit{
Starting from $x^0 \in \dom{G}$,  at each iteration $k \geq 0$, we construct an estimator $\widetilde{S}_{\gamma}^k$ satisfying Definition~\ref{de:ub_SG_estimator} with parameters $\rho \in (0, 1]$,  $ C \geq 0$, and $\hat{C} \geq 0$,  and then update
\vspace{-1ex}
\begin{equation}\label{eq:VROG4NE}
\begin{array}{lcl}
x^{k+1} &:= & x^k - \eta \widetilde{S}_{\gamma}^k,
\end{array}
\tag{VFR}
\vspace{-1ex}
\end{equation}
where $\eta > 0$ and $\gamma > 0$  are  determined below, $x^{-1} = x^{-2} := x^0$, and $\widetilde{S}_{\gamma}^0 := (1-\gamma)Gx^0$.}

At least two estimators  $\widetilde{S}_{\gamma}^k$: the \textit{Loopless-SVRG estimator} in \eqref{eq:SVRG_estimator} and the \textit{SAGA estimator} in \eqref{eq:SAGA_estimator}, can be used in our method \eqref{eq:VROG4NE}.
In terms of \textit{per-iteration complexity}, each iteration $k$ of \ref{eq:VROG4NE}, the loopless SVRG variant requires three mini-batch evaluations $G_{\Bc_k}w^k$, $G_{\Bc_k}x^k$, and $G_{\Bc_k}x^{k-1}$ of $G$, and occasionally computes one full evaluation $Gw^k$ of $G$ with the probability $\mbf{p}$.
It needs one more mini-batch evaluation $G_{\Bc_k}x^{k-1}$ compared to SVRG-type methods in optimization.
Similarly, the SAGA estimator also requires two mini-batch evaluations $G_{\Bc_k}x^k$ and $G_{\Bc_k}x^{k-1}$, which is one more mini-batch $G_{\Bc_k}x^{k-1}$ compared to SAGA-type methods in optimization, see, e.g., \cite{reddi2016fast}.
The \ref{eq:SAGA_estimator} estimator can avoid the occasional full-batch evaluation $Gw^k$ from \ref{eq:SVRG_estimator}, but as a compensation, we need to store a table  $\Tc_k := [\hat{G}_1^k, \hat{G}_2^k, \cdots, \hat{G}_n^k]$, which requires significant memory in the large-scale regime.

%%%%%%%%%%%%%%%%%%%%%%%%%%%%
\textbf{$\mathrm{(b)}$~Convergence guarantee.} 
Fixed $\gamma \in \left(\frac{1}{2}, 1\right)$, with $\rho$, $C$, and $\hat{C}$ as in Definition~\ref{de:ub_SG_estimator} we define
\begin{equation}\label{eq:VROG4NE_constants}
\hspace{-0.5ex}
\begin{array}{ll}
M := \frac{\gamma(1+5\gamma)}{3(2\gamma-1)} + \frac{1 + 6\gamma}{3(2\gamma -1) } \cdot \frac{C + \hat{C}}{\rho} \quad \text{and} \quad \delta := \frac{2\gamma-1}{8\sqrt{M}}.
\end{array}
\hspace{-2ex}
\end{equation}
%%%
Then, the following theorem states the convergence of \eqref{eq:VROG4NE}, whose proof is given in Supp. Doc. \ref{apdx:sec:VROG4NE}.
%
%%% Theorem 3.1.
\begin{theorem}\label{th:VROG4NE_convergence}
Let us fix $\gamma \in \left(\frac{1}{2}, 1\right)$, and define $M$ and $\delta$ as in \eqref{eq:VROG4NE_constants}.
Suppose that Assumptions~\ref{as:A0}, ~\ref{as:A1}, and ~\ref{as:A2} hold for \eqref{eq:NE} with some $\kappa \geq 0$ such that $L\kappa \leq \delta$.
Let $\sets{x^k}$ be generated by \eqref{eq:VROG4NE} using a learning rate $\eta > 0$ such that $\frac{8\kappa}{2\gamma-1} \leq \eta \leq \frac{1}{L\sqrt{M}}$.
Then, the following bounds hold:
\vspace{-1ex}
\begin{equation}\label{eq:SOG_convergence_bound1}
\begin{array}{lcl}
\dfrac{1}{K+1}\displaystyle\sum_{k=0}^{K} \Expk{ \norms{Gx^k}^2 }  & \leq &    \dfrac{ \Theta_1 \norms{x^0 - x^{\star}}^2 }{K + 1}, \vspace{0.25ex}\\
\dfrac{1}{K+1}\displaystyle\sum_{k=1}^{K} \Expk{ \norms{x^k - x^{k-1}}^2 } & \leq &  \dfrac{\Theta_2 \norms{x^0 - x^{\star}}^2 }{K + 1}  ,
\end{array}
\vspace{-1ex}
\end{equation}
where $\Theta_1 := \frac{2(1+L^2\eta^2)}{\gamma(1-\gamma)\eta^2}$ and $\Theta_2 := \frac{8(1 + L^2\eta^2)}{3(2\gamma-1) (1 - M L^2\eta^2)}$
\end{theorem}


Theorem~\ref{th:VROG4NE_convergence} only proves a $\BigO{1/K}$ convergence rate of both $\frac{1}{K+1}\sum_{k=0}^K \Expk{ \norms{Gx^{k}}^2}$ and $\frac{1}{K+1}\sum_{k=1}^K \Expk{ \norms{x^k - x^{k-1}}^2}$, but does not characterize the oracle complexity of \eqref{eq:VROG4NE}.
If we choose $\gamma := \frac{3}{4}$, then from \eqref{eq:VROG4NE_constants}, we have $M = \frac{57}{24} + \frac{11(C+\hat{C})}{3\rho}$ and $\delta = \frac{1}{16\sqrt{M}}$, which can simplify the bounds in Theorem~\ref{th:VROG4NE_convergence}.
In addition, it allows $\kappa > 0$ such that $L\kappa \leq \delta = \BigOs{\sqrt{\rho}}$, which means that $\kappa$ can be positive, but depends on $\sqrt{\rho}$.
This condition allows us to cover a class of nonmonotone operators $G$, where a weak-Minty solution exists as stated in  Assumption~\ref{as:A2}.

%%%%%%%%%%%%%%%%%%%%%%%%%%%%%%%%%%%%%%%%%%%
%%%% 4.2. Oracle Complexity Bounds for Specific Estimators.
%%%%%%%%%%%%%%%%%%%%%%%%%%%%%%%%%%%%%%%%%%%
\vspace{-0.25ex}
\beforesubsec
\subsection{Complexity Bounds of \ref{eq:VROG4NE} with SVRG and SAGA}\label{subsec:Oracle_complexity}
\aftersubsec
\vspace{0.3ex}
Let us first apply Theorem~\ref{th:VROG4NE_convergence} to the mini-batch SVRG estimator \eqref{eq:SVRG_estimator} in Section~\ref{sec:VRS_Estimator}.
For simplicity, we choose $\gamma := \frac{3}{4}$ and $\eta := \frac{1}{L\sqrt{M}}$, but any $\gamma \in \left(\frac{1}{2}, 1\right)$ still works.

%%% Corollary 3.2.
\begin{corollary}\label{co:SVRG_convergence1_v2}
Suppose that Assumptions~\ref{as:A0}, ~\ref{as:A1}, and ~\ref{as:A2} hold for \eqref{eq:NE} with $\kappa \geq 0$ as in Theorem \ref{th:VROG4NE_convergence}.
Let $\sets{x^k}$ be generated by \eqref{eq:VROG4NE} using \eqref{eq:SVRG_estimator}, $\gamma := \frac{3}{4}$, and $\eta := \frac{1}{L\sqrt{M}} \geq \frac{0.1440 \sqrt{b}\mbf{p}}{L}$, provided that $b\mbf{p}^2 \leq 1$.  
Then 
%, the following bound holds:
\vspace{-0.75ex}
\begin{equation}\label{eq:SFP_convergence_bound_v2}
%\begin{array}{ll}
\frac{1}{K+1}\sum_{k=0}^{K} \Expk{ \norms{Gx^k}^2 }   \leq     \frac{526 L^2  \norms{x^0 - x^{\star}}^2 }{ b\mbf{p}^2 (K + 1)}.
%\end{array}
\vspace{-0.5ex}
\end{equation}
For $\epsilon > 0$, if we choose $\mbf{p} :=  n^{-1/3}$ and $b := \lfloor n^{2/3} \rfloor$, then \eqref{eq:VROG4NE} requires $\Tc_{G_i} := n + \big\lfloor \frac{4\Gamma  L^2R_0^2 n^{2/3} }{\epsilon^2} \big\rfloor$ evaluations of $G_i$ to attain $\frac{1}{K+1}\sum_{k=0}^{K} \Expk{ \norms{Gx^k}^2 } \leq \epsilon^2$, where $\Gamma := 731$.
\end{corollary}

Corollary~\ref{co:SVRG_convergence1_v2} states that the oracle complexity of \eqref{eq:VROG4NE} is $\BigO{n + n^{2/3}\epsilon^{-2}}$, matching  (up to a constant) the one of SVRG in nonconvex optimization, see, e.g., \cite{allen2016variance,Reddi2016a}.
It improves by a factor $\BigO{n^{1/3}}$ compared to deterministic counterparts.
This complexity is known to be the best for SVRG so far without any additional enhancement (e.g., nested techniques \citep{zhou2018stochastic}) even for a special case of \eqref{eq:NE}: $Gx = \nabla{f}(x)$ in nonconvex optimization.

Note that  $\eta$ can be computed explicitly when $b$ and $\mbf{p}$ are given.
For example, if $n = 10000$ and we choose $\mbf{p} = n^{-1/3} = 0.0464$ and $b = \lfloor n^{2/3}\rfloor = 464$, then $\eta = \frac{0.1456}{L}$.
If $\mbf{p} = 0.1$, then $\eta = \frac{0.3038}{L}$.
Note that, in general,  we can choose appropriate $p := \BigOs{n^{-1/3}}$ and $b := \BigOs{n^{2/3}}$.

Alternatively, we can apply Theorem~\ref{th:VROG4NE_convergence} to \eqref{eq:SAGA_estimator}. 
%the mini-batch SAGA estimator \eqref{eq:SAGA_estimator}.

%%% Corollary 3.3.
\begin{corollary}\label{co:SVRG_convergence1_v3}
Suppose that Assumptions~\ref{as:A0}, ~\ref{as:A1}, and ~\ref{as:A2} hold for \eqref{eq:NE} with $\kappa \geq 0$ as in Theorem~\ref{th:VROG4NE_convergence}.
Let $\sets{x^k}$ be generated by \eqref{eq:VROG4NE} using  \eqref{eq:SAGA_estimator}, $\gamma := \frac{3}{4}$, and $\eta := \frac{1}{L\sqrt{M}} \geq \frac{0.1494 b^{3/2}}{nL}$, provided that $1\leq b \leq n^{2/3}$.
Then %, the following bound holds:
\vspace{-0.75ex}
\begin{equation}\label{eq:SFP_convergence_bound_v3}
%\begin{array}{ll}
\frac{1}{K+1}\sum_{k=0}^{K} \Expk{ \norms{Gx^k}^2 }   \leq     \frac{489 L^2 \norms{x^0 - x^{\star}}^2 }{ b\mbf{p}^2 (K + 1)}.
%\end{array}
\vspace{-0.5ex}
\end{equation}
Moreover, for a given $\epsilon > 0$, if we choose $b := \lfloor n^{2/3} \rfloor$, then  \eqref{eq:VROG4NE} requires $\Tc_{G_i} := n + \big\lfloor \frac{ 3\Gamma L^2R_0^2 n^{2/3}}{ \varepsilon^2} \big\rfloor$ evaluations of $G_i$ to achieve $\frac{1}{K+1}\sum_{k=0}^{K} \Expk{ \norms{Gx^k }^2 } \leq \epsilon^2$, where $\Gamma := 2816$.
\end{corollary}

Similar to Corollary \ref{co:SVRG_convergence1_v2}, the learning rate $\eta$ in Corollary~\ref{co:SVRG_convergence1_v3} can explicitly be computed if we know $n$ and $b$.
For instance, if $n = 10000$, and we choose $b = \lfloor n^{2/3}\rfloor$, then $\eta = \frac{0.1603}{L}$.

\vspace{-0.5ex}
If $\kappa = 0$, i.e. $G$ reduces to a star-monotone operator, then we can choose $\gamma \in \left(\frac{1}{2}, 1\right)$ and $\eta$ as:
\vspace{-0.25ex}
\begin{compactitem}
\item For SVRG, we have  $\eta \in \big(0, \frac{1}{L\sqrt{M}}\big]$. 
If we choose $\mbf{p} = \BigOs{n^{-1/3}}$ and $b = \BigOs{n^{2/3}}$, then $\eta = \BigOs{\frac{1}{L}}$;
\item For SAGA, we have $\eta \in \big(0, \frac{1}{L\sqrt{M}}\big]$.
If we choose $b = \BigOs{n^{2/3}}$, then $\eta = \BigOs{\frac{1}{L}}$.
\end{compactitem}
\vspace{-0.25ex}
%Clearly, in the star-monotone case ($\kappa=0$), our choice of $\mathbf{p}$ and $b$ in SVRG is different from  \cite{alacaoglu2021stochastic,alacaoglu2021forward}.
Hitherto, the constant factor $\Gamma$ in both corollaries is still relatively large, but it can be further improved  by refining our technical proofs (e.g., carefully using Young's inequality).

%%%%%%%%%%%%%%%%%%%%%%%%%%%%%%%%%%%%%%%
%% 4. Variance-Reduced Optimistic Gradient Methods
%%%%%%%%%%%%%%%%%%%%%%%%%%%%%%%%%%%%%%%
\beforesec
\section{A New Variance-Reduced FRBS Method for Nonmonotone Generalized Equations}\label{sec:VROG4NI}
\aftersec
\vspace{0.25ex}
In this section, we develop a new stochastic variance-reduced forward-reflected-backward splitting (FRBS) method to solve \eqref{eq:NI} under Assumptions~\ref{as:A0b}, \ref{as:A1}, and \ref{as:A2}.

%%%%%%%%%%%%%%%%%%%%%%%%%%%%%%%%%%%%%%%%%%%
%%%% 4.1. The algorithm and its convergence.
%%%%%%%%%%%%%%%%%%%%%%%%%%%%%%%%%%%%%%%%%%%
\beforesubsec
\subsection{The algorithm and its convergence}\label{subsec:VROG4NI_alg_convergence}
\aftersubsec
\vspace{0.5ex}
\textbf{$\mathrm{(a)}$~The variance-reduced FRBS method (VFRBS).}
Our scheme for solving \eqref{eq:NI} is as follows.
\textit{
Starting from $x^0 \in \dom{\Psi}$,  at each iteration $k \geq 0$, we generate an estimator $\widetilde{S}_{\gamma}^k$ satisfying Definition~\ref{de:ub_SG_estimator} and update
\vspace{-0.5ex}
\begin{equation}\label{eq:VROG4NI}
%\hspace{-2ex}
%\begin{array}{lcl}
x^{k+1} :=  x^k - \eta  \widetilde{S}_{\gamma}^k - \eta \big( \gamma v^{k+1} - (2\gamma - 1)v^k \big),
%\end{array}
\tag{VFRBS}
%\hspace{-1ex}
\vspace{-0.5ex}
\end{equation}
where $\eta > 0$ and $\gamma > 0$  are determined later, $v^k \in Tx^k$,  $x^{-1} = x^{-2} := x^0$, and $\widetilde{S}_{\gamma}^0 := (1-\gamma)Gx^0$.
}

%%% (b) Implementable version,
\textbf{$\mathrm{(b)}$~Implementable version.}
Since $v^{k+1} \in Tx^{k+1}$ appears on the RHS of \eqref{eq:VROG4NI}, using the resolvent  $J_{\gamma\eta T} := (\Id + \gamma\eta T)^{-1}$ of $T$, we can rewrite \eqref{eq:VROG4NI} equivalently to  
\vspace{-1ex}
\begin{equation}\label{eq:VROG4NI_impl}
\left\{\begin{array}{lcl}
y^{k+1} &:= & x^k - \eta  \widetilde{S}_{\gamma}^k + \frac{(2\gamma - 1)}{\gamma}(y^k - x^k), \\
x^{k+1} &:= &  J_{\gamma\eta T}\big( y^{k+1}\big). 
\end{array}\right.
\vspace{-1ex}
\end{equation}
Here, $y^0 \in \dom{\Psi}$ is given, and $x^0 = x^{-1} := J_{\gamma\eta T}(y^0)$.
This is an implementable variant of \eqref{eq:VROG4NI} using the resolvent $J_{\gamma\eta T}$.
Clearly, if $\gamma = \frac{1}{2}$, then \eqref{eq:VROG4NI_impl} reduces to 
\vspace{-0.5ex}
\begin{equation*} 
\begin{array}{lcl}
x^{k+1} :=   J_{(\eta/2) T}\big( x^k - \eta  \widetilde{S}_{1/2}^k \big),
\end{array}
\vspace{-0.5ex}
\end{equation*}
which can be viewed as a stochastic forward-reflected-backward splitting scheme.
However, our $\gamma \in \big(\frac{1}{2}, 1)$, making \eqref{eq:VROG4NI_impl} different from existing methods, even in the deterministic case.

Compared to \cite{alacaoglu2021stochastic}, \eqref{eq:VROG4NI_impl} requires only one $J_{\gamma\eta T}$ as in \cite{alacaoglu2022beyond}, while \cite{alacaoglu2021stochastic} needs more than ones.
Moreover, our estimator $\widetilde{S}^k_{\gamma}$ is also different from  \cite{alacaoglu2021stochastic}.  
Compared to \cite{beznosikov2023stochastic} and  \cite{alacaoglu2022beyond}, the term $\gamma^{-1}(2\gamma - 1)(y^k - x^k)$ makes it different from SGDA in \cite{beznosikov2023stochastic} and the golden-ratio method in \cite{alacaoglu2022beyond}, and also other existing deterministic methods.

%%%% (c) Approximate solution certification
\textbf{$\mathrm{(c)}$~Approximate solution certification.}
To certify an approximate solution of \eqref{eq:NI}, we note that its exact solution $x^{\star}\in\zer{\Psi}$ satisfies $\norms{Gx^{\star} + v^{\star}}^2 = 0$ for some $v^{\star} \in Tx^{\star}$.
Therefore, if $(x^k, v^k)$ satisfies $\Expk{ \norms{Gx^k + v^k}^2} \leq \epsilon^2$ for some $v^k \in Tx^k$, then we can say that $x^k$ is an $\epsilon$-solution of \eqref{eq:NI}.
Alternatively, we can define the following forward-backward splitting (FBS) residual of \eqref{eq:NI}:
\vspace{-1ex}
\begin{equation*} 
\begin{array}{lcl}
\Gc_{\eta}x := \eta^{-1}(x - J_{\eta}(x - \eta Gx)),
\end{array}
\vspace{-0.5ex}
\end{equation*}
for any given $\eta > 0$.
It is well-known that $x^{\star} \in \zer{\Psi}$ iff $\Gc_{\eta}x^{\star} = 0$.
Hence, if $\Expk{ \norms{\Gc_{\eta}x^k}^2 } \leq \epsilon^2$, then $x^k$ is also called an $\epsilon$-solution of \eqref{eq:NI}.
One can easily prove that $\norms{\Gc_{\eta}x^k} \leq \norms{Gx^k + v^k}$ for any $v^k \in Tx^k$.
Clearly,  the former metric implies the latter one.
Therefore, it is sufficient to only certify $\Expk{ \norms{Gx^k + v^k}^2} \leq \epsilon^2$, which implies $\Expk{ \norms{\Gc_{\eta}x^k}^2 } \leq \epsilon^2$.


%%%% (d) Convergence Analysis.
\textbf{$\mathrm{(d)}$~Convergence analysis.}
For simplicity of our presentation, for a given $\gamma \in \left(\frac{1}{2}, 1\right)$, with $\rho$, $C$, and $\hat{C}$ in Definition~\ref{de:ub_SG_estimator}, we define the following two parameters:
\vspace{-0.75ex}
\begin{equation}\label{eq:VROG4NI_const}
\begin{array}{ll}
M := 4\gamma^2 + \frac{ 4\gamma}{1-\gamma} \cdot \frac{C + \hat{C}}{\rho} \quad \text{and} \quad \delta := \frac{\gamma(2\gamma-1)}{(3\gamma-1)\sqrt{M}}.
\end{array}
\vspace{-0.75ex}
\end{equation}
Then, Theorem~\ref{th:VROG4NI_convergence} below states the convergence of \eqref{eq:VROG4NI}, whose proof can be found in Supp. Doc.~\ref{apdx:sec:VROG4NI}.

%%% Theorem 4.1.
\begin{theorem}\label{th:VROG4NI_convergence}
Let us fix $\gamma \in \left(\frac{1}{2}, 1\right)$, and define $M$ and $\delta$ as in \eqref{eq:VROG4NI_const}.
Suppose that Assumptions~\ref{as:A0},~\ref{as:A0b}, ~\ref{as:A1}, and ~\ref{as:A2} hold for \eqref{eq:NI} for some $\kappa \geq 0$ such that $L\kappa <  \delta$.
Let $\sets{x^k}$ be generated by \eqref{eq:VROG4NI} using a fixed learning rate $\eta$ such that $\frac{(3\gamma-1)\kappa}{\gamma(2\gamma-1)} < \eta \leq \frac{1}{L\sqrt{M}}$.
Then, we have
\vspace{-1ex}
\begin{equation}\label{eq:VROG4NI_convergence_bound1}
\hspace{-0.25ex}
\begin{array}{lcl}
\dfrac{1}{K+1}\displaystyle\sum_{k=0}^{K} \Expk{ \norms{Gx^{k}+v^k}^2 }  & \leq &    \dfrac{ \hat{\Theta}_1 \hat{R}_0^2 }{\eta^2 ( K+1) }, \vspace{1ex}\\
\dfrac{1}{K+1}\displaystyle\sum_{k=0}^{K} \Expk{ \norms{x^k - x^{k-1}}^2 } & \leq &  \dfrac{ \hat{\Theta}_2 \hat{R}_0^2 }{ K + 1},
\end{array}
\hspace{-3.0ex}
\vspace{-1ex}
\end{equation}
where $\hat{R}_0^2$, $\hat{\Theta}_1$, and $\hat{\Theta}_2$ are respectively given by
\vspace{-1ex}
\begin{equation*}
\begin{array}{lcl}
\hat{R}_0^2 & := & \norms{x^0 - x^{\star}}^2 + \gamma^2\eta^2\norms{Gx^0 + v^0}^2, \vspace{0ex}\\
\hat{\Theta}_1 & := & \frac{(3\gamma-1)\eta }{(1-\gamma)[ \gamma(2\gamma-1)\eta - (3\gamma-1)\kappa]}, \vspace{0ex}\\
\hat{\Theta}_2 & := & \frac{4(3\gamma-1) \hat{R}_0^2 }{(1- \gamma)(1-ML^2\eta^2)}.
\end{array}
\vspace{-1ex}
\end{equation*}
\end{theorem}
The bounds in Theorem~\ref{th:VROG4NI_convergence} are similar to Theorem~\ref{th:VROG4NE_convergence}, but their proof relies on a new Lyapunov function.
Note that the condition on $L\kappa$ still depends on $\rho$ as $L\kappa  \leq \delta = \BigOs{\sqrt{\rho}}$.

%%%%%%%%%%%%%%%%%%%%%%%%%%%%%%%%%%%%%%%%%%%
%%%% 4.2. Oracle Complexity Bounds for Specific Estimators.
%%%%%%%%%%%%%%%%%%%%%%%%%%%%%%%%%%%%%%%%%%%
\beforesubsec
\subsection{Complexity of \ref{eq:VROG4NI} with SVRG and SAGA {\!\!\!}}\label{subsec:VROG2_oracle_complexity}
\aftersubsec
\vspace{0.5ex}
Similar to Section~\ref{sec:VROG4NE}, we can apply Theorem~\ref{th:VROG4NI_convergence} for the mini-batch SVRG estimator in Section~\ref{sec:VRS_Estimator}.

%%% Corollary 3.2.
\begin{corollary}\label{co:VROG2_SVRG_convergence1_v2}
Suppose that Assumptions~\ref{as:A0},~\ref{as:A0b} ~\ref{as:A1}, and ~\ref{as:A2} hold for \eqref{eq:NI} with $\kappa \geq 0$ as in Theorem \ref{th:VROG4NI_convergence}.
Let $\sets{x^k}$ be generated by \eqref{eq:VROG4NI} using the SVRG estimator \eqref{eq:SVRG_estimator}, $\gamma \in \left( \frac{1}{2}, 1\right)$, and $\eta := \frac{1}{L\sqrt{M}} \geq \frac{\sigma \sqrt{b}\mbf{p}}{L}$ with $\sigma := \frac{\sqrt{1-\gamma}}{2\sqrt{8+\gamma+7\gamma^2}}$, provided that $b\mbf{p}^2 \leq 1$.
Then, we have
\vspace{-0.75ex}
\begin{equation}\label{eq:svrg_bound_for_NI} 
%\begin{array}{ll}
\frac{1}{K+1}\sum_{k=0}^{K} \Expk{ \norms{Gx^{k} + v^k}^2 }   \leq     \frac{ \hat{\Theta}_1 L^2 \hat{R}^2_0 }{\sigma^2 b\mbf{p}^2 (K + 1)}, 
%\end{array}
\vspace{-0.75ex}
\end{equation}
where $\hat{R}^2_0 := \norms{x^0 - x^{\star}}^2 + \gamma^2\eta^2 \norms{Gx^0 + v^0}^2$.

For given $\epsilon > 0$, if we choose $\mbf{p} := n^{-1/3}$ and $b := \lfloor n^{2/3} \rfloor$, then \eqref{eq:VROG4NI} requires $\Tc_{G_i} := n + \big\lfloor \frac{4\Gamma L^2\hat{R}^2_0 n^{2/3}}{\varepsilon^2} \big\rfloor$ evaluations of $G_i$ and $\Tc_{T} = \big\lfloor  \frac{\Gamma L^2 \hat{R}_0^2 }{ \epsilon^2 } \big\rfloor$ evaluations of $J_{\gamma\eta T}$ to achieve $\frac{1}{K+1}\sum_{k=0}^{K} \Expk{ \norms{Gx^{k} \! + \! v^k}^2 } \leq \epsilon^2$, where $\Gamma := \frac{\hat{\Theta}_1}{\sigma^2}$.
\end{corollary}

Alternatively, we can apply Theorem~\ref{th:VROG4NI_convergence} to the mini-batch SAGA estimator \eqref{eq:SAGA_estimator} in Section~\ref{sec:VRS_Estimator}.

%%% Corollary 3.3.
\begin{corollary}\label{co:VROG2_SAGA_convergence1_v3}
Suppose that Assumptions~\ref{as:A0},~\ref{as:A0b}, ~\ref{as:A1}, and ~\ref{as:A2} hold for \eqref{eq:NI} with $\kappa \geq 0$ as in Theorem~\ref{th:VROG4NI_convergence}.
Let $\sets{x^k}$ be generated by \eqref{eq:VROG4NI} using the SAGA estimator \eqref{eq:SAGA_estimator}, $\gamma \in \left( \frac{1}{2}, 1\right)$, and $\eta := \frac{1}{L\sqrt{M}} \geq \frac{\sigma b^{3/2}}{nL}$ with $\sigma := \frac{\sqrt{1-\gamma}}{2\sqrt{\gamma(10+\gamma+7\gamma^2)}}$, provided that $1 \leq b \leq n^{2/3}$.
Then 
\vspace{-1ex}
\begin{equation}\label{eq:saga_bound_for_NI} 
%\begin{array}{ll}
\frac{1}{K+1}\sum_{k=0}^{K} \Expk{ \norms{Gx^{k} + v^k}^2 }   \leq     \frac{n^2 \hat{\Theta}_1 L^2 \hat{R}^2_0 }{\sigma^2b^3 (K + 1)},  
%\end{array}
\vspace{-1ex}
\end{equation}
where $\hat{R}^2_0 := \norms{x^0 - x^{\star}}^2 + \gamma^2\eta^2 \norms{Gx^0 + v^0}^2$.

For a given $\epsilon > 0$, if we choose $b := \lfloor n^{2/3} \rfloor$, then \eqref{eq:VROG4NI} requires $\Tc_{G_i} := n + \big\lfloor \frac{3\Gamma L^2\hat{R}^2_0 n^{2/3}}{\varepsilon^2} \big\rfloor$ evaluations of $G_i$ and $\Tc_{T} = \big\lfloor \frac{\Gamma L^2\hat{R}_0^2}{\epsilon^2} \big\rfloor$ evaluations of $J_{\gamma\eta T}$   to achieve $\frac{1}{K+1}\sum_{k=0}^{K} \Expk{ \norms{Gx^{k}+v^k}^2 } \leq \epsilon^2$, where $\Gamma := \frac{\hat{\Theta}_1}{\sigma^2}$.
\end{corollary}
Similar to Subsection~\ref{subsec:Oracle_complexity}, when $\gamma$, $n$, $b$, and $\mbf{p}$ are given, we can compute concrete values of the theoretical learning rate $\eta$ in both corollaries.
They are larger than the corresponding lower bounds given in these corollaries.

%%%%%%%%%%%%%%%%%%%%%%%%%%%%%%%%%%%%%%
%%% 6. Numerical Experiments.
%%%%%%%%%%%%%%%%%%%%%%%%%%%%%%%%%%%%%%
\beforesec
\section{Numerical Experiments}\label{sec:NumExps}
\aftersec
We provide two examples to illustrate \eqref{eq:VROG4NE} and \eqref{eq:VROG4NI} and compare them with other methods.

%%%% Example 1.
\beforesubsec
\subsection{Nonconvex-nonconcave minimax optimization}\label{subsec:example1}
\aftersubsec
We consider the following  nonconvex-nonconcave minimax optimization problem as a special case of \eqref{eq:minimax}:
\vspace{-0.75ex}
\begin{equation}\label{eq:minimax_exam1}
\hspace{-1ex}
\begin{array}{ll}
& {\displaystyle\min_{u \in \R^{p_1}}\max_{v\in\R^{p_2}}} \big\{ \Lc(u, v) := \varphi(u) +  \Hc(u, v) -  \psi(v) \big\}, 
\end{array}
\hspace{-1ex}
\vspace{-0.75ex}
\end{equation}
where $\Hc(u, v) := \frac{1}{n}\sum_{i=1}^n \Hc_i(u, v) = \frac{1}{n}\sum_{i=1}^n[ u^TA_iu +  u^TL_iv - v^TB_iv + b_i^{\top}u - c_i^{\top}v]$
such that $A_i \in \R^{p_1\times p_1}$ and $B_i \in \R^{p_2\times p_2}$ are symmetric matrices, $L_i \in \R^{p_1\times p_2}$, $b_i \in \R^{p_1}$, and $c_i \in \R^{p_2}$; $\varphi$ and $\psi$ are two proper, closed, and convex functions.
The optimality of \eqref{eq:minimax_exam1} becomes \eqref{eq:NI} (see Supp. Doc. \ref{apdx:sec:num_experiments}).
In our experiments, we choose  $A_i$ and $B_i$ to be not positive semidefinite such that Assumption~\ref{as:A2} holds.
Thus, \eqref{eq:minimax_exam1} is nonconvex-nonconcave.

We generate $A_i = Q_iD_iQ_i^T$ for a given orthonormal matrix $Q_i$ and a diagonal matrix $D_i$, where its entries $D_i^j$ are generated from standard normal distribution and clipped by $\max\sets{D_i^j, -0.1}$.
The matrix $B_i$ is also generated by the same way, while $L_i$, $b_i$, and $c_i$ are generated from standard normal distribution.
Hence, $\mathbf{G}$ in \eqref{eq:NI} is not symmetric and also not positive semidefinite.

\textbf{The unconstrained case.}
We implement three variants of \eqref{eq:VROG4NE}: \texttt{VFR-svrg} (double-loop SVRG), \texttt{LVFR-svrg} (loopless SVRG), \texttt{VFR-saga} (using SAGA estimator) in Python to solve \eqref{eq:minimax_exam1} when both $\varphi$ and $\psi$ are vanished, i.e. its optimality is a special case of \eqref{eq:NE}.
We also compare our methods with the deterministic optimistic gradient method (\texttt{OG}) in \cite{daskalakis2018training}, the variance-reduced FRBS scheme (\texttt{VFRBS}) in \cite{alacaoglu2022beyond}, and the variance-reduced extragradient algorithm (\texttt{VEG}) in \cite{alacaoglu2021stochastic}.
We select the parameters as suggested by our theory, while choosing appropriate parameters for \texttt{OG}, \texttt{VFRBS}, and \texttt{VEG}.
The details of this experiment, including generating data and specific choice of parameters, are given in Supp. Doc. \ref{apdx:sec:num_experiments}.

The relative residual norm $\norms{Gx^k}/\norms{Gx^0}$ against the number of epochs averaged on $10$ problem instances is revealed in Figure~\ref{fig:unconstr_minimax_exam1} for two datasets $(p, n) = (100, 5000)$ and $(p, n) = (200, 10000)$. 

\begin{figure}[ht]
\vspace{-2ex}
\begin{center}
\centerline{\includegraphics[width=0.8\columnwidth]{figures/unconstr_minimax_exam1_L}}
\centerline{\includegraphics[width=0.8\columnwidth]{figures/unconstr_minimax_exam1_R}}
\vspace{-1.5ex}
\caption{Comparison of $6$ algorithms to solve the unconstrained \eqref{eq:minimax_exam1} on $2$ experiments  (The average of $10$ runs).}
\label{fig:unconstr_minimax_exam1}
\end{center}
\vspace{-4ex}
\end{figure}

Clearly, with these experiments, three SVRG variants of our method \eqref{eq:VROG4NI} work well and significantly outperform other competitors. 
The \texttt{LVFR-svrg} variant of \eqref{eq:VROG4NI} seems to work best, while \texttt{VFRBS} and  \texttt{VEG} still cannot beat the deterministic algorithm \texttt{OG} in this example.

%%%%%%%%%%%%%%%%%%%%%%%%%%%%%%%%%%%%%%%%
%%% E.2.2. The constrained case
%%%%%%%%%%%%%%%%%%%%%%%%%%%%%%%%%%%%%%%%
\textbf{The constrained case.}
We now adding two simplex constraints $u \in \Delta_{p_1}$ and $v \in \Delta_{p_2}$ to \eqref{eq:minimax_exam1}, where $\Delta_p := \sets{u \in \R^p_{+}: \sum_{i=1}^pu_i = 1}$ is the standard simplex in $\R^p$. 
These constraints are common in bilinear games. 
To handle these constraints, we set $\varphi(u) := \delta_{\Delta_{p_1}}(u)$ and $\psi(v) := \delta_{\Delta_{p_2}}(v)$ as the indicators of $\Delta_{p_1}$ and $\Delta_{p_2}$, respectively.

Again, we run 6 algorithms for solving this constrained case of \eqref{eq:minimax_exam1}  using the same parameters as \textbf{the unconstrained case}.
We report the relative norm of the FBS residual $\norms{\Gc_{\eta}x^k}/\norms{\Gc_{\eta}x^0}$ against the number of epochs. 
The results are revealed in Figure~\ref{fig:bilinear_minimax_experiment2_mt} for two datasets $(p, n) = (100, 5000)$ and $(p, n) = (200, 10000)$. 
\begin{figure}[ht]
\vspace{-0ex}
\centering
\centerline{\includegraphics[width=0.8\columnwidth]{figures/constr_minimax_exam2_b_left}}
\centerline{\includegraphics[width=0.8\columnwidth]{figures/constr_minimax_exam2_b_right}}
\vspace{-1.5ex}
\caption{
The performance of $6$ algorithms to solve the constrained \eqref{eq:minimax_exam1} on $2$ experiments  (The average of $10$ runs).
}
\label{fig:bilinear_minimax_experiment2_mt}
\vspace{-3ex}
\end{figure}

Clearly, with these experiments, both SVRG variants of our method \eqref{eq:VROG4NI} work well and significantly outperform other competitors. 
The SVRG variant (\texttt{VFR-svrg}) of \eqref{eq:VROG4NI} seems to work best, while our \texttt{VFR-saga} has a similar performance as \texttt{VEG}.
Again, we also see that \texttt{VFRBS} tends to have a similar performance as \texttt{OG}.

%%%% Example 2.
\beforesubsec
\subsection{Logistic regression with ambiguous features}\label{subsec:example2}
\aftersubsec
We consider the following minimax optimization problem arising from a regularized logistic regression with ambiguous features (see Supp. Doc. \ref{apdx:sec:num_experiments} for more details):
\vspace{-1ex}
\begin{equation}\label{eq:logistic_reg_exam}
\hspace{-1ex}
\begin{array}{ll}
{\displaystyle\min_{w \in\R^d}} {\displaystyle\max_{z \in \R^m }} \Big\{ \mathcal{L}(w, z) \! := \! & \frac{1}{N} \sum_{i=1}^{N} \sum_{j=1}^m z_j  \ell ( \iprods{X_{ij}, w}, y_i ) \vspace{0.5ex}\\
& + {~} \tau R(w) - \delta_{\Delta_m}(z) \Big\},
\end{array}
\hspace{-5ex}
\vspace{-1ex}
\end{equation}
where  $\ell(\tau, s) := \log(1 + \exp(\tau)) - s\tau$ is the standard logistic loss,  $R(w) := \norms{w}_1$ is an $\ell_1$-norm regularizer, $\tau > 0$ is a regularization parameter, and $\delta_{\Delta_m}$ is the indicator of $\Delta_m$ to handle the constraint $z \in \Delta_m$.
Then, the optimality condition of \eqref{eq:logistic_reg_exam} can be cast into  \eqref{eq:NI}, where $x := [w, z]$.

We implement three variants of \eqref{eq:VROG4NI} to solve \eqref{eq:logistic_reg_exam}: \texttt{VFR-svrg}, \texttt{LVFR-svrg}, and \texttt{VFR-saga}.
We also compare our methods with \texttt{OG}, \texttt{VFRBS}, and \texttt{VEG} as in Subsection~\ref{subsec:example1}.
We cary out a mannual tuning procedure to select appropriate learning rates for all methods.
We test these algorithms on two real datasets: \texttt{a9a} (134 features and 3561 samples) and \texttt{w8a} (311 features and 45546 samples)  downloaded from \texttt{LIBSVM} \citep{CC01a}.
We first normalize the feature vector $\hat{X}_i$  and add a column of  all ones to address the bias term.
To generate ambiguous features, we take the nominal feature vector $\hat{X}_i$ and add a random noise generated from a normal distribution of zero mean and variance of $\sigma^2 = 0.5$.
In our test, we choose $\tau := 10^{-3}$ and $m := 10$.
The relative FBS residual norm $\norms{\Gc_{\eta}x^k}/\norms{\Gc_{\eta}x^0}$ against the epochs is plotted in Figure~\ref{fig:logistic_regression_exam2} for both datasets.

\begin{figure}[ht]
\vspace{-1ex}
\begin{center}
\centerline{\includegraphics[width=0.8\columnwidth]{figures/logistic_reg_exam2_L}}
\centerline{\includegraphics[width=0.8\columnwidth]{figures/logistic_reg_exam2_R}}
\vspace{-1.5ex}
\caption{Comparison of $6$ algorithms to solve \eqref{eq:logistic_reg_exam} on two real datasets: \texttt{a8a}  and \texttt{w8a}.}
\label{fig:logistic_regression_exam2}
\end{center}
\vspace{-3.5ex}
\end{figure}

%%{\onecolumn
%\begin{figure}[ht!]
%%\vspace{-1ex}
%\centering
%\includegraphics[width=0.9\textwidth]{figures/unconstr_minimax_exam1}
%\vspace{-2ex}
%\caption{
%Comparison of $6$ algorithms to solve \eqref{eq:minimax_exam1} on $2$ experiments  (The average of $10$ runs).
%}
%\label{fig:unconstr_minimax_exam1}
%\vspace{-2.5ex}
%\end{figure}

%\begin{figure}[ht!]
%%\vspace{-1ex}
%\centering
%\includegraphics[width=0.9\textwidth]{figures/logistic_reg_exam2}
%\vspace{-1.25ex}
%\caption{
%Comparison of $6$ algorithms to solve \eqref{eq:logistic_reg_exam} on two real datasets: \texttt{a8a}  and \texttt{w8a}.
%}
%\label{fig:logistic_regression_exam2}
%\vspace{-1ex}
%\end{figure}
%%}

As we can observe from Figure~\ref{fig:logistic_regression_exam2} that three variants \texttt{VFR-svrg}, \texttt{LVFR-svrg}, and \texttt{VFR-saga} have similar performance and are better than their competitors.
Among three competitors, \texttt{VFRBS} still works well, and is much better than \texttt{OG} and \texttt{VEG}.
The deterministic method, \texttt{OG}, is the worst one in terms of oracle complexity.
In this test, \texttt{VEG} has a larger learning rate than ours and \texttt{VFRBS}.

%%%%%%%%%%%%%%%%%%%%%%%%%%%%%%%%%%%%%%%%%%%%%%%%%%%%%%%%%%%%
%%% 7. Conclusions
%%%%%%%%%%%%%%%%%%%%%%%%%%%%%%%%%%%%%%%%%%%%%%%%%%%%%%%%%%%%
\beforesec
\section{Conclusions}\label{sec:conclusions}
\aftersec
\vspace{0.25ex}
We develop two new variance-reduced algorithms based on the forward-reflected-backward splitting method to tackle both root-finding problems \eqref{eq:NE} and  \eqref{eq:NI}.
These methods encompass both SVRG and SAGA estimators as special cases.
By carefully selecting the parameters, our algorithms achieve the state-of-the-art oracle complexity for attaining an $\epsilon$-solution, matching the state-of-the-art complexity bounds observed in nonconvex optimization methods using SVRG and SAGA. 
While the first scheme resembles a stochastic variant of the optimistic gradient method, the second one is entirely novel and distinct from existing approaches, even their deterministic counterparts.
We have validated our methods through numerical examples, and the results demonstrate promising performance compared to existing techniques under careful parameter selections.


%We have developed two novel classes of variance-reduced forward-reflected-based methods to solve \eqref{eq:NE} and \eqref{eq:NI}, respectively, which cover both SVRG and SAGA estimators as special cases.
%Under appropriate choices of parameters, our method achieve the best-known oracle complexity to reach an $\epsilon$-solution as often seen from nonconvex optimization methods.
%While our first scheme can be viewed as a stochastic variant of the optimistic gradient method, the second algorithm is new, and different from existing ones, even for deterministic counterparts.
%We have tested our methods on different numerical examples, and our numerical results show that under the same parameter configuration, our methods show some promising performance compared to their competitors. 


%%%%%%%%%%%%%%%%%%%%%%%%%%%%%%%%%%%%%%%%%%%%%%%%%%%%%%%%%%%%
%%% END OF THE MAIN PAPER.
%%%%%%%%%%%%%%%%%%%%%%%%%%%%%%%%%%%%%%%%%%%%%%%%%%%%%%%%%%%%


% In the unusual situation where you want a paper to appear in the
% references without citing it in the main text, use \nocite
\section*{Impact Statement}
This paper proposes new algorithms with rigorous convergence guarantees and complexity estimates for solving a broad class of large-scale problems. 
These problems cover many fundamental challenges and applications in optimization, machine learning, and related fields as special cases. 
We believe that our new algorithms have the potential to make a significant impact in machine learning and related areas. Additionally, there are various potential societal consequences of our work, though none that we feel require specific emphasis at this time.

%This paper proposes new algorithms  with rigorous convergence guarantees and complexity estimates for solving a broad class of large-scale problems, which covers many fundamental problems and applications in optimization, machine learning, and related fields as special cases.
%We believe that our new algorithms can potentially make great impact in machine learning and related areas. 
%In addition, there are many potential societal consequences of our work, none which we feel must be specifically highlighted here. 

\section*{Acknowledgements}
This work was partly supported by the National Science Foundation (NSF): NSF-RTG grant No. NSF DMS-2134107 and the Office of Naval Research (ONR), grant No. N00014-23-1-2588.

\nocite{langley00}

%\bibliography{example_paper}
%\bibliography{/Users/quoctd/Dropbox/E-Books/tran_bibtex_new.bib}
\bibliographystyle{icml2025}
\begin{thebibliography}{74}
\providecommand{\natexlab}[1]{#1}
\providecommand{\url}[1]{\texttt{#1}}
\expandafter\ifx\csname urlstyle\endcsname\relax
  \providecommand{\doi}[1]{doi: #1}\else
  \providecommand{\doi}{doi: \begingroup \urlstyle{rm}\Url}\fi

\bibitem[Alacaoglu \& Malitsky(2022)Alacaoglu and
  Malitsky]{alacaoglu2021stochastic}
Alacaoglu, A. and Malitsky, Y.
\newblock Stochastic variance reduction for variational inequality methods.
\newblock In \emph{Conference on Learning Theory}, pp.\  778--816. PMLR, 2022.

\bibitem[Alacaoglu et~al.(2021)Alacaoglu, Malitsky, and
  Cevher]{alacaoglu2021forward}
Alacaoglu, A., Malitsky, Y., and Cevher, V.
\newblock Forward-reflected-backward method with variance reduction.
\newblock \emph{Comput. Optim. Appl.}, 80\penalty0 (2):\penalty0 321--346,
  2021.

\bibitem[Alacaoglu et~al.(2023)Alacaoglu, B{\"o}hm, and
  Malitsky]{alacaoglu2022beyond}
Alacaoglu, A., B{\"o}hm, A., and Malitsky, Y.
\newblock Beyond the golden ratio for variational inequality algorithms.
\newblock \emph{J. Mach. Learn. Res}, 24\penalty0 (172):\penalty0 1--33, 2023.

\bibitem[Allen-Zhu \& Hazan(2016)Allen-Zhu and Hazan]{allen2016variance}
Allen-Zhu, Z. and Hazan, E.
\newblock Variance reduction for faster non-convex optimization.
\newblock In \emph{International conference on machine learning}, pp.\
  699--707. PMLR, 2016.

\bibitem[Arjovsky et~al.(2017)Arjovsky, Chintala, and
  Bottou]{arjovsky2017wasserstein}
Arjovsky, M., Chintala, S., and Bottou, L.
\newblock Wasserstein generative adversarial networks.
\newblock In \emph{International Conference on Machine Learning}, pp.\
  214--223, 2017.

\bibitem[Azar et~al.(2017)Azar, Osband, and Munos]{azar2017minimax}
Azar, M.~G., Osband, I., and Munos, R.
\newblock Minimax regret bounds for reinforcement learning.
\newblock In \emph{International Conference on Machine Learning}, pp.\
  263--272. PMLR, 2017.

\bibitem[Bauschke \& Combettes(2017)Bauschke and Combettes]{Bauschke2011}
Bauschke, H.~H. and Combettes, P.
\newblock \emph{Convex analysis and monotone operators theory in {H}ilbert
  spaces}.
\newblock Springer-Verlag, 2nd edition, 2017.

\bibitem[Bauschke et~al.(2020)Bauschke, Moursi, and
  Wang]{bauschke2020generalized}
Bauschke, H.~H., Moursi, W.~M., and Wang, X.
\newblock Generalized monotone operators and their averaged resolvents.
\newblock \emph{Math. Program.}, pp.\  1--20, 2020.

\bibitem[Ben-Tal et~al.(2009)Ben-Tal, El~Ghaoui, and Nemirovski]{Ben-Tal2009}
Ben-Tal, A., El~Ghaoui, L., and Nemirovski, A.
\newblock \emph{{R}obust optimization}.
\newblock Princeton University Press, 2009.

\bibitem[Bertsimas \& Caramanis(2011)Bertsimas and Caramanis]{Bertsimas2011}
Bertsimas, D.~Brown, D. and Caramanis, C.
\newblock {T}heory and {A}pplications of {R}obust {O}ptimization.
\newblock \emph{SIAM Review}, 53\penalty0 (3):\penalty0 464--501, 2011.

\bibitem[Beznosikov et~al.(2023)Beznosikov, Gorbunov, Berard, and
  Loizou]{beznosikov2023stochastic}
Beznosikov, A., Gorbunov, E., Berard, H., and Loizou, N.
\newblock Stochastic gradient descent-ascent: Unified theory and new efficient
  methods.
\newblock In \emph{International Conference on Artificial Intelligence and
  Statistics}, pp.\  172--235. PMLR, 2023.

\bibitem[Bhatia \& Sridharan(2020)Bhatia and Sridharan]{bhatia2020online}
Bhatia, K. and Sridharan, K.
\newblock Online learning with dynamics: {A} minimax perspective.
\newblock \emph{Advances in Neural Information Processing Systems},
  33:\penalty0 15020--15030, 2020.

\bibitem[B{\"o}hm(2022)]{bohm2022solving}
B{\"o}hm, A.
\newblock Solving nonconvex-nonconcave min-max problems exhibiting weak {M}inty
  solutions.
\newblock \emph{Transactions on Machine Learning Research}, 2022.

\bibitem[Bot et~al.(2019)Bot, Mertikopoulos, Staudigl, and
  Vuong]{bot2019forward}
Bot, R.~I., Mertikopoulos, P., Staudigl, M., and Vuong, P.~T.
\newblock Forward-backward-forward methods with variance reduction for
  stochastic variational inequalities.
\newblock \emph{arXiv preprint arXiv:1902.03355}, 2019.

\bibitem[Burachik \& Iusem(2008)Burachik and Iusem]{reginaset2008}
Burachik, R.~S. and Iusem, A.
\newblock \emph{Set-Valued Mappings and Enlargements of Monotone Operators}.
\newblock New York: Springer, 2008.

\bibitem[Cai et~al.(2022)Cai, Song, Guzm\'{a}n, and
  Diakonikolas]{cai2022stochastic}
Cai, X., Song, C., Guzm\'{a}n, C., and Diakonikolas, J.
\newblock A stochastic {H}alpern iteration with variance reduction for
  stochastic monotone inclusion problems.
\newblock In \emph{Proceedings of the 12th International Conference on Learning
  Representations (ICLR 2022)}, 2022.
\newblock URL \url{https://openreview.net/forum?id=BRZos-8TpCf}.

\bibitem[Cai et~al.(2024)Cai, Alacaoglu, and Diakonikolas]{cai2023variance}
Cai, X., Alacaoglu, A., and Diakonikolas, J.
\newblock Variance reduced halpern iteration for finite-sum monotone
  inclusions.
\newblock In \emph{The 12th International Conference on Learning
  Representations (ICLR)}, pp.\  1--33, 2024.

\bibitem[Cai \& Zheng(2023)Cai and Zheng]{cai2022baccelerated}
Cai, Y. and Zheng, W.
\newblock {A}ccelerated {S}ingle-{C}all {M}ethods for {C}onstrained {M}in-{M}ax
  {O}ptimization.
\newblock In \emph{The 11th International Conference on Learning
  Representations, ICLR 2023}. The Eleventh International Conference on
  Learning Representations, ICLR 2023, 2023.

\bibitem[Carmon et~al.(2019)Carmon, Jin, Sidford, and Tian]{carmon2019variance}
Carmon, Y., Jin, Y., Sidford, A., and Tian, K.
\newblock Variance reduction for matrix games.
\newblock \emph{Advances in Neural Information Processing Systems}, 32, 2019.

\bibitem[Cevher \& V{\~u}(2021)Cevher and V{\~u}]{cevher2021reflected}
Cevher, V. and V{\~u}, B.
\newblock A reflected forward-backward splitting method for monotone inclusions
  involving {L}ipschitzian operators.
\newblock \emph{Set-Valued and Variational Analysis}, 29\penalty0 (1):\penalty0
  163--174, 2021.

\bibitem[Chang \& Lin(2011)Chang and Lin]{CC01a}
Chang, C.-C. and Lin, C.-J.
\newblock {LIBSVM}: {A} library for {S}upport {V}ector {M}achines.
\newblock \emph{ACM Transactions on Intelligent Systems and Technology},
  2:\penalty0 27:1--27:27, 2011.

\bibitem[Chavdarova et~al.(2019)Chavdarova, Gidel, Fleuret, and
  Lacoste-Julien]{chavdarova2019reducing}
Chavdarova, T., Gidel, G., Fleuret, F., and Lacoste-Julien, S.
\newblock Reducing noise in gan training with variance reduced extragradient.
\newblock \emph{Advances in Neural Information Processing Systems},
  32:\penalty0 393--403, 2019.

\bibitem[Combettes \& Pennanen(2004)Combettes and
  Pennanen]{combettes2004proximal}
Combettes, P.~L. and Pennanen, T.
\newblock Proximal methods for cohypomonotone operators.
\newblock \emph{SIAM J. Control Optim.}, 43\penalty0 (2):\penalty0 731--742,
  2004.

\bibitem[Cui \& Shanbhag(2021)Cui and Shanbhag]{cui2021analysis}
Cui, S. and Shanbhag, U.
\newblock On the analysis of variance-reduced and randomized projection
  variants of single projection schemes for monotone stochastic variational
  inequality problems.
\newblock \emph{Set-Valued and Variational Analysis}, 29\penalty0 (2):\penalty0
  453--499, 2021.

\bibitem[Cutkosky \& Orabona(2019)Cutkosky and Orabona]{Cutkosky2019}
Cutkosky, A. and Orabona, F.
\newblock Momentum-based variance reduction in non-convex {SGD}.
\newblock In \emph{Advances in Neural Information Processing Systems}, pp.\
  15210--15219, 2019.

\bibitem[Daskalakis et~al.(2018)Daskalakis, Ilyas, Syrgkanis, and
  Zeng]{daskalakis2018training}
Daskalakis, C., Ilyas, A., Syrgkanis, V., and Zeng, H.
\newblock Training {GANs} with {O}ptimism.
\newblock In \emph{International Conference on Learning Representations (ICLR
  2018)}, 2018.

\bibitem[Davis(2022)]{davis2022variance}
Davis, D.
\newblock Variance reduction for root-finding problems.
\newblock \emph{Math. Program.}, pp.\  1--36, 2022.

\bibitem[Defazio et~al.(2014)Defazio, Bach, and Lacoste-Julien]{Defazio2014}
Defazio, A., Bach, F., and Lacoste-Julien, S.
\newblock {SAGA}: {A} fast incremental gradient method with support for
  non-strongly convex composite objectives.
\newblock In \emph{Advances in Neural Information Processing Systems (NIPS)},
  pp.\  1646--1654, 2014.

\bibitem[Diakonikolas(2020)]{diakonikolas2020halpern}
Diakonikolas, J.
\newblock {H}alpern iteration for near-optimal and parameter-free monotone
  inclusion and strong solutions to variational inequalities.
\newblock In \emph{Conference on Learning Theory}, pp.\  1428--1451. PMLR,
  2020.

\bibitem[Diakonikolas et~al.(2021)Diakonikolas, Daskalakis, and
  Jordan]{diakonikolas2021efficient}
Diakonikolas, J., Daskalakis, C., and Jordan, M.
\newblock Efficient methods for structured nonconvex-nonconcave min-max
  optimization.
\newblock In \emph{International Conference on Artificial Intelligence and
  Statistics}, pp.\  2746--2754. PMLR, 2021.

\bibitem[Driggs et~al.(2020)Driggs, Ehrhardt, and
  Sch{\"o}nlieb]{driggs2019accelerating}
Driggs, D., Ehrhardt, M.~J., and Sch{\"o}nlieb, C.-B.
\newblock Accelerating variance-reduced stochastic gradient methods.
\newblock \emph{Math. Program.}, (online first):\penalty0 1--45, 2020.

\bibitem[Facchinei \& Pang(2003)Facchinei and Pang]{Facchinei2003}
Facchinei, F. and Pang, J.-S.
\newblock \emph{Finite-dimensional variational inequalities and complementarity
  problems}, volume 1-2.
\newblock Springer-Verlag, 2003.

\bibitem[Goodfellow et~al.(2014)Goodfellow, Pouget-Abadie, Mirza, Xu,
  Warde-Farley, Ozair, Courville, and Bengio]{goodfellow2014generative}
Goodfellow, I., Pouget-Abadie, J., Mirza, M., Xu, B., Warde-Farley, D., Ozair,
  S., Courville, A., and Bengio, Y.
\newblock Generative adversarial nets.
\newblock In \emph{Advances in neural information processing systems}, pp.\
  2672--2680, 2014.

\bibitem[Gorbunov et~al.(2022)Gorbunov, Berard, Gidel, and
  Loizou]{gorbunov2022stochastic}
Gorbunov, E., Berard, H., Gidel, G., and Loizou, N.
\newblock Stochastic extragradient: General analysis and improved rates.
\newblock In \emph{International Conference on Artificial Intelligence and
  Statistics}, pp.\  7865--7901. PMLR, 2022.

\bibitem[Gower et~al.(2021)Gower, Richt{\'a}rik, and Bach]{gower2021stochastic}
Gower, R.~M., Richt{\'a}rik, P., and Bach, F.
\newblock Stochastic quasi-gradient methods: {V}ariance reduction via
  {J}acobian sketching.
\newblock \emph{Math. Program.}, 188\penalty0 (1):\penalty0 135--192, 2021.

\bibitem[Hanzely et~al.(2018)Hanzely, Mishchenko, and
  Richt{\'a}rik]{hanzely2018sega}
Hanzely, F., Mishchenko, K., and Richt{\'a}rik, P.
\newblock {SEGA}: {V}ariance reduction via gradient sketching.
\newblock In \emph{Advances in Neural Information Processing Systems}, pp.\
  2082--2093, 2018.

\bibitem[Huang et~al.(2022)Huang, Wang, and Zhang]{huang2022accelerated}
Huang, K., Wang, N., and Zhang, S.
\newblock An accelerated variance reduced extra-point approach to finite-sum vi
  and optimization.
\newblock \emph{arXiv preprint arXiv:2211.03269}, 2022.

\bibitem[Iusem et~al.(2017)Iusem, Jofr{\'e}, Oliveira, and
  Thompson]{iusem2017extragradient}
Iusem, A.~N., Jofr{\'e}, A., Oliveira, R.~I., and Thompson, P.
\newblock Extragradient method with variance reduction for stochastic
  variational inequalities.
\newblock \emph{SIAM J. Optim.}, 27\penalty0 (2):\penalty0 686--724, 2017.

\bibitem[Johnson \& Zhang(2013)Johnson and Zhang]{SVRG}
Johnson, R. and Zhang, T.
\newblock Accelerating stochastic gradient descent using predictive variance
  reduction.
\newblock In \emph{NIPS}, pp.\  315--323, 2013.

\bibitem[Juditsky et~al.(2011)Juditsky, Nemirovski, and
  Tauvel]{juditsky2011solving}
Juditsky, A., Nemirovski, A., and Tauvel, C.
\newblock Solving variational inequalities with stochastic mirror-prox
  algorithm.
\newblock \emph{Stochastic Systems}, 1\penalty0 (1):\penalty0 17--58, 2011.

\bibitem[Kannan \& Shanbhag(2019)Kannan and Shanbhag]{kannan2019optimal}
Kannan, A. and Shanbhag, U.~V.
\newblock Optimal stochastic extragradient schemes for pseudomonotone
  stochastic variational inequality problems and their variants.
\newblock \emph{Comput. Optim. Appl.}, 74\penalty0 (3):\penalty0 779--820,
  2019.

\bibitem[Kotsalis et~al.(2022)Kotsalis, Lan, and Li]{kotsalis2022simple}
Kotsalis, G., Lan, G., and Li, T.
\newblock Simple and optimal methods for stochastic variational inequalities,
  i: operator extrapolation.
\newblock \emph{SIAM J. Optim.}, 32\penalty0 (3):\penalty0 2041--2073, 2022.

\bibitem[Kovalev et~al.(2020)Kovalev, Horvath, and Richtarik]{kovalev2019don}
Kovalev, D., Horvath, S., and Richtarik, P.
\newblock {D}on't jump through hoops and remove those loops: {SVRG} and
  {K}atyusha are better without the outer loop.
\newblock In \emph{Algorithmic Learning Theory}, pp.\  451--467. PMLR, 2020.

\bibitem[Lee \& Kim(2021)Lee and Kim]{lee2021fast}
Lee, S. and Kim, D.
\newblock Fast extra gradient methods for smooth structured
  nonconvex-nonconcave minimax problems.
\newblock \emph{Thirty-fifth Conference on Neural Information Processing
  Systems (NeurIPs2021)}, 2021.

\bibitem[Levy et~al.(2020)Levy, Carmon, Duchi, and Sidford]{levy2020large}
Levy, D., Carmon, Y., Duchi, J.~C., and Sidford, A.
\newblock Large-scale methods for distributionally robust optimization.
\newblock \emph{Advances in Neural Information Processing Systems},
  33:\penalty0 8847--8860, 2020.

\bibitem[Loizou et~al.(2021)Loizou, Berard, Gidel, Mitliagkas, and
  Lacoste-Julien]{loizou2021stochastic}
Loizou, N., Berard, H., Gidel, G., Mitliagkas, I., and Lacoste-Julien, S.
\newblock Stochastic gradient descent-ascent and consensus optimization for
  smooth games: {C}onvergence analysis under expected co-coercivity.
\newblock \emph{Advances in Neural Information Processing Systems},
  34:\penalty0 19095--19108, 2021.

\bibitem[Luo \& Tran-Dinh(2022)Luo and Tran-Dinh]{luo2022last}
Luo, Y. and Tran-Dinh, Q.
\newblock Extragradient-type methods for co-monotone root-finding problems.
\newblock \emph{(UNC-STOR Technical Report)}, 2022.

\bibitem[Madry et~al.(2018)Madry, Makelov, Schmidt, Tsipras, and
  Vladu]{madry2018towards}
Madry, A., Makelov, A., Schmidt, L., Tsipras, D., and Vladu, A.
\newblock Towards deep learning models resistant to adversarial attacks.
\newblock In \emph{International Conference on Learning Representations}, 2018.

\bibitem[Malitsky(2015)]{malitsky2015projected}
Malitsky, Y.
\newblock Projected reflected gradient methods for monotone variational
  inequalities.
\newblock \emph{SIAM J. Optim.}, 25\penalty0 (1):\penalty0 502--520, 2015.

\bibitem[Malitsky \& Tam(2020)Malitsky and Tam]{malitsky2020forward}
Malitsky, Y. and Tam, M.~K.
\newblock A forward-backward splitting method for monotone inclusions without
  cocoercivity.
\newblock \emph{SIAM J. Optim.}, 30\penalty0 (2):\penalty0 1451--1472, 2020.

\bibitem[Namkoong \& Duchi(2016)Namkoong and Duchi]{namkoong2016stochastic}
Namkoong, H. and Duchi, J.
\newblock Stochastic gradient methods for distributionally robust optimization
  with f-divergences.
\newblock \emph{Advances in neural information processing systems}, 29, 2016.

\bibitem[Nesterov(2007)]{Nesterov2007a}
Nesterov, Y.
\newblock Dual extrapolation and its applications to solving variational
  inequalities and related problems.
\newblock \emph{Math. Program.}, 109\penalty0 (2--3):\penalty0 319--344, 2007.

\bibitem[Nguyen et~al.(2017)Nguyen, Liu, Scheinberg, and
  Tak{\'a}{\v{c}}]{nguyen2017sarah}
Nguyen, L.~M., Liu, J., Scheinberg, K., and Tak{\'a}{\v{c}}, M.
\newblock {SARAH}: A novel method for machine learning problems using
  stochastic recursive gradient.
\newblock In \emph{Proceedings of the 34th International Conference on Machine
  Learning}, pp.\  2613--2621, 2017.

\bibitem[Palaniappan \& Bach(2016)Palaniappan and
  Bach]{palaniappan2016stochastic}
Palaniappan, B. and Bach, F.
\newblock Stochastic variance reduction methods for saddle-point problems.
\newblock In \emph{Advances in Neural Information Processing Systems}, pp.\
  1416--1424, 2016.

\bibitem[Pethick et~al.(2022)Pethick, Patrinos, Fercoq, and
  Cevher]{pethick2022escaping}
Pethick, T., Patrinos, P., Fercoq, O., and Cevher, V.
\newblock Escaping limit cycles: {G}lobal convergence for constrained
  nonconvex-nonconcave minimax problems.
\newblock In \emph{International Conference on Learning Representations}, 2022.

\bibitem[Pethick et~al.(2023)Pethick, Fercoq, Latafat, Patrinos, and
  Cevher]{pethick2023solving}
Pethick, T., Fercoq, O., Latafat, P., Patrinos, P., and Cevher, V.
\newblock Solving stochastic weak {M}inty variational inequalities without
  increasing batch size.
\newblock In \emph{Proceedings of International Conference on Learning
  Representations (ICLR)}, pp.\  1--34, 2023.

\bibitem[Phelps(2009)]{phelps2009convex}
Phelps, R.~R.
\newblock \emph{Convex functions, monotone operators and differentiability},
  volume 1364.
\newblock Springer, 2009.

\bibitem[Popov(1980)]{popov1980modification}
Popov, L.~D.
\newblock A modification of the {A}rrow-{H}urwicz method for search of saddle
  points.
\newblock \emph{Math. notes of the Academy of Sciences of the USSR},
  28\penalty0 (5):\penalty0 845--848, 1980.

\bibitem[Reddi et~al.(2016{\natexlab{a}})Reddi, Hefny, Sra, P{\'{o}}czos, and
  Smola]{Reddi2016a}
Reddi, S.~J., Hefny, A., Sra, S., P{\'{o}}czos, B., and Smola, A.~J.
\newblock Stochastic variance reduction for nonconvex optimization.
\newblock In \emph{ICML}, pp.\  314--323, 2016{\natexlab{a}}.

\bibitem[Reddi et~al.(2016{\natexlab{b}})Reddi, Sra, P{\'o}czos, and
  Smola]{reddi2016fast}
Reddi, S.~J., Sra, S., P{\'o}czos, B., and Smola, A.
\newblock Fast incremental method for smooth nonconvex optimization.
\newblock In \emph{2016 IEEE 55th Conference on Decision and Control (CDC)},
  pp.\  1971--1977. IEEE, 2016{\natexlab{b}}.

\bibitem[Robbins \& Monro(1951)Robbins and Monro]{RM1951}
Robbins, H. and Monro, S.
\newblock A stochastic approximation method.
\newblock \emph{The Annals of Mathematical Statistics}, 22\penalty0
  (3):\penalty0 400--407, 1951.

\bibitem[Rockafellar \& Wets(1997)Rockafellar and Wets]{Rockafellar1997}
Rockafellar, R. and Wets, R. J.-B.
\newblock \emph{{V}ariational {A}nalysis}.
\newblock Springer-Verlag, 1997.

\bibitem[Ryu \& Yin(2022)Ryu and Yin]{ryu2022large}
Ryu, E. and Yin, W.
\newblock \emph{Large-scale convex optimization: {A}lgorithms \& analyses via
  monotone operators}.
\newblock Cambridge University Press, 2022.

\bibitem[Ryu \& Boyd(2016)Ryu and Boyd]{ryu2016primer}
Ryu, E.~K. and Boyd, S.
\newblock Primer on monotone operator methods.
\newblock \emph{Appl. Comput. Math}, 15\penalty0 (1):\penalty0 3--43, 2016.

\bibitem[Tran-Dinh(2023{\natexlab{a}})]{tran2023extragradient}
Tran-Dinh, Q.
\newblock {E}xtragradient-{T}ype {M}ethods with
  $\mathcal{O}(1/k)$-{C}onvergence {R}ates for {C}o-{H}ypomonotone
  {I}nclusions.
\newblock \emph{J. Global Optim.}, pp.\  1--25, 2023{\natexlab{a}}.

\bibitem[Tran-Dinh(2023{\natexlab{b}})]{tran2023sublinear}
Tran-Dinh, Q.
\newblock {S}ublinear {C}onvergence {R}ates of {E}xtragradient-{T}ype
  {M}ethods: {A} {S}urvey on {C}lassical and {R}ecent {D}evelopments.
\newblock \emph{arXiv preprint arXiv:2303.17192}, 2023{\natexlab{b}}.

\bibitem[Tran-Dinh(2024)]{tran2024accelerated}
Tran-Dinh, Q.
\newblock {A}ccelerated {V}ariance-{R}educed {F}orward-{R}eflected {M}ethods
  for {R}oot-{F}inding {P}roblems.
\newblock \emph{arXiv preprint arXiv:2406.02413}, 2024.

\bibitem[Tran-Dinh \& Luo(2025)Tran-Dinh and Luo]{tran2022accelerated}
Tran-Dinh, Q. and Luo, Y.
\newblock {R}andomized {B}lock-{C}oordinate {O}ptimistic {G}radient
  {A}lgorithms for {R}oot-{F}inding {P}roblems.
\newblock \emph{Math. Oper. Res.}, in press, 2025.

\bibitem[Tran-Dinh et~al.(2019)Tran-Dinh, Pham, Phan, and
  Nguyen]{Tran-Dinh2019}
Tran-Dinh, Q., Pham, H.~N., Phan, T.~D., and Nguyen, M.~L.
\newblock Hybrid stochastic gradient descent algorithms for stochastic
  nonconvex optimization.
\newblock \emph{Preprint: arXiv:1905.05920}, 2019.

\bibitem[Tran-Dinh et~al.(2022)Tran-Dinh, Pham, Phan, and
  Nguyen]{Tran-Dinh2019a}
Tran-Dinh, Q., Pham, N.~H., Phan, D.~T., and Nguyen, L.~M.
\newblock A hybrid stochastic optimization framework for stochastic composite
  nonconvex optimization.
\newblock \emph{Math. Program.}, 191:\penalty0 1005--1071, 2022.

\bibitem[Yousefian et~al.(2018)Yousefian, Nedi{\'c}, and
  Shanbhag]{yousefian2018stochastic}
Yousefian, F., Nedi{\'c}, A., and Shanbhag, U.~V.
\newblock On stochastic mirror-prox algorithms for stochastic cartesian
  variational inequalities: {R}andomized block coordinate and optimal averaging
  schemes.
\newblock \emph{Set-Valued and Variational Analysis}, 26:\penalty0 789--819,
  2018.

\bibitem[Yu et~al.(2022)Yu, Lin, Mazumdar, and Jordan]{yu2022fast}
Yu, Y., Lin, T., Mazumdar, E.~V., and Jordan, M.
\newblock Fast distributionally robust learning with variance-reduced min-max
  optimization.
\newblock In \emph{International Conference on Artificial Intelligence and
  Statistics}, pp.\  1219--1250. PMLR, 2022.

\bibitem[Zhang et~al.(2021)Zhang, Yang, and Ba{\c{s}}ar]{zhang2021multi}
Zhang, K., Yang, Z., and Ba{\c{s}}ar, T.
\newblock Multi-agent reinforcement learning: {A} selective overview of
  theories and algorithms.
\newblock \emph{Handbook of reinforcement learning and control}, pp.\
  321--384, 2021.

\bibitem[Zhou et~al.(2018)Zhou, Xu, and Gu]{zhou2018stochastic}
Zhou, D., Xu, P., and Gu, Q.
\newblock Stochastic nested variance reduction for nonconvex optimization.
\newblock In \emph{Proceedings of the 32nd International Conference on Neural
  Information Processing Systems}, pp.\  3925--3936. Curran Associates Inc.,
  2018.

\end{thebibliography}

%\bibliography{/Users/quoc/Dropbox/E-Books/tran_bibtex_new.bib}

%\end{document}
%%%%%%%%%%%%%%%%%%%%%%%%%%%%%%%%%%%%%%%%%%%%%%%%%%%%%%%%%%%%%%%%%%%%%%%%%%%%%%%
%%%%%%%%%%%%%%%%%%%%%%%%%%%%%%%%%%%%%%%%%%%%%%%%%%%%%%%%%%%%%%%%%%%%%%%%%%%%%%%
% APPENDIX
%%%%%%%%%%%%%%%%%%%%%%%%%%%%%%%%%%%%%%%%%%%%%%%%%%%%%%%%%%%%%%%%%%%%%%%%%%%%%%%
%%%%%%%%%%%%%%%%%%%%%%%%%%%%%%%%%%%%%%%%%%%%%%%%%%%%%%%%%%%%%%%%%%%%%%%%%%%%%%%
\newpage
\appendix
\onecolumn

%%%%%%%%%%%%%%%%%%%%%%%%%%%%%%%%%%%%%%%%%%%%%%%%%%%%%%%
%%% TITLE
%%%%%%%%%%%%%%%%%%%%%%%%%%%%%%%%%%%%%%%%%%%%%%%%%%%%%%%
\vspace{-3ex}
\begin{center}
\hrulefill\vspace{1ex}\\
{\Large Supplementary Document:} \vspace{0.5ex}\\
\textbf{\Large
Variance-Reduced Forward-Reflected-Backward Splitting Methods for \vspace{0.75ex}\\ Nonmonotone Generalized Equations
}\\
\hrulefill
\end{center}
\vspace{2ex}

%%%%%%%%%%%%%%%%%%%%%%%%%%%%%%%%%%%%%%%%%%%%%%%%%%%%%%%
%%% The supplementary document
%%%%%%%%%%%%%%%%%%%%%%%%%%%%%%%%%%%%%%%%%%%%%%%%%%%%%%%
Due to space limit, some parts of our algorithmic construction and theory are not described in detail and motivated in the main text.
This supplementary document aims at providing more details of the algorithmic construction, motivation, related work, technical proofs, and additional experiments related to our methods.

%%%%%%%%%%%%%%%%%%%%%%%%%%%%%%%%%%%%%%%%%%%%%%%%%%%%%%%
%%%% A. The Proof of Technical Results in Section 2.
%%%%%%%%%%%%%%%%%%%%%%%%%%%%%%%%%%%%%%%%%%%%%%%%%%%%%%%
\beforesec
\section{Further Discussion of Related Work and Assumptions}\label{apdix:sec:further_related_work}
\aftersec
Let us further expand our discussion of related work in the main text.
Then, we show that our assumptions, Assumptions~\ref{as:A1} and \ref{as:A2}, indeed cover nonmonotone problems.

\beforesubsec
\subsection{Further Discussion of Related Work}\label{apdix:subsec:further_related_work}
\aftersubsec
As we already discussed in the introduction of the main text, both standard stochastic approximation and variance-reduction methods have been broadly studied for \eqref{eq:NE} and \eqref{eq:NI}, including \citep{juditsky2011solving,kotsalis2022simple,pethick2023solving}.
In this section, we further discuss some other related work to \eqref{eq:NE} and \eqref{eq:NI}, their special cases, and equivalent forms.

%%%%% (a)~Beyond monotonicity.
\noindent\textbf{$\mathrm{(a)}$~Beyond monotonicity.}
Classical  methods such as extragradient, prox-mirror, and projective schemes often relax the monotonicity to star-monotonicity, and other forms such as pseudo-monotonicity and quasi-monotonicity \citep{Konnov2001,noor2003extragradient,noor1999wiener,vuong2018weak}.
These assumptions are certainly weaker than the monotonicity and can cover some wider classes of problems, including some nonmonotone subclasses. 
Another extension of monotonicity is the weak-Minty solution condition in Assumption~\ref{as:A2}, which was proposed in early work, perhaps in the most recent one such as  \citep{diakonikolas2021efficient}, as an extension of the star-monotonicity and star-weak-monotonicity assumptions. 
Other following-up works include \citep{bohm2022solving,gorbunov2022extragradient,luo2022last}.
A comprehensive survey for extragradient-type methods using the weak-Minty solution condition can be found in \citep{tran2023sublinear,TranDinh2025a,tran2025accelerated}.
The monotonicity has also been extended to a weak monotonicity, or related, prox-regularity \citep{Rockafellar1997} (in particular, weak-convexity).
Other types of  hypo-monotonicity or co-monotonicity concepts can be found, e.g., in \citep{bauschke2020generalized}.
These concepts have been exploited to develop algorithms for solving  \eqref{eq:NE} and \eqref{eq:NI} and their special cases.
For stochastic methods, extensions beyond monotonicity have been also extensively explored.
For instance, some further structures beyond monotonicity such as weak solution were exploited for MVIs in \citep{song2020optimistic}, a pseudo-monotonicity was used in  \citep{boct2021minibatch,kannan2019optimal} for stochastic VIPs, a two-sided Polyak-{\L}ojasiewicz condition was extended to VIP in  \citep{yang2020global} to tackle a class on nonconvex-nononcave minimax problems, an expected co-coercivity was used \citep{loizou2021stochastic}, and a strongly star-monotone was further exploited in \citep{gorbunov2022stochastic}.
While these structures are occasionally used in different works, the relation between them is still largely elusive. 
In addition, their relation to concrete applications is still not well studied.

%%%%% (b)~Further discussion on stochastic methods.
\noindent\textbf{$\mathrm{(b)}$~Further discussion on stochastic methods.}
Under the monotonicity, several authors have exploited  the stochastic approximation approach \citep{RM1951} to develop stochastic variants for solving \eqref{eq:NE} and \eqref{eq:NI} and their special cases.
For example, a stochastic Mirror-Prox was proposed in \citep{juditsky2011solving}, which has convergence on a gap function, but requires a bounded domain assumption.
This approach was later extended to the extragradient method under additional assumptions in \citep{mishchenko2020revisiting}. 
In  \citep{hsieh2019convergence}, the authors discussed several methods for solving MVIs, a special case of \eqref{eq:NI}, including stochastic methods.
They experimented on numerical examples and showed that the norm of the operator can asymptotically converge for unconstrained MVIs  with a double learning rate.
In the last few years, there were many works focusing on developing stochastic methods for solving \eqref{eq:NE} and \eqref{eq:NI}, and their special cases using different techniques such as  single-call stochastic schemes in \citep{hsieh2019convergence}, non-accelerated and accelerated variance reduction with  Halpern-type iterations in \citep{cai2023variance,cai2022stochastic}, co-coercive structures in \citep{beznosikov2023stochastic}, and bilinear game models in \citep{li2022convergence}.
%%

%%%%%% (c)~The challenge of using biased estimators in algorithms for solving \eqref{eq:NI} and related problems
\noindent\textbf{$\mathrm{(c)}$~The challenge of using biased estimators in algorithms for solving \eqref{eq:NI} and related problems.}
In optimization, especially in nonconvex optimization, stochastic methods using biased estimators such as SARAH, Nested SVRG, Hybrid-SGD, and STORM can achieve better, even ``optimal'' oracle complexity compared to unbiased ones such as standard SVRG and SAGA, see, e.g., \citep{Cutkosky2019,driggs2019bias,Pham2019,Tran-Dinh2019a}.  
However, it becomes challenging in root-finding algorithms for solving \eqref{eq:NI} and related problems such as minimax optimization and VIPs.
One main reason for this is that the convergence analysis of these optimization methods relies on the objective function as a key metric to prove convergence guarantee and to estimate oracle complexity. 
However, in \eqref{eq:NI} and \eqref{eq:NE}, we do not have such an object, making it difficult to process the biased terms, including product terms  such as $\iprods{e^k, x^{k+1} - x^{\star}}$ and $\iprods{e^k, x^{k+1} - x^k}$, where $e^k$ is a bias rendered from the underlying stochastic estimator.
This is currently one of the main obstacles to move from using unbiased to biased estimators in root-finding algorithms, including our methods in this paper.

%%%%%% (d)~Comparison to  \citep{cai2023variance}.
\noindent\textbf{$\mathrm{(d)}$~Comparison to  \citep{cai2023variance}.}
Among many existing works, perhaps,  \citep{cai2023variance} is one of the most recent works that develops variance-reduction methods for solving \eqref{eq:NI} and achieves the state-of-the-art oracle complexity.
However,  \citep{cai2023variance} explores a different approach than ours, which relies on some recent development of the Halpern fixed-point iteration and a biased SARAH estimator.
Let us clarify the differences of this work and our paper here.
Algorithm 1 in \citep{cai2023variance} is a single-loop and achieves a better oracle complexity.
However, it requires a much stronger assumption, Assumption 3, which is a co-coercive condition. 
Note that this assumption excludes the well-known bilinear matrix game, or the synthetic WGAN model \eqref{eq:syn_WGAN_sup} below.
Section 4 of \citep{cai2023variance} studies both the monotone and the co-hypomonotone cases of \eqref{eq:NI}. 
The main idea is to reformulate \eqref{eq:NI} into a resolvent equation $J_{\eta(G+T)}x = 0$ and then apply a deterministic variant of Algorithm 1 to this equation, where $J_{\eta(G+T)}$ is co-coercive.
However, exactly evaluating $J_{\eta(G+T)}$ is impractical, one needs to approximate it by an appropriate algorithm.
For instance,  \citep{cai2023variance} suggests to use  the variance-reduced FRBS method in \citep{alacaoglu2022beyond} to approximate this resolvent, leading to a double loop algorithm. 
Note that this method also relies on a unbiased estimator, namely SVRG.
This approach is not a direct variance-reduced method (i.e., the inner loop can be any algorithm) as ours or Algorithm 1 of  \citep{cai2023variance}. 
Moreover, practically implementing as well as  rigorously analyzing an inexact double loop algorithm, when the inner loop is also a stochastic method, is often very challenging and technical as it is difficult to conduct a stopping criterion of the inner loop, and to select appropriate parameters. 
Nevertheless, our algorithms developed in this paper are simple to implement and applicable to both \eqref{eq:NE} and \eqref{eq:NI} whose weak-Minty solution exists.
These problems are broader than the ones in \citep{cai2023variance}.
We also believe that our oracle complexity in this paper can be further improved by exploiting enhancement techniques such as nested trick or multiple loops as done in \cite{cai2023variance,zhou2018stochastic}.


%In \citep{cai2023variance}, the authors exploited a loopless SARAH estimator to develop a Halpern-type algorithm for solving a co-coercive equation as a special case of \eqref{eq:NE}. This class of problems is much restrictive compared to ours, and it does not cover the well-known bilinear matrix game model.
%Then, the authors developed a double-loop strategy to weaken their assumption from the co-coerciveness to monotonicity and Lipschitz continuity, and then to co-hypomonotonicity and Lipschitz continuity. 
%Note that this is a different approach than ours in this paper, and their algorithms are complicated to implement. 
%In contrast, our methods use constant learning rates and solve a weaker class of problems compared to \citep{cai2023variance}. 
%In addition, our algorithms are very simple compared to the double-loop ones as in \citep{cai2023variance}.

%%%%% (e)~Randomized coordinate and cyclic coordinate methods for \eqref{eq:NE} and \eqref{eq:NI}.
\noindent\textbf{$\mathrm{(e)}$~Randomized coordinate and cyclic coordinate methods for \eqref{eq:NE} and \eqref{eq:NI}.}
Together with stochastic algorithms for solving \eqref{eq:NE} and \eqref{eq:NI} and their special cases, randomized coordinate methods have also been proposed to solve these problems, including \citep{combettes2018asynchronous,combettes2015stochastic,peng2016arock}.
Recent works on randomized coordinate and cyclic coordinate methods can be found, e.g., in \citep{chakrabarti2024block,cui2021analysis,hamedani2018iteration,song2023cyclic,tran2022accelerated,yousefian2018stochastic}.
These methods are not directly related to our work, but they can be considered as a dual form of stochastic methods in certain settings such as convex-concave minimax problems.
Studying relations between randomized coordinate methods and stochastic algorithms for \eqref{eq:NE} and \eqref{eq:NI} appears to be an interesting research topic.


\beforesubsec
\subsection{A Nonmonotone Example \eqref{eq:NI}}\label{apdix:subsec:nonmonotone_examples}
\aftersubsec
As an example of \eqref{eq:NI}, we can consider the following linear operators 
\vspace{-0.5ex}
\begin{equation*}
Gx := \mbf{G}x + \mbf{g} \quad \textrm{and} \quad Tx := \mbf{T}x, 
\vspace{-0.5ex}
\end{equation*}
where $\mbf{G}$ and $\mbf{T}$ are given square matrices in $\R^{p\times p}$ and $\mbf{g} \in\R^p$ is a given vector.
Clearly, for any $\mbf{G}$, $G$ is $L$-Lipschitz continuous with $L := \norms{\mbf{G}}$ (the operator norm of $\mbf{G}$).
Our goal is to choose $\mbf{G}$ and $\mbf{T}$ such that $\Psi{x} := Gx + Tx$ is nonmonotone and satisfies Assumption~\ref{as:A2}.
\begin{compactitem}
\item Clearly, we can choose $\mbf{G}$ and $\mbf{T}$ such that $\frac{1}{2}(\mbf{G} + \mbf{G}^{\top})$ is positive semidefinite and $\frac{1}{2}(\mbf{G}+\mbf{G}^{\top} +\mbf{T} + \mbf{T}^{\top})$ is not positive semidefinite.
This shows that $\Psi$ is nonmonotone.
\item
Now, Assumption~\ref{as:A2} holds if $\Psi$ is $\kappa$-co-hypomonotone, i.e. there exists $\kappa > 0$ such that $\iprods{u - v, x - y} \geq -\kappa\norms{u - v}^2$, for any $(x, u), (y, v) \in \gra{\Psi}$.
In the linear case, this condition is equivalent to
\vspace{-0.5ex}
\begin{equation*}
\mbf{S} := \frac{1}{2}(\mbf{G} + \mbf{G}^{\top} + \mbf{T} + \mbf{T}^{\top}) + \kappa(\mbf{G} + \mbf{T})^{\top}(\mbf{G} + \mbf{T}) \quad \textrm{is positive semidefinite.}
\vspace{-0.5ex}
\end{equation*}
However, since $\frac{1}{2}(\mbf{G} + \mbf{G}^{\top})$ is positive semidefinite, this condition holds if $\frac{1}{2}(\mbf{T} + \mbf{T}^{\top}) + \kappa(\mbf{G} + \mbf{T})^{\top}(\mbf{G} + \mbf{T})$ is positive semidefinite. 
In particular, $\Psi$ satisfies Assumption~\ref{as:A2}.
\end{compactitem}
For example, given any $\epsilon > 0$, we choose 
\vspace{-0.5ex}
\begin{equation*}
\mbf{G} := \begin{bmatrix}0 & 1 \\ 1 & 0\end{bmatrix}, \quad \mbf{T} := \begin{bmatrix}-\epsilon & 0 \\ 0 & 0\end{bmatrix}, \quad \textrm{and} \quad \kappa := \epsilon > 0.
\vspace{-0.5ex}
\end{equation*}
Then, it is clear that $\mbf{G}$ is symmetric and positive semidefinite, but $\mbf{G}+\mbf{T}$ is symmetric and not positive semidefinite.
Thus, $\Psi{x} = Gx + Tx$ is nonmonotone.
Moreover, $Gx = \mbf{G}x + \mbf{g}$ is $L$-Lipschitz continuous with $L = 1$.

Next, we check Assumption~\ref{as:A2}.
Clearly, $\mbf{G}$ is symmetric and positive semidefinite.
Moreover, since $\mbf{G}$ and $\mbf{T}$ are symmetric, we have 
\begin{equation*}
\mbf{M} := \frac{1}{2}(\mbf{T} + \mbf{T}^{\top}) + \kappa(\mbf{G} + \mbf{T})^{\top}(\mbf{G} + \mbf{T}) = 
\begin{bmatrix}-\epsilon & 0 \\ 0 & 0\end{bmatrix} + \kappa\begin{bmatrix}-\epsilon & 1 \\ 1 & 0\end{bmatrix}^{\top}\begin{bmatrix}-\epsilon & 1 \\ 1 & 0\end{bmatrix} 
=  \begin{bmatrix}-\epsilon + \kappa(1 + \epsilon^2) & -\kappa\epsilon \\ -\kappa \epsilon & \kappa\end{bmatrix}.
\end{equation*}
If we choose $\kappa = \epsilon$, then one can easily check that $\mbf{M}$ is positive semidefinite.
Hence, we can conclude that $\Psi$ is $\kappa$-co-hypomonotone with $\kappa = \epsilon > 0$.
In particular, $\Psi$ satisfies Assumption~\ref{as:A2}.
In addition, we can choose $\epsilon$ sufficiently small such that $L\kappa$ is sufficiently small, which fulfills the condition $L\kappa \leq \delta$ in Theorem~\ref{th:VROG4NI_convergence}.

 
%%%%%%%%%%%%%%%%%%%%%%%%%%%%%%%%%%%%%%%%%%%%%%%%%%%%%%%
%%%% A. The Proof of Technical Results in Section 2.
%%%%%%%%%%%%%%%%%%%%%%%%%%%%%%%%%%%%%%%%%%%%%%%%%%%%%%%
\beforesec
\section{The Proof of Technical Results in Section~\ref{sec:VRS_Estimator}}
\aftersec
This supplementary section provides the full proof of Lemma~\ref{le:SVGR_estimator} and Lemma~\ref{le:SAGA_estimator}.

\textbf{Further discussion of the FR quantity.}
Let us recall our quantity $S^k_{\gamma}$ defined by \eqref{eq:Sk_op} as follows: 
\begin{equation}\label{eq:Sk_op_sup}
S_{\gamma}^k := Gx^k  - \gamma Gx^{k-1}.
\tag{FRO}
\end{equation}
As we mentioned earlier, $\gamma$ plays a crucial role in our methods  as $\gamma \in \big(\frac{1}{2}, 1\big)$.
If $\gamma = \frac{1}{2}$, then we can write $S_{1/2}^k = \frac{1}{2}Gx^k + \frac{1}{2}(Gx^k - Gx^{k-1}) = \frac{1}{2}[ 2Gx^k - Gx^{k-1}]$ used in both the forward-reflected-backward splitting (FRBS) method \citep{malitsky2020forward} and the optimistic gradient method  \citep{daskalakis2018training}.

Note that if we write $Gx^k - Gx^{k-1} =  \hat{J}_G(x^k)(x^k - x^{k-1})$ by the Mean-Value Theorem, where $\hat{J}_G(x^k) := \int_0^1\nabla{G}(x^{k-1} + \tau(x^k - x^{k-1}))d\tau$, then $S^k_{\gamma} = (1-\gamma)G(x^k) + \gamma \hat{J}_G(x^k)(x^k - x^{k-1})$.
Clearly, if $\gamma$ is small, then $S_{\gamma}^k$ can be considered as an approximation of $Gx^k$ augmented by a second-order correction term $\gamma\hat{J}_G(x^k)(x^k - x^{k-1})$ (called Hessian-driven damping term or second-order dissipative term) widely used in dynamical systems for convex optimization, see, e.g., \citep{adly2021first,attouch2020convergence}.
These two viewpoints motivate the use of our new operator $S^k_{\gamma}$, not only in our \eqref{eq:VROG4NE} and \eqref{eq:VROG4NI}, but in other methods such as accelerated algorithms.
Thus, the results in Section~\ref{sec:VRS_Estimator} are of independent interest, and can potentially be used to develop other methods. 

\textbf{Other possible stochastic estimators for $S^k_{\gamma}$.}
One natural idea to construct an unbiased estimator for $S^k$ is to use an increasing mini-batch stochastic estimator as $\widetilde{S}_{\gamma}^k := \frac{1}{b_k}\sum_{i\in\Bc_k}[ G_ix^k - \gamma G_ix^{k-1}]$, where $\Bc_k$ is an increasing mini-batch in $[n]$, with $b_k := \vert \Bc_k\vert \geq \frac{b_{k-1}}{1-\rho_k} \geq b_{k-1}$, see, e.g., \citep{iusem2017extragradient}.
While this idea may work well for the general expectation case $Gx = \Expsk{\xi}{\mathbf{G}(x, \xi)}$, it may not be an ideal choice for the finite-sum operator \eqref{eq:finite_sum_form} as $b_k \leq n$, which requires to stop increasing after finite iterations (i.e. $\BigO{\frac{\ln(n)}{-\ln(1-\rho)}}$ iterations).
Other stochastic approximations may also fall into our class in Definition~\ref{de:ub_SG_estimator} such as JacSketch \citep{gower2021stochastic}, SEGA \citep{hanzely2018sega}, and quantized and compressed estimators ( see, e.g., \citep{horvath2023stochastic}).

%%%%%%%%%%%%%%%%%%%%%%%%%%%%%%%%
%%% A.1. Proof of Lemma 1.
%%%%%%%%%%%%%%%%%%%%%%%%%%%%%%%%
\beforesubsec
\subsection{Proof of Lemma~\ref{le:SVGR_estimator}: Loopless-SVRG Estimator}\label{subsec:le:SVGR_estimator}
\aftersubsec
Let us further expand Lemma~\ref{le:SVGR_estimator} in detail as follows and then provide its full proof.

%%% Lemma B.1.
\begin{lemma}\label{le:SVGR_estimator_full}
Let $S_{\gamma}^k := Gx^k - \gamma Gx^{k-1}$ be defined by \eqref{eq:Sk_op} and $\widetilde{S}_{\gamma}^k$ be generated by \eqref{eq:SVRG_estimator}.
We define
\vspace{-0.5ex}
\begin{equation}\label{eq:Deltak_quantity_full}
\begin{array}{lcl}
\Delta_k &:= &  \frac{1}{nb} \sum_{i=1}^n \Expk{ \norms{ G_ix^k - \gamma G_ix^{k-1} - (1 - \gamma) G_iw^k }^2 }.
\end{array}
\vspace{-0.5ex}
\end{equation}
Then, we have 
\begin{equation}\label{eq:SVRG_variance}
\begin{array}{lcl}
\Expsk{k}{\widetilde{S}_{\gamma}^k} & = & S_{\gamma}^k \equiv Gx^k - \gamma Gx^{k-1}, \\
\Expk{ \norms{\widetilde{S}_{\gamma}^k - S_{\gamma}^k }^2 } & \leq &   \Delta_k - \frac{1}{b} \Expk{ \norms{Gx^k - \gamma Gx^{k-1} - (1 - \gamma)Gw^k }^2 } \leq \Delta_k, \\
\Delta_k & \leq & \big( 1 - \frac{\mbf{p}}{2} \big) \Delta_{k-1}   +  \frac{(4 - 6\mbf{p} + 3\mbf{p}^2 )}{ nb \mbf{p} } \sum_{i=1}^n \Expk{ \norms{G_ix^k - G_ix^{k-1}}^2 } \\
&& + {~}  \frac{2 \gamma^2 (2 - 3\mbf{p} + \mbf{p}^2 ) }{nb \mbf{p}  } \sum_{i=1}^n \Expk{ \norms{G_ix^{k-1} - G_ix^{k-2} }^2 }.
\end{array}
\end{equation}
Consequently, the SVRG estimator $\widetilde{S}_{\gamma}^k$ constructed by \eqref{eq:SVRG_estimator}  satisfies Definition~\ref{de:ub_SG_estimator} with $\Delta_k$ in \eqref{eq:Deltak_quantity_full}, $\rho := \frac{\mbf{p} }{2} \in (0, 1]$, $C := \frac{ 4 - 6\mbf{p} + 3\mbf{p}^2 }{b\mbf{p}}$, and $\hat{C} := \frac{4\gamma^2(2 - 3\mbf{p} + \mbf{p}^2 )}{b\mbf{p}}$.
\end{lemma}

%%% Proof of Lemma B.1.
\begin{proof}
It is well-known, see, e.g., \citep{SVRG}, that $\widetilde{S}_{\gamma}^k$ is an unbiased estimator of $S^k$ conditioned on $\Fc_k$, we have $\Expsk{k}{\widetilde{S}_{\gamma}^k}  =  S_{\gamma}^k$.

Next, let $X_i := G_ix^k - \gamma G_ix^{k-1} - (1-\gamma)G_iw^k$ for any $i \in [n]$.
Then, we have $\Expsk{k}{X_i} = Gx^k - \gamma Gx^{k-1} -  (1-\gamma)Gw^k$ for any $i \in [n]$.
Since $\Bc_k$ is in $\Fc_k$, using the property of expectation, we can derive
\begin{equation*} 
\begin{array}{lcl}
\Expsk{k}{ \norms{  \widetilde{S}_{\gamma}^k - S_{\gamma}^k }^2 } & \overset{\tiny\eqref{eq:SVRG_estimator}}{=} & \Expsk{k}{\norms{ \frac{1}{b}\sum_{i\in\Bc_k}X_i  -  [ Gx^k + \gamma Gx^{k-1} -  (1-\gamma)Gw^k ] }^2 } \\
& = & \Expsk{k}{ \big\Vert \frac{1}{b}\sum_{i\in\Bc_k}\big[ X_i  -  \Expsk{k}{X_i} \big] \big\Vert^2 } \\
&\overset{\myeqc{1}}{ = } & \frac{1}{b^2} \Expsk{k}{ \sum_{i\in\Bc_k}\norms{X_i -  \Expsk{k}{X_i} }^2 }   \\
&\overset{\myeqc{2}}{ = } & \frac{1}{b^2} \Expsk{k}{ \sum_{i\in\Bc_k}\norms{  G_ix^k - \gamma G_ix^{k-1} - (1-\gamma)G_iw^k }^2 } - \frac{1}{b} \big[ \Expsk{k}{X_i} \big]^2  \\
& = & \frac{1}{n b}  \sum_{i = 1}^n \norms{  G_ix^k - \gamma G_ix^{k-1} - (1-\gamma)G_iw^k }^2 - \frac{1}{b} \big[ \Expsk{k}{X_i} \big]^2.
\end{array}
\end{equation*}
Here, $\myeqc{1}$ holds due to the i.i.d. property of $\Bc_k$, and  $\myeqc{2}$ holds since $\Expsk{k}{\norms{X_i - \Expsk{k}{X_i}}^2} = \Expsk{k}{\norms{X_i}^2} - \big( \Expsk{k}{ X_i} \big)^2$.
This estimate implies the second line of \eqref{eq:SVRG_variance} by taking the total expectation $\Expk{\cdot}$ both sides and the definition of $\Delta_k$ from \eqref{eq:Deltak_quantity_full}.

Now, from \eqref{eq:SG_estimator_w} and \eqref{eq:Deltak_quantity_full}, we can show that
\begin{equation*}
\begin{array}{lcl}
\Delta_k & \overset{\eqref{eq:Deltak_quantity_full}}{:=} &  \frac{1}{n b}  \sum_{i = 1}^n \Expk{  \norms{  G_ix^k - \gamma G_ix^{k-1} - (1-\gamma)G_iw^k }^2 } \\
& \overset{\eqref{eq:SG_estimator_w}}{ = } &  \frac{(1 - \mbf{p})}{n b}  \sum_{i = 1}^n \Expk{  \norms{  G_ix^k - \gamma G_ix^{k-1} - (1 - \gamma)G_iw^{k-1} }^2 } \\
&& +  {~}  \frac{ \mbf{p}}{n b}  \sum_{i = 1}^n \Expk{  \norms{  G_ix^k - \gamma G_ix^{k-1} - (1-\gamma)G_ix^{k-1} }^2 } \\
&\overset{\myeqc{1}}{ \leq } &  \frac{(1+c)(1- \mbf{p} )}{nb}  \sum_{i = 1}^n \Expk{ \norms{G_ix^{k-1} - \gamma G_ix^{k-2}   - (1-\gamma)G_iw^{k-1}}^2 }\\
&& + {~} \frac{(1 + c)(1 - \mbf{p} )}{c n b}  \sum_{i = 1}^n \Expk{ \norms{   G_ix^k - \gamma G_ix^{k-1}  -  [G_ix^{k-1} - \gamma G_ix^{k-2}] }^2 } \\
&& + {~} \frac{ \mbf{p} }{ n b }  \sum_{i = 1}^n \Expk { \norms{ G_ix^k - G_ix^{k-1} }^2 } \\
&\overset{\myeqc{2}}{ \leq } &  \frac{ (1+c)(1- \mbf{p} )}{n b}   \sum_{i = 1}^n \Expk{ \norms{G_ix^{k-1} - \gamma G_ix^{k-2}   - (1-\gamma)G_iw^{k-1}}^2 }\\
&& + {~} \frac{2(1 + c)(1 - \mbf{p} ) \gamma^2  }{n b c  }  \sum_{i = 1}^n \Expk{ \norms{   G_ix^{k-1} - G_ix^{k-2} }^2 } \\
&& + {~} \frac{1}{n b}\big[ \mbf{p} + \frac{2(1 + c)(1 - \mbf{p} )}{c} \big]   \sum_{i = 1}^n \Expk{ \norms{ G_ix^k - G_ix^{k-1} }^2 } \\
& =  &  (1+c)(1- \mbf{p} )  \Delta_{k-1}  + \frac{2(1 + c)(1 - \mbf{p} ) \gamma^2  }{n b c }  \sum_{i = 1}^n \Expk{ \norms{   G_ix^{k-1} - G_ix^{k-2} }^2 } \\
&& + {~} \frac{1}{n b}\big[ \mbf{p} + \frac{2(1 + c)(1 - \mbf{p} )}{c} \big]   \sum_{i = 1}^n \Expk{ \norms{ G_ix^k - G_ix^{k-1} }^2 }.
\end{array}
\end{equation*}
Here, in both inequalities $\myeqc{1}$ and $\myeqc{2}$, we have used Young's inequality twice.
If we choose $c := \frac{ \mbf{p} }{2(1 - \mbf{p} )}$, then $(1+c)(1 - \mbf{p} ) = 1 - \frac{\mbf{p}}{2}$, $\frac{(1+c)(1 - \mbf{p})}{c} = (1 - \mbf{p})\big(1 + \frac{2(1 - \mbf{p} )}{\mbf{p} }\big) = \frac{(2- \mbf{p} )(1- \mbf{p} )}{\mbf{p} } = \frac{2 - 3 \mbf{p} + \mbf{p}^2}{ \mbf{p} }$, % \leq \frac{2}{ \mbf{p} }$, 
and $\frac{2(1+c)(1 - \mbf{p} )}{c} + \mbf{p}  = \frac{4 - 6 \mbf{p} + 3 \mbf{p}^2}{ \mbf{p} }$. % \leq \frac{4}{ \mbf{p} }$.
Hence, we obtain 
\begin{equation*}
\begin{array}{lcl}
\Delta_k  & \leq &  \big(1 -  \frac{\mbf{p}}{2} \big)  \Delta_{k-1}  +   \frac{(4 - 6\mbf{p} + 3\mbf{p}^2 )}{ nb  \mbf{p} } \sum_{i=1}^n \Expk{ \norms{G_ix^k - G_ix^{k-1}}^2 } \\
&& + {~}  \frac{2 \gamma^2(2 - 3\mbf{p} + \mbf{p}^2 ) }{ nb \mbf{p}  } \sum_{i=1}^n \Expk{ \norms{G_ix^{k-1} - G_ix^{k-2} }^2 }.
\end{array}
\end{equation*}
This is exactly the last inequality of  \eqref{eq:SVRG_variance}.
\end{proof}
%%% End of Proof.

%%%%%%%%%%%%%%%%%%%%%%%%%%%%%%
%%% A.2. Proof of Lemma 2.2.
%%%%%%%%%%%%%%%%%%%%%%%%%%%%%%
\beforesubsec
\subsection{Proof of Lemma~\ref{le:SAGA_estimator}: SAGA estimator}\label{subsec:le:SAGA_estimator}
\aftersubsec
Similarly, we also further expand Lemma~\ref{le:SAGA_estimator} in detail as follows and then provide its full proof.

%%% Lemma 2.2.
\begin{lemma}\label{le:SAGA_estimator_full}
Let $S^k_{\gamma} := Gx^k - \gamma Gx^{k-1}$ be defined by \eqref{eq:Sk_op} and $\widetilde{S}_{\gamma}^k$ be generated by the SAGA estimator \eqref{eq:SAGA_estimator}, and $e^k := \widetilde{S}^k_{\gamma} - S^k_{\gamma}$.
We consider the following quantity:
\begin{equation}\label{eq:Deltak_hat_quantity_full}
\begin{array}{lcl}
 \Delta_k &:= &  \frac{1}{nb} \sum_{i=1}^n \Expk{ \norms{ G_ix^k - \gamma G_ix^{k-1} - (1-\gamma) \hat{G}_{i}^k }^2 }.
\end{array}
\end{equation}
Then, we have 
\begin{equation}\label{eq:SAGA_variance}
\hspace{-4ex}
\begin{array}{lcl}
\Expsk{k}{\widetilde{S}_{\gamma}^k} & = & S_{\gamma}^k \equiv Gx^k - \gamma Gx^{k-1}, \\
\Expk{ \norms{\widetilde{S}_{\gamma}^k - S_{\gamma}^k }^2 }  & \leq &   \Delta_k - \frac{1}{b} \Expk{ \big\Vert Gx^k - \gamma Gx^{k-1} -  \frac{(1-\gamma)}{n}\sum_{i =1}^n\hat{G}^k_i \big\Vert^2 } \leq \Delta_k, \\
\Delta_k & \leq &  \big(1 -  \frac{b}{2n} \big)  \Delta_{k-1}   +   \frac{[ 2(n - b)(2n+b) + b^2] }{ n^2b^2 } \sum_{i=1}^n \Expk{ \norms{G_ix^k - G_ix^{k-1}}^2 } \\
&& + {~}  \frac{2(n - b)(2n+b) \gamma^2 }{ n^2b^2 } \sum_{i=1}^n \Expk{ \norms{G_ix^{k-1} - G_ix^{k-2} }^2 }.
\end{array}
\hspace{-4ex}
\end{equation}
Consequently, the SAGA estimator $\widetilde{S}_{\gamma}^k$ constructed by \eqref{eq:SAGA_estimator} satisfies Definition~\ref{de:ub_SG_estimator} with $\Delta_k$ in \eqref{eq:Deltak_hat_quantity_full}, $\rho := \frac{b}{2n} \in (0, 1]$, $C :=  \frac{[ 2(n - b)(2n+b) + b^2] }{ n b^2}$, and $\hat{C} := \frac{2(n - b)(2n+b) \gamma^2}{ nb^2}$.
\end{lemma}

%%% Proof of Lemma B.2.
\begin{proof}
It is well-known, see, e.g., \citep{Defazio2014}, that $\widetilde{S}^k_{\gamma}$ defined by \eqref{eq:SAGA_estimator} is an unbiased estimator of $S^k$.
Indeed, we have $\Expsk{k}{\hat{G}_{\Bc_k}^k} = \frac{1}{n}\sum_{i=1}^n\hat{G}_i^k$, $\Expsk{k}{G_{\Bc_k}x^k} = Gx^k$, and $\Expsk{k}{G_{\Bc_k}x^{k-1}} = Gx^{k-1}$.
Using these relations and the definition of $\widetilde{S}^k$, we have
\begin{equation*}
\begin{array}{lcl}
\Expsk{k}{\widetilde{S}^k} & = &  \Expsk{k}{\frac{(1-\gamma)}{n}\sum_{i=1}^n\hat{G}_i^k} - (1-\gamma)\Expsk{k}{\hat{G}_{\Bc_k}^k} + \Expsk{k}{G_{\Bc_k}x^k} - \gamma \Expsk{k}{G_{\Bc_k}x^{k-1}} \\
& = & \frac{(1-\gamma)}{n}\sum_{i=1}^n\hat{G}_i^k - \frac{(1-\gamma)}{n}\sum_{i=1}^n\hat{G}_i^k + Gx^k - \gamma Gx^{k-1} \\
& = &  Gx^k - \gamma Gx^{k-1} \\
& = & S^k.
\end{array}
\end{equation*}
Hence, $\widetilde{S}^k$ is an unbiased estimator of $S^k$.

Next, let $X_i := G_ix^k - \gamma G_ix^{k-1} - (1-\gamma)\hat{G}_i^k$ for any $i \in [n]$.
Then, we have $\Expsk{k}{X_i} = Gx^k - \gamma Gx^{k-1} -  \frac{(1-\gamma)}{n}\sum_{i=1}^n\hat{G}^k_i$ for any $i \in [n]$.
Therefore, we can derive
\begin{equation*} 
\begin{array}{lcl}
\Expsk{k}{ \norms{  \widetilde{S}_{\gamma}^k - S_{\gamma}^k }^2 } & = & \Expsk{k}{\norms{ \frac{1}{b}\sum_{i\in\Bc_k}X_i  -  \big[ Gx^k + \gamma Gx^{k-1}  -  \frac{(1-\gamma)}{n}\sum_{i=1}^n \hat{G}_i^k \big] }^2 } \\
& = & \Expsk{k}{\norms{ \frac{1}{b}\sum_{i\in\Bc_k}X_i  -  \Expsk{k}{X_i} ] }^2 } \\
& = & \frac{1}{b^2} \Expsk{k}{ \sum_{i\in\Bc_k}\norms{X_i - \Expsk{k}{X_i} }^2 } \\
& = & \frac{1}{b^2} \Expsk{k}{ \sum_{i\in\Bc_k}\norms{  G_ix^k - \gamma G_ix^{k-1} - (1-\gamma)\hat{G}_i^k }^2 } - \frac{1}{b} \big[ \Expsk{k}{X_i} \big]^2 \\
& = & \frac{1}{n b}  \sum_{i = 1}^n \norms{  G_ix^k - \gamma G_ix^{k-1} - (1-\gamma)\hat{G}_i^k }^2 - \frac{1}{b} \big[ \Expsk{k}{X_i} \big]^2.
\end{array}
\end{equation*}
This implies the second line of \eqref{eq:SAGA_variance} by taking the total expectation $\Expk{\cdot}$ both sides.

Now, from \eqref{eq:SAGA_ref_points} and \eqref{eq:Deltak_hat_quantity_full} and the rule \eqref{eq:SAGA_ref_points}, for any $c > 0$, by Young's inequality, we can show that
\begin{equation*}
\begin{array}{lcl}
\Delta_k & \overset{\eqref{eq:Deltak_hat_quantity_full}}{ := } &  \frac{1}{n b}  \sum_{i = 1}^n \Expk{  \norms{  G_ix^k - \gamma G_ix^{k-1} - (1-\gamma)\hat{G}_i^k }^2 } \\
& \overset{ \eqref{eq:SAGA_ref_points} }{ = } & \big(1 - \frac{b}{n}\big)  \frac{1}{n b}  \sum_{i = 1}^n \Expk{  \norms{  G_ix^k - \gamma G_ix^{k-1} - (1-\gamma)\hat{G}_i^{k-1} }^2 } \\
&& + {~}  \frac{b}{n} \cdot  \frac{1}{n b}  \sum_{i = 1}^n \Expk{  \norms{  G_ix^k - \gamma G_ix^{k-1} - (1-\gamma) G_ix^{k-1} }^2 } \\
& \leq &   \frac{(1+c)}{ n b} \big(1 - \frac{b}{n}\big)  \sum_{i = 1}^n \Expk{  \norms{  G_ix^{k-1} - \gamma G_ix^{k-2} - (1-\gamma) \hat{G}_i^{k-1} }^2 } \\
&& + {~} \frac{(1 + c)}{c n b}\big(1 - \frac{b}{n}\big)  \sum_{i = 1}^n \Expk{  \norms{  G_ix^k - \gamma G_ix^{k-1} - ( G_ix^{k-1} -  \gamma G_ix^{k-2}) }^2 } \\
&& + {~}  \frac{1}{n^2}   \sum_{i = 1}^n \Expk{  \norms{  G_ix^k -  G_ix^{k-1} }^2 } \\
%%
&  \leq  &  (1 + c) \big(1 - \frac{b}{n}\big) \Delta_{k-1} + \big[ \frac{1}{n^2} + \big(1 - \frac{b}{n}\big)\frac{2(1+c)}{c n b} \big]  \sum_{i = 1}^n \Expk{  \norms{  G_ix^k - G_ix^{k-1} }^2  } \\
&& +  {~} \frac{2(1+c)\gamma^2}{c n b} \big(1 - \frac{b}{n}\big)  \sum_{i = 1}^n \Expk{ \norms{  G_ix^{k-1} - G_ix^{k-2}  }^2 }.
\end{array}
\end{equation*}
If we choose $c := \frac{b}{2n} \in (0, 1)$, then $(1 - \frac{b}{n})(1+c) = 1 - \frac{b}{2n} - \frac{b^2}{2n^2} \leq 1 - \frac{b}{2n}$.
Hence,  we can further upper bound the last inequality as 
\begin{equation*}
\begin{array}{lcl}
\Delta_k & \leq &  \big(1 -  \frac{b}{2n} \big)  \Delta_{k-1}  +   \frac{[ 2(n - b)(2n + b) +  b^2] }{ n^2 b^2 } \sum_{i=1}^n \Expk{ \norms{G_ix^k - G_ix^{k-1}}^2 } \\
&& + {~}  \frac{2(n - b)(2n + b) \gamma^2 }{ n^2b^2  } \sum_{i=1}^n \Expk{ \norms{G_ix^{k-1} - G_ix^{k-2} }^2 }.
\end{array}
\end{equation*}
This is exactly the last inequality of  \eqref{eq:SAGA_variance}.
%\Eproof
\end{proof}
%%% End of Proof.

%%%%%%%%%%%%%%%%%%%%%%%%%%%%%%%%%%%%%%%%%%%%%%%%%%%%%%%
%%% B. Convergence Analysis of VROG for NE: Technical Proofs.
%%%%%%%%%%%%%%%%%%%%%%%%%%%%%%%%%%%%%%%%%%%%%%%%%%%%%%%
\section{Convergence Analysis of \ref{eq:VROG4NE} for \eqref{eq:NE}: Technical Proofs}\label{apdx:sec:VROG4NE}
To analyze our \eqref{eq:VROG4NE} scheme, we introduce the following two functions:
\begin{equation}\label{eq:VROG4NE_Lfuncs}
\begin{array}{lcl}
\Lc_k & := &  \norms{x^k + \gamma\eta Gx^{k-1} - x^{\star}}^2 + \mu \norms{x^k - x^{k-1}}^2, \\
\Ec_k &:= & \Lc_k + \frac{\eta^2(1+\mu)(1-\rho)}{\rho} \Delta_{k-1} +   \frac{L^2\eta^2\hat{C}(1+\mu)}{ \rho } \norms{x^{k-1} - x^{k-2}}^2,
\end{array}
\end{equation}
where $\mu$ is a given positive parameter, $\rho$, $C$, $\hat{C}$, and $\Delta_k$ are given in Definition~\ref{de:ub_SG_estimator},  and $x^{-2} = x^{-1} = x^0$.
Clearly, we have $\Lc_k \geq 0$ and $\Ec_k \geq 0$ for all $k \geq 0$ a.s.

One key step to analyze the convergence of \eqref{eq:VROG4NE} is to prove a descent property of $\Ec_k$ defined by \eqref{eq:VROG4NE_Lfuncs}.
The following lemma provides such a key estimate to prove the convergence of \eqref{eq:VROG4NE}.

%%% Lemma 3.1.
\begin{lemma}\label{le:VROG4NE_key_estimate1}
Suppose that Assumptions~\ref{as:A1} and \ref{as:A2} hold for \eqref{eq:NE}.
Let $\sets{x^k}$ be generated by \eqref{eq:VROG4NE} and $\Ec_k$ be defined by \eqref{eq:VROG4NE_Lfuncs} for any $\gamma \in [0, 1]$.
Then, with $M := \frac{\gamma(1 + \mu - \gamma)}{\mu} + \frac{ (1+\mu)(C + \hat{C})}{\mu \rho }$, we have
\begin{equation}\label{eq:VROG4NE_key_estimate1}
\begin{array}{lcl}
\Expk{ \Ec_k } - \Expk{ \Ec_{k+1} } & \geq &  \mu \left(  1 -  M \cdot L^2\eta^2  \right) \Expk{ \norms{x^k - x^{k-1}}^2 }  \\
&& + {~}  \eta (1-\gamma) \big[ \eta \big( 2\gamma - 1 - \mu \big) - 2\kappa \big]  \Expk{ \norms{Gx^k}^2 } \\
&& + {~} \eta^2 \gamma(1-\gamma)(1+\mu) \Expk{ \norms{Gx^{k-1}}^2 }.
\end{array}
\end{equation}
\end{lemma}

%%% Proof of Lemma 3.1.
\begin{proof}
First, using $x^{k+1} := x^k - \eta \widetilde{S}_{\gamma}^k$ from \eqref{eq:VROG4NE}, we can expand
\begin{equation*} 
\begin{array}{lcl}
\norms{x^{k+1} + \gamma\eta Gx^{k} - x^{\star}}^2  & \overset{\tiny\eqref{eq:VROG4NE}}{=} & \norms{x^k - x^{\star} + \gamma\eta Gx^k - \eta \widetilde{S}_{\gamma}^k }^2 \\
& = & \norms{x^k - x^{\star}}^2 + 2\gamma\eta\iprods{Gx^k, x^k - x^{\star}} +  \gamma^2\eta^2\norms{Gx^k}^2  \\
&& - {~}  2\eta \iprods{\widetilde{S}_{\gamma}^k, x^k - x^{\star}} - 2\gamma\eta^2 \iprods{Gx^k, \widetilde{S}_{\gamma}^k} + \eta^2\norms{\widetilde{S}_{\gamma}^k}^2.
\end{array}
\end{equation*}
Second, it is obvious to show that
\begin{equation*} 
\begin{array}{lcl}
\norms{x^k + \gamma\eta Gx^{k-1} - x^{\star}}^2 & = & \norms{x^k - x^{\star}}^2 + 2\gamma\eta\iprods{Gx^{k-1}, x^k - x^{\star}} + \gamma^2\eta^2\norms{Gx^{k-1}}^2.
\end{array}
\end{equation*}
Third, using again $x^{k+1} := x^k - \eta \widetilde{S}_{\gamma}^k$ from \eqref{eq:VROG4NE}, we can show that
\begin{equation*} 
\begin{array}{lcl}
\norms{x^{k+1} - x^k}^2 &= & \eta^2\norms{\widetilde{S}_{\gamma}^k}^2.
\end{array}
\end{equation*}
Combining three expressions above, and using $\Lc_k$ from \eqref{eq:VROG4NE_Lfuncs}, we can establish  that 
\begin{equation}\label{eq:VROG4NE_proof1} 
\begin{array}{lcl}
\Lc_k - \Lc_{k+1}  & = & \norms{x^k + \gamma\eta Gx^{k-1} - x^{\star}}^2 - \norms{x^{k+1} + \gamma\eta Gx^{k} - x^{\star}}^2 \\
&& + {~} \mu\norms{x^k - x^{k-1}}^2 - \mu \norms{x^{k+1} - x^k}^2\\
& = & 2\gamma\eta\iprods{Gx^{k-1}, x^k - x^{\star}} - 2\gamma\eta \iprods{ Gx^k, x^k - x^{\star}}   +  \gamma^2\eta^2\norms{Gx^{k-1}}^2 \\
&& - {~} \gamma^2\eta^2\norms{Gx^k}^2 +   2\eta \iprods{\widetilde{S}^k_{\gamma}, x^k - x^{\star}} + 2\gamma\eta^2 \iprods{Gx^k, \widetilde{S}_{\gamma}^k} \\
&& +  {~} \mu \norms{x^k - x^{k-1}}^2 - \eta^2(1 + \mu) \norms{\widetilde{S}_{\gamma}^k}^2.
\end{array}
\end{equation}
Next, since $\Expsk{k}{\widetilde{S}^k_{\gamma}} = S_{\gamma}^k \equiv Gx^k - \gamma Gx^{k-1}$ as shown in the first line of \eqref{eq:ub_SG_estimator} of Definition~\ref{de:ub_SG_estimator}.
Moreover, since $\widetilde{S}^k_{\gamma}$ is conditionally independent of $x^k-x^{\star}$ and $Gx^k$ w.r.t. the $\sigma$-field $\Fc_k$, we have 
\begin{equation*} 
\begin{array}{lcl}
\Expsk{k}{ \iprods{\widetilde{S}^k_{\gamma}, x^k - x^{\star}} } &= & \iprods{Gx^k, x^k - x^{\star}} - \gamma\iprods{Gx^{k-1}, x^k - x^{\star}}, \\
2\Expsk{k}{ \iprods{\widetilde{S}^k_{\gamma}, Gx^k} } & = &  2 \norms{Gx^k}^2 - 2\gamma \iprods{Gx^{k-1}, Gx^k} \\
& = & (2 - \gamma) \norms{Gx^k}^2 - \gamma\norms{Gx^{k-1}}^2 + \gamma\norms{Gx^k - Gx^{k-1}}^2.
\end{array}
\end{equation*}
Taking the conditional expectation $\Expsk{k}{\cdot}$ both sides of \eqref{eq:VROG4NE_proof1} and using the last two expressions, we can show that
\begin{equation*} 
\begin{array}{lcl}
\Lc_k - \Expsk{k}{ \Lc_{k+1} }
& = & 2\gamma\eta\iprods{Gx^{k-1}, x^k - x^{\star}}  - 2\gamma\eta \iprods{Gx^k, x^k - x^{\star} } +  \gamma^2\eta^2\norms{Gx^{k-1}}^2 \\
&& - {~} \gamma^2\eta^2\norms{Gx^k}^2 + 2\eta \Expsk{k}{ \iprods{\widetilde{S}^k_{\gamma}, x^k - x^{\star}} } + 2\gamma\eta^2 \Expsk{k}{ \iprods{Gx^k, \widetilde{S}^k_{\gamma} }}  \\
&& - {~}  \eta^2(1 + \mu)  \Expsk{k}{ \norms{\widetilde{S}^k _{\gamma} }^2 }  +   \mu \norms{x^k - x^{k-1}}^2 \\
%& = & 2\gamma \eta\iprods{Gx^{k-1}, x^k - x^{\star}}  - 2\gamma\eta \iprods{Gx^k, x^k - x^{\star} }   +  \gamma^2\eta^2\norms{Gx^{k-1}}^2 \\
%&&  - {~} \gamma^2\eta^2\norms{Gx^k}^2  +  2\eta \big[  \iprods{Gx^k, x^k - x^{\star}} - \gamma\iprods{Gx^{k-1}, x^k - x^{\star}} \big] \\
%&& + {~}  \gamma\eta^2 \big[ (2 - \gamma) \norms{Gx^k}^2 - \gamma\norms{Gx^{k-1}}^2 + \gamma\norms{Gx^k - Gx^{k-1}}^2 \big] \\
%&& - {~}  \eta^2(1 + \mu) \Expsk{k}{ \norms{\widetilde{S}^k_{\gamma} }^2 }  +   \mu \norms{x^k - x^{k-1}}^2 \\
%%%%%
%& = & 2(\beta - \eta\gamma)\iprods{Gx^{k-1}, x^k - x^{\star}} + 2( \eta - \beta ) \iprods{Gx^k, x^k - x^{\star}} \\
%&& + {~} \beta (\beta - \eta\gamma) \norms{Gx^{k-1}}^2 + \beta \big[  (2 - \gamma)\eta - \beta \big] \norms{Gx^k}^2 \\
%&& + {~} \beta\eta\gamma\norms{Gx^k - Gx^{k-1}}^2  -  \eta^2(1 + \mu)  \Expsk{k}{ \norms{\widetilde{S}^k}^2 }  +  \mu \norms{x^k - x^{k-1}}^2 \\
%%%%%
& = & 2\eta( 1 - \gamma ) \iprods{Gx^k, x^k - x^{\star}} + 2\gamma(1-\gamma)\eta^2 \norms{Gx^k}^2 \\
&& + {~} \gamma^2\eta^2 \norms{Gx^k - Gx^{k-1}}^2  -  \eta^2(1 + \mu)  \Expsk{k}{ \norms{\widetilde{S}^k_{\gamma}}^2 }  +  \mu \norms{x^k - x^{k-1}}^2.
\end{array}
\end{equation*}
Since $\widetilde{S}^k_{\gamma}$ is an unbiased estimator of $S^k_{\gamma}$, if we denote $e^k := \widetilde{S}^k_{\gamma} - S^k_{\gamma}$, then we have $\Expsk{k}{e^k} = 0$.
Hence, we can show that $\Expsk{k}{ \norms{\widetilde{S}^k_{\gamma} }^2 } = \Expsk{k}{ \norms{S^k_{\gamma} + e^k}^2 } = \norms{S^k_{\gamma} }^2 + 2 \Expsk{k}{\iprods{e^k, S^k_{\gamma} }} + \Expsk{k}{ \norms{e^k}^2 } = \Expsk{k}{ \norms{e^k}^2 } + \norms{S^k_{\gamma} }^2$.
Using this relation and $S^k_{\gamma} = Gx^k - \gamma Gx^{k-1}$, we can show that
\begin{equation*}
\begin{array}{lcl}
\Expsk{k}{ \norms{\widetilde{S}^k_{\gamma} }^2 } &= & \norms{S^k_{\gamma} }^2 + \Expsk{k}{ \norms{e^k}^2 }  = \norms{Gx^k  - \gamma Gx^{k-1}}^2 + \Expsk{k}{ \norms{e^k}^2 } \\
&= & \norms{Gx^k}^2 - 2\gamma \iprods{Gx^k, Gx^{k-1}} + \gamma^2\norms{Gx^{k-1}}^2 + \Expsk{k}{ \norms{e^k}^2 } \\
& = & (1-\gamma)\norms{Gx^k}^2 - \gamma(1-\gamma)\norms{Gx^{k-1}}^2 + \gamma\norms{Gx^k - Gx^{k-1}}^2 + \Expsk{k}{ \norms{e^k}^2 }.
\end{array}
\end{equation*}
Substituting this expression into the last estimate, we can show that
\begin{equation*} 
\begin{array}{lcl}
\Lc_k - \Expsk{k}{ \Lc_{k+1} }
& = &  2\eta(1-\gamma)  \iprods{Gx^k, x^k - x^{\star}} + \eta^2(1 - \gamma) \big(2\gamma - 1 -\mu \big)  \norms{Gx^k}^2 \\
&& + {~} \eta^2\gamma(1-\gamma)(1+\mu)  \norms{Gx^{k-1}}^2 -  \gamma \eta^2( 1 +  \mu - \gamma) \norms{Gx^k - Gx^{k-1}}^2  \\
&& - {~}  \eta^2(1 + \mu)  \Expsk{k}{ \norms{e^k}^2 }  +  \mu \norms{x^k - x^{k-1}}^2.
\end{array}
\end{equation*}
Taking the total expectation $\Expk{\cdot}$ both sides of this expression, we get
\begin{equation*} 
\begin{array}{lcl}
\Expk{ \Lc_k } - \Expk{ \Lc_{k+1} } & = &   2(1-\gamma)\eta \Expk{  \iprods{Gx^k, x^k - x^{\star}} } + \eta^2\gamma(1-\gamma)(1+\mu) \Expk{ \norms{Gx^{k-1}}^2  }  \\
&& + {~} \eta^2(1 - \gamma) \big(2\gamma - 1 -\mu \big) \Expk{ \norms{Gx^k}^2 }  +  \mu \Expk{ \norms{x^k - x^{k-1}}^2 } \\
&& - {~}   \gamma \eta^2(1 + \mu - \gamma) \Expk{ \norms{Gx^k - Gx^{k-1}}^2  } -  \eta^2(1 + \mu) \Expk{ \norms{e^k}^2 }  .
\end{array}
\end{equation*}
By Young's inequality in  $\myeqc{1}$ and \eqref{eq:L_smooth} of Assumption~\ref{as:A1}, we have
\begin{equation}\label{eq:L_smooth2}
\begin{array}{lcl}
\norms{Gx^k - Gx^{k-1}}^2 & = &  \Vert \frac{1}{n}\sum_{i=1}^n[G_ix^k - G_ix^{i-1}]\Vert^2 \overset{ \myeqc{1} }{ \leq }  \frac{1}{n}\sum_{i=1}^n\norms{G_ix^k - G_ix^{k-1}}^2 \\
& \overset{\tiny\eqref{eq:L_smooth}}{ \leq } & L^2\norms{x^k - x^{k-1}}^2.
\end{array}
\end{equation}
Utilizing this inequality, $ \iprods{Gx^k, x^k - x^{\star}} \geq -\kappa \norms{Gx^k}^2$ from Assumption~\ref{as:A2} with $T = 0$, and $\Expk{ \norms{e^k}^2 } \leq  \Delta_k$ from \eqref{eq:ub_SG_estimator}, we can bound the last expression as
\begin{equation}\label{eq:SFP_proof5}  
\hspace{-2ex}
\begin{array}{lcl}
\Expk{ \Lc_k } - \Expk{ \Lc_{k+1} } & \geq &  {~} \big[ \mu -   L^2\eta^2\gamma(1+\mu-\gamma) \big] \Expk{ \norms{x^k - x^{k-1}}^2 } \\
&& + {~} (1+\mu)\gamma(1-\gamma)\eta^2 \Expk{ \norms{Gx^{k-1}}^2 }  \\
&& + {~}   \eta (1-\gamma) \big[ \eta \big( 2\gamma - 1 - \mu \big) - 2\kappa \big] \Expk{ \norms{Gx^k}^2 } - \eta^2(1 + \mu) \Delta_k.
\end{array}
\hspace{-2ex}
\end{equation}
By the third line of  \eqref{eq:ub_SG_estimator} in Definition~\ref{de:ub_SG_estimator} and again \eqref{eq:L_smooth}, we have
\begin{equation*} 
\begin{array}{lcl}
\Delta_{k} & \leq &  (1 - \rho ) \Delta_{k-1}   + C L^2 \Expk{ \norms{x^{k} - x^{k-1}}^2 }  +  \hat{C} L^2  \Expk{ \norms{x^{k-1} - x^{k-2}}^2 }.
\end{array}
\end{equation*}
Rearranging this inequality, we get
\begin{equation*} 
\begin{array}{lcl}
\Delta_{k} & \leq &   \big( \frac{1-\rho}{\rho} \big) \big( \Delta_{k-1} - \Delta_{k} \big)  + \frac{\hat{C} L^2}{ \rho } \big[ \Expk{ \norms{x^{k-1} - x^{k-2}}^2  } -  \Expk{ \norms{x^k - x^{k-1}}^2 } \big] \\
&& + {~}  \frac{(C + \hat{C})L^2}{ \rho }  \Expk{ \norms{x^{k} - x^{k-1}}^2 }.
\end{array}
\end{equation*}
Substituting this inequality into \eqref{eq:SFP_proof5}, we can show that
\begin{equation*} 
\begin{array}{lcl}
\Expk{ \Lc_k } - \Expk{ \Lc_{k+1} } & \geq &  \Big[ \mu -   L^2\eta^2\gamma(1 + \mu - \gamma) - \frac{L^2\eta^2(1+\mu)(C + \hat{C})}{ \rho } \Big] \Expk{ \norms{x^k - x^{k-1}}^2 } \\
&& + {~}  \eta (1-\gamma) \big[ \eta \big( 2\gamma - 1 - \mu \big) - 2 \kappa \big]  \Expk{ \norms{Gx^k}^2 }  \\
&& + {~} (1+\mu)\gamma(1-\gamma)\eta^2 \Expk{ \norms{Gx^{k-1}}^2 }  \\
&& - {~} \frac{L^2\eta^2\hat{C}(1+\mu)}{ \rho } \big[ \Expk{ \norms{x^{k-1} - x^{k-2}}^2 }  - \Expk{ \norms{x^k - x^{k-1}}^2 } \big] \\
&& - {~} \frac{\eta^2(1+\mu)(1-\rho)}{\rho} \big( \Delta_{k-1} - \Delta_{k} \big).
\end{array}
\end{equation*}
Rearranging this inequality and using $\Ec_k$ from \eqref{eq:VROG4NE_Lfuncs}, we obtain \eqref{eq:VROG4NE_key_estimate1}.
\end{proof}
%%% End of Proof.

Now, we are ready to prove our first main result, Theorem~\ref{th:VROG4NE_convergence} in the main text.

%%% Proof of Theorem 3.1.
\begin{proof}[\textbf{Proof of Theorem~\ref{th:VROG4NE_convergence}}]
Let us denote by $M := \frac{\gamma(1 + \mu - \gamma)}{ \mu} + \frac{ (1+\mu)(C + \hat{C})}{ \rho \mu }$.
Then, to keep the right-hand side of  \eqref{eq:VROG4NE_key_estimate1} positive, we need to choose the parameters such that  $L^2\eta^2 \leq \frac{1}{M}$ and $\eta \geq \frac{2\kappa}{2\gamma - 1 - \mu}$.
These two conditions lead to $\frac{4L^2\kappa^2}{(2\gamma - 1 - \mu)^2} \leq L^2\eta^2 \leq \frac{1}{M}$.

Now, for a given $\gamma \in \left(\frac{1}{2}, 1 \right)$, let us choose $\mu := \frac{3(2\gamma - 1)}{4} > 0$.
Then, the last condition holds if $L\kappa \leq \delta := \frac{2\gamma-1}{8\sqrt{M}}$ as stated in Theorem~\ref{th:VROG4NE_convergence}.
In this case, we have $M = \frac{\gamma(1+5\gamma)}{3(2\gamma-1)} + \frac{1 + 6\gamma}{3(2\gamma -1) } \cdot \frac{C + \hat{C}}{\rho}$ as stated in \eqref{eq:VROG4NE_constants}.
Hence, we can choose $\frac{8\kappa}{2\gamma-1} \leq \eta \leq \frac{1}{L\sqrt{M}}$ as claimed in Theorem~\ref{th:VROG4NE_convergence}. % \eqref{eq:SOG_para_choice}.

Next, utilizing $\mu + 1 = \frac{1 + 6\gamma}{2} \geq 1$ and $\mu = \frac{3(2\gamma-1)}{4}$,  \eqref{eq:VROG4NE_key_estimate1} reduces to
\begin{equation*} 
\begin{array}{lcl}
\Expk{\Ec_k} - \Expk{ \Ec_{k+1} } & \geq & \frac{3(2\gamma - 1)}{4} \left( 1 - M \cdot L^2\eta^2  \right) \Expk{ \norms{x^k - x^{k-1}}^2  } + \gamma(1-\gamma)  \eta^2 \Expk{ \norms{Gx^{k-1}}^2 }.
\end{array}
\end{equation*}
Averaging this inequality from $k:=0$ to $k:=K$, we obtain
\begin{equation*} 
\left\{\begin{array}{lcl}
\frac{1}{K+1}\sum_{k=0}^{K} \Expk{ \norms{Gx^{k-1}}^2 }  & \leq &  \frac{ \mathbb{E}[\Ec_0] }{\gamma(1-\gamma)\eta^2(K + 1)}, \vspace{0.5ex}\\
\frac{(1 - M L^2\eta^2)}{K+1}\sum_{k=0}^{K} \Expk{ \norms{x^k - x^{k-1}}^2 } & \leq & \frac{ 4 \mathbb{E}[ \Ec_0] }{3(2\gamma-1)(K + 1)}.
\end{array}\right.
\end{equation*}
Finally, since $x^{-1} = x^{-2} = x^0$, and $\widetilde{S}_{\gamma}^0$ is chosen as $\widetilde{S}_{\gamma}^0 := (1-\gamma)Gx^0$, we have $\Delta_{-1} = \Delta_0 = 0$.
Using this fact, $Gx^{\star} = 0$, the Lipschitz continuity of $G$,  $\rho \in [0, 1]$, and $\gamma < 1$, we can show that
\begin{equation*}
\begin{array}{lcl}
\Expk{ \Ec_0 } & = & \Expk{ \norms{x^0 + \eta\gamma G(x^{0}) - x^{\star}}^2 } +  \frac{ \eta^2(1+\mu)(1-\rho)}{\rho} \Delta_0 \\
& \leq &  2 \Expk{ \norms{x^0 - x^{\star}}^2 } + 2\eta^2\gamma^2 \Expk{ \norms{Gx^0 - Gx^{\star}}^2}  +  \frac{(1+6\gamma) \eta^2}{4\rho} \Delta_0 \\
&  \leq & 2(1 + L^2\eta^2\gamma^2) \Expk{ \norms{x^0 - x^{\star}}^2 }  +  \frac{(1+6\gamma) \eta^2}{4 \rho} \Delta_0 \\
& \leq  &  2\left(1 + L^2\eta^2\right) \norms{x^0 - x^{\star}}^2. 
\end{array}
\end{equation*}
Substituting this upper bound into the above estimates, we get the second bound of \eqref{eq:SOG_convergence_bound1}.
For the first bound, we replace $k-1$ by $k$, and $K$ by $K+1$, using $\norms{Gx^0}^2 \leq 2\norms{Gx^0}^2$, and then multiplying both sides of the result by $\frac{K+2}{K+1}$ to obtain the first line of \eqref{eq:SOG_convergence_bound1}.
\end{proof}
%%% End of Proof.

Next, we restate Corollary \ref{co:SVRG_convergence1_v2} for the case $\gamma \in \left(\frac{1}{2}, 1\right)$ instead of $\gamma = \frac{3}{4}$ as in the main text.
Then, we derive the proof of Corollary \ref{co:SVRG_convergence1_v2} from this result by fixing $\gamma = \frac{3}{4}$.

%%% Corollary 3.2.
\begin{corollary}\label{co:SVRG_convergence1_v2_full}
Suppose that Assumptions~\ref{as:A0}, ~\ref{as:A1}, and ~\ref{as:A2} hold for \eqref{eq:NE} with $\kappa \geq 0$ as in Theorem \ref{th:VROG4NE_convergence}.
Let $\sets{x^k}$ be generated by \eqref{eq:VROG4NE} using the SVRG estimator \eqref{eq:SVRG_estimator}, $\gamma \in \left(\frac{1}{2}, 1\right)$, and  
\begin{equation}\label{eq:SFP_para_choice_v2_full}
\begin{array}{l}
\eta := \frac{1}{L\sqrt{M}},  \quad \text{where} \quad \Lambda :=  \frac{4(1+\gamma^2)(2 - 3\mbf{p}) + 2(3+2\gamma^2)\mbf{p}^2}{b\mbf{p}^2} \ \ \text{and} \ \  M := \frac{\gamma(1+5\gamma)}{3(2\gamma-1)} + \frac{1 + 6\gamma}{3(2\gamma -1) } \cdot \Lambda.
\end{array}
\end{equation}
Then, we have $\eta \geq \frac{\sigma \sqrt{b}\mbf{p}}{L}$ for $\sigma := \frac{\sqrt{3(2\gamma-1)}}{\sqrt{8 + 49\gamma + 13\gamma^2 + 48\gamma^3}}$, and the following bound holds:
\begin{equation}\label{eq:SFP_convergence_bound_v2_full}
\begin{array}{ll}
\frac{1}{K+1}\sum_{k=0}^{K} \Expk{ \norms{Gx^k}^2 }   \leq     \frac{2(1 + L^2\eta^2) R_0^2 }{ \gamma(1-\gamma) \eta^2 (K + 1)},  \quad\text{where} \quad R_0 :=  \norms{x^0 - x^{\star}}.
\end{array}
\end{equation}
For a given tolerance $\epsilon > 0$, if we choose $\mbf{p} :=  n^{-1/3}$ and $b := \lfloor n^{2/3} \rfloor$, then \eqref{eq:VROG4NE} requires $\Tc_{G_i} := n + \big\lfloor\frac{4\Gamma  L^2R_0^2 n^{2/3} }{\epsilon^2} \big\rfloor$ evaluations of $G_i$ to achieve $\frac{1}{K+1}\sum_{k=0}^{K} \Expk{ \norms{Gx^k }^2 } \leq \epsilon^2$, where $\Gamma := \frac{2(5\gamma^2+7\gamma-3)(8 + 49\gamma + 13\gamma^2 + 48\gamma^3)}{3\gamma^2(2\gamma-1)(1-\gamma)(1+5\gamma)}$.
\end{corollary}


%%% Proof of Corollary 3.2.
\begin{proof}% [\textbf{Proof of Corollary~\ref{co:SVRG_convergence1_v2}}]
For the SVRG estimator \eqref{eq:SVRG_estimator}, by Lemma~\ref{le:SVGR_estimator}, we have $\rho := \frac{\mbf{p}}{2} \in (0, 1]$, $C := \frac{4-6\mbf{p} + 3\mbf{p}^2 }{b\mbf{p}}$, and $\hat{C} := \frac{2\gamma^2(2 - 3\mbf{p} + \mbf{p}^2)}{b\mbf{p}}$.
Therefore, we can compute $\Lambda := \frac{C+\hat{C}}{\rho} = \frac{4(1+\gamma^2)(2 - 3\mbf{p}) + 2(3+2\gamma^2)\mbf{p}^2}{b\mbf{p}^2} \leq \frac{8(1+\gamma^2)}{b\mbf{p}^2}$, and $M$ in Theorem~\ref{th:VROG4NE_convergence} as $M :=    \frac{\gamma(1+5\gamma)}{3(2\gamma-1)} + \frac{1 + 6\gamma}{3(2\gamma -1) } \cdot \Lambda \leq \frac{\gamma(1+5\gamma)}{3(2\gamma-1)} + \frac{8(1 + 6\gamma)(1+\gamma^2)}{3(2\gamma -1) b\mbf{p}^2 }$ as stated in \eqref{eq:SFP_para_choice_v2_full}.
The estimate \eqref{eq:SFP_convergence_bound_v2_full} is exactly the first line of \eqref{eq:SOG_convergence_bound1}.

Now, suppose that $b\mbf{p}^2 \leq 1$.
Then, by \eqref{eq:SFP_para_choice_v2_full}, we have $M \leq  \frac{8 + 49\gamma + 13\gamma^2 + 48\gamma^3}{3(2\gamma -1) b\mbf{p}^2 }$.
Therefore, if we choose $\eta := \frac{1}{L\sqrt{M}}$, then $\eta$ satisfies the conditions of Theorem~\ref{th:VROG4NE_convergence}, provided that $L\rho \leq \delta$.
Moreover, we have $\eta \geq \frac{\sqrt{3(2\gamma-1)}\sqrt{b}\mbf{p}}{\sqrt{8 + 49\gamma + 13\gamma^2 + 48\gamma^3}} = \frac{\sigma \sqrt{b}\mbf{p}}{L}$, where $\sigma := \frac{\sqrt{3(2\gamma-1)}}{\sqrt{8 + 49\gamma + 13\gamma^2 + 48\gamma^3}}$.

From \eqref{eq:SFP_convergence_bound_v2_full}, to guarantee $\frac{1}{K+1}\sum_{k=0}^K\Exp{\norms{Gx^k}^2} \leq \epsilon^2$, we need to impose $\frac{2(1 + L^2\eta^2)R_0^2}{\gamma(1-\gamma) \eta^2(K+1)} \leq \epsilon^2$, where $R_0 := \norms{x^0 - x^{\star}}$.
However, since $1+ L^2\eta^2 = 1 + \frac{1}{M} \leq \frac{5\gamma^2 + 7\gamma-3}{\gamma(1+5\gamma)}$ and $\eta \geq \frac{\sigma \sqrt{b}\mbf{p}}{ L }$, the last condition holds if we choose $K :=  \left\lfloor \Gamma \cdot \frac{L^2 R_0^2}{ b\mbf{p}^2\epsilon^2} \right\rfloor$, where $\Gamma := \frac{2(5\gamma^2+7\gamma-3)}{\sigma^2\gamma^2(1-\gamma)(1+5\gamma)} = \frac{2(5\gamma^2+7\gamma-3)(8 + 49\gamma + 13\gamma^2 + 48\gamma^3)}{3\gamma^2(2\gamma-1)(1-\gamma)(1+5\gamma)}$.

Finally, note that, at each iteration $k$, \eqref{eq:VROG4NE} requires $3$ mini-batches of size $b$, and occasionally compute the full batch $Gw^k$, leading to the cost of $n\mbf{p} + 3b$.
The total complexity is 
\begin{equation*}
\begin{array}{lcl}
\Tc_{c} & := & K(n\mbf{p} + 3b) = \frac{ \Gamma L^2 R_0^2(n\mbf{p} + 3b)}{ b\mbf{p}^2\epsilon^2} = \frac{\Gamma L^2R_0^2}{\epsilon^2}\big(\frac{n}{b\mbf{p}} + \frac{3}{ \mbf{p}^2} \big).
\end{array}
\end{equation*}
If we choose $b := \lfloor n^{2/3} \rfloor$ and $\mbf{p} := n^{-1/3}$, then $b\mbf{p}^2 = 1$ and $\Tc_{c} = \frac{4\Gamma n^{2/3} L^2R_0^2}{\epsilon^2}$.
For the SVRG estimator \eqref{eq:SVRG_estimator}, one needs to compute $Gw^0$, which requires $n$ evaluations of $G_i$.
Hence, the total complexity of the algorithm is $\Tc_{G_i} := n + \left\lfloor\frac{4\Gamma n^{2/3} L^2R_0^2}{\epsilon^2}\right\rfloor$ as stated.
\end{proof}
%%% End of Proof.

%%% Proof of Corollary 3.2.
\begin{proof}[\textbf{Proof of Corollary~\ref{co:SVRG_convergence1_v2}}]
Since we fix $\gamma := \frac{3}{4}$, we can easily compute $\sigma := \frac{\sqrt{3(2\gamma-1)}}{\sqrt{8 + 49\gamma + 13\gamma^2 + 48\gamma^3}} \approx 0.144025 \geq 0.1440$ and $\Gamma := \frac{2(5\gamma^2+7\gamma-3)(8 + 49\gamma + 13\gamma^2 + 48\gamma^3)}{3\gamma^2(2\gamma-1)(1-\gamma)(1+5\gamma)} \approx 730.736842 \leq 731$.
Therefore, we obtain $\eta \geq \frac{0.1440 \sqrt{b}\mbf{p}}{L}$ and $\Tc_{G_i} := n + \big\lfloor \frac{4\Gamma n^{2/3} L^2R_0^2}{\epsilon^2} \big\rfloor$, where $\Gamma := 731$.
Moreover, \eqref{eq:SFP_convergence_bound_v2_full} reduces to $\frac{1}{K+1}\sum_{k=0}^{K} \Expk{ \norms{Gx^k}^2 }   \leq     \frac{32(1 + 0.1440^2) L^2 R_0^2 }{3\cdot 0.1440^2 b\mbf{p}^2 (K + 1)} \leq \frac{526 \cdot L^2R_0^2}{b\mbf{p}^2(K+1)}$.
\end{proof}
%%% End of Proof.


Finally, we also restate Corollary~\ref{co:SVRG_convergence1_v3} for the case $\gamma \in \left(\frac{1}{2}, 1\right)$ and then derive the oracle complexity of Corollary~\ref{co:SVRG_convergence1_v3}  from this result by fixing $\gamma := \frac{3}{4}$.

%%% Corollary 3.3.
\begin{corollary}\label{co:SVRG_convergence1_v3_full}
Suppose that Assumptions~\ref{as:A0}, ~\ref{as:A1}, and ~\ref{as:A2} hold for \eqref{eq:NE} with $\kappa \geq 0$ as in Theorem~\ref{th:VROG4NE_convergence}.
Let $\sets{x^k}$ be generated by \eqref{eq:VROG4NE} using the SAGA estimator \eqref{eq:SAGA_estimator}, $\gamma \in \left(\frac{1}{2}, 1\right)$, and
\begin{equation}\label{eq:SFP_para_choice_v3_full}
\begin{array}{l}
\eta := \frac{1}{L\sqrt{M}}, \quad \text{where} \quad \Lambda := \frac{2}{b} + \frac{4(1+\gamma^2)(n-b)(2n+b)}{b^3} \quad \text{and} \quad  M := \frac{\gamma(1+5\gamma)}{3(2\gamma-1)} + \frac{1 + 6\gamma}{3(2\gamma -1) } \cdot \Lambda.
\end{array}
\end{equation}
Then, we have $\eta \geq \frac{\sigma b^{3/2}}{n L}$ for  $\sigma :=  \frac{\sqrt{3(2\gamma-1)}}{\sqrt{10 + 61\gamma + 13\gamma^2 + 48\gamma^3}}$, and the following bound holds:
\begin{equation}\label{eq:SFP_convergence_bound_v3_full}
\begin{array}{ll}
\frac{1}{K+1}\sum_{k=0}^{K} \Expk{ \norms{Gx^k}^2 }   \leq     \frac{2(1 + L^2\eta^2) R_0^2 }{ \gamma(1-\gamma) \eta^2 (K + 1)},  \quad\text{where} \quad R_0 :=  \norms{x^0 - x^{\star}}.
\end{array}
\end{equation}
Moreover, for a given tolerance $\epsilon > 0$, if we choose $b := \lfloor n^{2/3} \rfloor$, then  \eqref{eq:VROG4NE} requires $\Tc_{G_i} := n + \big\lfloor \frac{ 3\Gamma L^2R_0^2 n^{2/3}}{ \varepsilon^2} \big\rfloor$ evaluations of $G_i$ to achieve $\frac{1}{K+1}\sum_{k=0}^{K} \Expk{ \norms{Gx^k }^2 } \leq \epsilon^2$, where $\Gamma := \frac{2(7\gamma + 5\gamma^2-3)(10 + 61\gamma + 13\gamma^2 + 48\gamma^3)}{3\gamma^2(1-\gamma)(2\gamma-1)(1+5\gamma)}$.
\end{corollary}

%%% Proof of Corollary 3.3.
\begin{proof}
Since we use the SAGA estimator \eqref{eq:SAGA_estimator}, we have  $\rho := \frac{b}{2n} \in (0, 1]$, $C :=  \frac{[ 2(n - b)(2n+b) + b^2] }{ n b^2}$, and $\hat{C} := \frac{2(n - b)(2n+b) \gamma^2}{ nb^2}$.
In this case, since $b \geq 1$, we can easily show that $\Lambda := \frac{C + \hat{C}}{\rho} = \frac{2}{b} + \frac{4(1+\gamma^2)(n-b)(2n+b)}{b^3} \leq 2 + \frac{8(1+\gamma^2)n^2}{b^3}$.
Hence, $M$ in Theorem~\ref{th:VROG4NE_convergence} reduces to 
\begin{equation*}
\begin{array}{lcl}
M := \frac{\gamma(1+5\gamma)}{3(2\gamma-1)} + \frac{1 + 6\gamma}{3(2\gamma -1) } \cdot \Lambda \leq  \frac{2 + 13\gamma + 5\gamma^2}{3(2\gamma-1)} +  \frac{8(1+\gamma^2)(1 + 6\gamma)n^2}{3(2\gamma -1) b^3}.
\end{array}
\end{equation*}
Suppose that $1 \leq b \leq n^{2/3}$.
Then, one can prove that $M \leq \left[ \frac{2 + 13\gamma + 5\gamma^2 }{3(2\gamma-1)} +  \frac{8(1+\gamma^2)(1 + 6\gamma)}{3(2\gamma -1)} \right] \frac{n^2}{b^3} = \frac{(10 + 61\gamma + 13\gamma^2 + 48\gamma^3)n^2}{3(2\gamma-1)b^3} = \frac{n^2}{ \sigma^2b^3}$, where $\sigma :=  \frac{\sqrt{3(2\gamma-1)}}{\sqrt{10 + 61\gamma + 13\gamma^2 + 48\gamma^3}}$.
Hence, if we choose $\eta := \frac{1}{L\sqrt{M}}$, then we get $\eta \geq \frac{\sigma b^{3/2}}{nL}$ as stated.
Moreover, we obtain \eqref{eq:SFP_convergence_bound_v2_full} from the first line of \eqref{eq:SOG_convergence_bound1} as before.

Now, for $\eta := \frac{1}{L\sqrt{M}} \geq \frac{\sigma b^{3/2}}{nL}$, from \eqref{eq:SFP_convergence_bound_v2_full}, to guarantee $\frac{1}{K+1}\sum_{k=0}^K\Exp{\norms{Gx^k}^2} \leq \epsilon^2$, we need to impose $\frac{2(1 + L^2\eta^2)R_0^2}{\gamma(1-\gamma) \eta^2(K+1)} \leq \epsilon^2$, where $R_0 := \norms{x^0 - x^{\star}}$.
Since $1+ L^2\eta^2 = 1 + \frac{1}{M} \leq \frac{7\gamma + 5\gamma^2-3}{\gamma(1+5\gamma)}$ and $\eta \geq  \frac{\sigma b^{3/2}}{n L}$, the last condition holds if we choose $K :=  \left\lfloor \Gamma \cdot \frac{L^2 R_0^2 n^2}{ b^3 \epsilon^2} \right\rfloor$, where $\Gamma := \frac{2(7\gamma + 5\gamma^2-3)}{\sigma^2\gamma^2(1-\gamma)(1+5\gamma)} = \frac{2(7\gamma + 5\gamma^2-3)(10 + 61\gamma + 13\gamma^2 + 48\gamma^3)}{3\gamma^2(1-\gamma)(2\gamma-1)(1+5\gamma)}$.

Finally, at each iteration $k$, \eqref{eq:VROG4NE} requires $3$ mini-batches of size $b$, leading to the cost of $3b$ per iteration.
Hence, the total complexity is 
\begin{equation*}
\begin{array}{lcl}
\Tc_{c} & := &  3bK = \big\lfloor \frac{ 3\Gamma L^2 R_0^2 n^2 }{ b^2 \epsilon^2} \big\rfloor.
\end{array}
\end{equation*}
If we choose $b := \lfloor n^{2/3} \rfloor$, then $\Tc_{c} = \big\lfloor \frac{3\Gamma  L^2R_0^2 n^{2/3} }{\epsilon^2} \big\rfloor$.
For the SAGA estimator \eqref{eq:SAGA_estimator}, one needs to compute $Gw^0$, which requires $n$ evaluations of $G_i$.
Hence, the total complexity of the algorithm is $\Tc_{G_i} := n + \big\lfloor\frac{3\Gamma L^2R_0^2  n^{2/3}  }{\epsilon^2} \big\rfloor$.
\end{proof}
%%% End of Proof.

%%% Proof of Corollary 3.2.
\begin{proof}[\textbf{Proof of Corollary~\ref{co:SVRG_convergence1_v3}}]
Since we choose $\gamma := \frac{3}{4}$, we have $\sigma := \frac{\sqrt{3(2\gamma-1)}}{\sqrt{10 + 61\gamma + 13\gamma^2 + 48\gamma^3}} = 0.14948 \geq 0.1494$ and $\Gamma := \frac{2(7\gamma + 5\gamma^2-3)(10 + 61\gamma + 13\gamma^2 + 48\gamma^3)}{3\gamma^2(1-\gamma)(2\gamma-1)(1+5\gamma)} = 2815.8 \leq 2816$.
Applying the results of Corollary~\ref{co:SVRG_convergence1_v3_full}, we obtain our conclusions in Corollary~\ref{co:SVRG_convergence1_v3}.
Moreover, \eqref{eq:SFP_convergence_bound_v3_full} reduces to $\frac{1}{K+1}\sum_{k=0}^{K} \Expk{ \norms{Gx^k}^2 }   \leq     \frac{32(1 + 0.494^2) L^2 R_0^2 }{3\cdot 0.1494^2 b\mbf{p}^2 (K + 1)} \leq \frac{489 \cdot L^2R_0^2}{b\mbf{p}^2(K+1)}$ as stated.
\end{proof}
%%%% End of Proof.

%%%%%%%%%%%%%%%%%%%%%%%%%%%%%%%%%%%%%%%%%%%%%%%%%%%%%%%
%%% C. Convergence Analysis of VROG2 for NI: Technical Proofs.
%%%%%%%%%%%%%%%%%%%%%%%%%%%%%%%%%%%%%%%%%%%%%%%%%%%%%%%
\beforesec
\section{Convergence Analysis of \ref{eq:VROG4NI} for \eqref{eq:NI}: Technical Proofs}\label{apdx:sec:VROG4NI}
\aftersec
One key step to analyze the convergence of \eqref{eq:VROG4NI} is to construct an appropriate potential function.
For this purpose, we introduce the following function:
\vspace{-0.5ex}
\begin{equation}\label{eq:VROG4NI_Lfuncs}
\begin{array}{lcl}
\Lc_k & := &  \norms{x^k + \gamma\eta (Gx^{k-1} + v^{k}) - x^{\star}}^2 + \mu \norms{x^k - x^{k-1} + \gamma\eta(Gx^{k-1} + v^k) }^2,
\end{array}
\vspace{-0.25ex}
\end{equation}
where $\mu > 0$ is a given parameter  and $v^k \in Tx^k$ is given.
This function is then combined with $\Ec_k$ from \eqref{eq:VROG4NE_Lfuncs} to establish the convergence of \eqref{eq:VROG4NI}.

Let us first state and prove Lemma~\ref{le:VROG2_key_bounds}, which provides a key estimate for our convergence analysis of \eqref{eq:VROG4NI} in Theorem~\ref{th:VROG4NI_convergence}.

%%% Lemma 3.1.
\begin{lemma}\label{le:VROG2_key_bounds}
Suppose that Assumption~\ref{as:A1} holds for \eqref{eq:NI}.
Let $\sets{x^k}$ be generated by \eqref{eq:VROG4NI}, $\Lc_k$ be defined by \eqref{eq:VROG4NI_Lfuncs}, and $\Ec_k$ be defined by  \eqref{eq:VROG4NE_Lfuncs}.
Then, we have
\begin{equation}\label{eq:VROG4NI_key_est1} 
\hspace{-4ex}
\begin{array}{lcl}
\Lc_k - \Expsk{k}{ \Lc_{k+1} }
& \geq & 2(1-\gamma)\eta\iprods{Gx^{k} + v^k, x^k - x^{\star}}  + (1+\mu)(1-\gamma)(2\gamma - 1)\eta^2 \norms{Gx^k + v^k}^2 \\
&& + {~}  \gamma[ 1-\gamma - \mu(3\gamma-1)]\eta^2\norms{Gx^{k-1} + v^k}^2 - (1+\mu)\eta^2 \Expsk{k}{ \norms{e^k}^2 } \\
&& + {~} \frac{1}{2}\big[ \mu -  2(1+\mu)\gamma(1-\gamma)L^2\eta^2 \big] \norms{x^k - x^{k-1} }^2.
\end{array} 
\hspace{-4ex}
\end{equation}
If, additionally, Assumption~\ref{as:A2} holds for \eqref{eq:NI}, then we have
\begin{equation}\label{eq:VROG4NI_key_est2}
\hspace{-4ex}
\begin{array}{lcl}
\Expk{ \Ec_k } - \Expk{ \Ec_{k+1} } 
& \geq &  \frac{1}{2} \Big[ \mu -  2(1+\mu)\gamma(1-\gamma)L^2\eta^2 - \frac{2L^2\eta^2(1+\mu)(C + \hat{C})}{ \rho } \Big] \Expk{ \norms{x^k - x^{k-1}}^2 } \\
&& + {~} \gamma[ 1 - \gamma  - \mu(3\gamma-1)]\eta^2 \Expk{ \norms{Gx^{k-1} + v^k }^2 }  \\
&& + {~} (1-\gamma)\eta \big[  (1+\mu)(2\gamma - 1)\eta - 2\kappa \big]  \Expk{ \norms{Gx^k + v^k}^2 }.
\end{array}
\hspace{-4ex}
\end{equation}
\end{lemma}

%%% Proof of Lemma 3.1.
\begin{proof}
Let us introduce two notations $w^k := Gx^k + v^k$ and $\hat{w}^k := Gx^{k-1} + v^k$, where $v^k \in Tx^k$.
We also recall $S^k_{\gamma} := Gx^k - \gamma Gx^{k-1}$ and $e^k := \widetilde{S}^k_{\gamma} - S_{\gamma}^k$ from \eqref{eq:Sk_op}.
Then, it is obvious that $\widetilde{S}^k_{\gamma} = S_{\gamma}^k + e^k = Gx^k - \gamma Gx^{k-1} + e^k$.

Now, using $\widetilde{S}^k_{\gamma} =  Gx^k - \gamma Gx^{k-1} + e^k$, it follows from \eqref{eq:VROG4NI} that
\begin{equation}\label{eq:VROG4NI_key_bound_proof1} 
\begin{array}{lcl}
x^{k+1}  & = & x^k - \eta  \widetilde{S}_{\gamma}^k - \gamma\eta  v^{k+1} - (2\gamma - 1)\eta v^k \\
%& = & x^k - \eta Gx^k + \eta\gamma Gx^{k-1} - \gamma\eta  v^{k+1} - (2\gamma - 1)\eta v^k -  \eta e^k \\
& = & x^k - \gamma\eta(Gx^k + v^{k+1}) -  (1-\gamma)\eta ( Gx^k + v^k) + \eta\gamma ( Gx^{k-1} +  v^k) -  \eta e^k \\
& = & x^k - \gamma\eta\hat{w}^{k+1} -   (1-\gamma)\eta w^k + \gamma\eta \hat{w}^k -  \eta e^k.
\end{array} 
\end{equation}
Then, using \eqref{eq:VROG4NI_key_bound_proof1} and $\hat{w}^{k+1} = Gx^k + v^{k+1}$, we can show that
\begin{equation*} 
\begin{array}{lcl}
\Tc_{[1]} &:= & \norms{x^{k+1} + \gamma\eta(Gx^{k} + v^{k+1}) - x^{\star}}^2 = \norms{x^{k+1} - x^{\star} + \gamma \eta \hat{w}^{k+1}}^2 \\
& \overset{\tiny\eqref{eq:VROG4NI_key_bound_proof1}}{ = } & \norms{ x^k - \gamma \eta \hat{w}^{k+1} -  (1-\gamma)\eta w^k + \gamma \eta \hat{w}^k -  \eta e^k - x^{\star} + \gamma \eta \hat{w}^{k+1} }^2 \\
%& = & \norms{ x^k - x^{\star} -  (1-\gamma)\eta w^k + \gamma\eta \hat{w}^k -  \eta e^k }^2 \\
%& = & \norms{x^k - x^{\star}}^2 + \eta^2\norms{ (1-\gamma) w^k - \gamma\hat{w}^k}^2 + \eta^2\norms{e^k}^2 \\
%&& - {~} 2\eta\iprods{(1-\gamma) w^k - \gamma \hat{w}^k, x^k - x^{\star}}  -  2\eta\iprods{e^k, x^k - x^{\star}}  + 2\eta^2 \iprods{e^k, (1-\gamma) w^k - \gamma \hat{w}^k} \\
& = & \norms{x^k - x^{\star}}^2  - 2(1-\gamma)\eta \iprods{w^k, x^k - x^{\star}} + 2\gamma\eta\iprods{\hat{w}^k, x^k - x^{\star}}  + \eta^2\norms{e^k}^2 \\
&& + {~}  (1-\gamma)^2\eta^2\norms{w^k}^2 - 2\gamma(1-\gamma)\eta^2\iprods{w^k, \hat{w}^k} + \gamma^2\eta^2\norms{\hat{w}^k}^2 \\
&& -  {~}  2\eta\iprods{e^k, x^k - x^{\star}}  + 2(1-\gamma)\eta^2 \iprods{e^k, w^k} - 2\gamma\eta^2 \iprods{e^k, \hat{w}^k}.
\end{array} 
\end{equation*}
Alternatively, using $\hat{w}^k = Gx^{k-1} + v^k$, we also have
\begin{equation*} 
\begin{array}{lcl}
\Tc_{[2]} &:= & \norms{x^k + \gamma \eta(Gx^{k-1} + v^k) - x^{\star}}^2 = \norms{x^k - x^{\star} + \gamma\eta\hat{w}^k }^2 \\
& = & \norms{x^k - x^{\star}}^2 + 2\gamma\eta\iprods{\hat{w}^k, x^k - x^{\star}} + \gamma^2\eta^2\norms{\hat{w}^k}^2.
\end{array} 
\end{equation*}
Subtracting $\Tc_{[1]}$ from $\Tc_{[2]}$, we can show that
\begin{equation}\label{eq:VROG4NI_key_bound_proof2} 
\begin{array}{lcl}
\Tc_{[3]} &:= & \norms{x^k + \gamma \eta(Gx^{k-1} + v^k) - x^{\star}}^2 - \norms{x^{k+1} + \gamma\eta(Gx^k + v^{k+1}) - x^{\star}}^2 \\
& = & 2(1-\gamma)\eta\iprods{w^k, x^k - x^{\star}} -  (1-\gamma)^2\eta^2\norms{w^k}^2 + 2\gamma(1-\gamma)\eta^2\iprods{w^k, \hat{w}^k} \\
&& +  {~}  2\eta\iprods{e^k, x^k - x^{\star}}  - 2(1-\gamma)\eta^2 \iprods{e^k, w^k} + 2\gamma\eta^2 \iprods{e^k, \hat{w}^k} - \eta^2\norms{e^k}^2 \\
%& = & 2(1-\gamma)\eta\iprods{w^k, x^k - x^{\star}} -  (1-\gamma)^2\eta^2\norms{w^k}^2 \\
%&& + {~} \gamma(1-\gamma)\eta^2 \big[ \norms{w^k}^2 + \norms{\hat{w}^k}^2 - \norms{w^k - \hat{w}^k}^2 \big] \\
%&& +  {~}  2\eta\iprods{e^k, x^k - x^{\star}}  - 2(1-\gamma)\eta^2 \iprods{e^k, w^k} + 2\gamma\eta^2 \iprods{e^k, \hat{w}^k} - \eta^2\norms{e^k}^2 \\
& = & 2(1-\gamma)\eta \iprods{w^k, x^k - x^{\star}}  +  (1-\gamma)(2\gamma - 1) \eta^2 \norms{w^k}^2 \\
&& + {~}  \gamma(1-\gamma)\eta^2\norms{\hat{w}^k}^2 -  \gamma(1-\gamma)\eta^2\norms{w^k - \hat{w}^k}^2 \\
&& +  {~}  2\eta\iprods{e^k, x^k - x^{\star}}  - 2(1-\gamma)\eta^2 \iprods{e^k, w^k} + 2\gamma\eta^2 \iprods{e^k, \hat{w}^k} - \eta^2\norms{e^k}^2.
\end{array} 
\end{equation}
Next, using again $\hat{w}^{k+1} = Gx^k + v^{k+1}$ and \eqref{eq:VROG4NI_key_bound_proof1}, we have
\begin{equation*} 
\begin{array}{lcl}
\Tc_{[4]} &:= & \norms{x^{k+1} - x^k + \gamma \eta(Gx^k + v^{k+1}) }^2 = \norms{x^{k+1} - x^k + \gamma \eta\hat{w}^{k+1} }^2 \\ 
& \overset{\tiny\eqref{eq:VROG4NI_key_bound_proof1}}{ = } & \eta^2 \norms{(1-\gamma) w^k - \gamma \hat{w}^k + e^k}^2 \\
& = & (1-\gamma)^2 \eta^2 \norms{w^k}^2 - 2\gamma(1-\gamma)\eta^2\iprods{w^k, \hat{w}^k} + \gamma^2\eta^2\norms{\hat{w}^k}^2 \\
&& + {~} \eta^2\norms{e^k}^2 + 2(1-\gamma)\eta^2\iprods{e^k, w^k} - 2\gamma\eta^2\iprods{e^k, \hat{w}^k} \\
%& = & (1-\gamma)^2\eta^2 \norms{w^k}^2 + \gamma^2\eta^2\norms{\hat{w}^k}^2 - \gamma(1-\gamma)\eta^2\big[ \norms{w^k}^2 + \norms{\hat{w}^k}^2 - \norms{w^k - \hat{w}^k}^2 \big] \\
%&& + {~} \eta^2\norms{e^k}^2 + 2(1-\gamma)\eta^2\iprods{e^k, w^k} - 2\gamma\eta^2\iprods{e^k, \hat{w}^k} \\
& = & -(1-\gamma)(2\gamma-1)\eta^2 \norms{w^k}^2 + \gamma(2\gamma-1)\eta^2\norms{\hat{w}^k}^2 + \gamma(1-\gamma)\eta^2 \norms{w^k - \hat{w}^k}^2 \\
&& + {~} \eta^2\norms{e^k}^2 + 2(1-\gamma)\eta^2\iprods{e^k, w^k} - 2\gamma\eta^2\iprods{e^k, \hat{w}^k}.
\end{array} 
\end{equation*}
Moreover, by the Cauchy-Schwarz inequality in $\myeqc{1}$ and Young's inequality in $\myeqc{2}$, we can prove that
\begin{equation*} 
\begin{array}{lcl}
\norms{x^k - x^{k-1} + \gamma \eta \hat{w}^k }^2 & = & \norms{x^k - x^{k-1}}^2 + 2\gamma\eta\iprods{\hat{w}^k, x^k - x^{k-1}} + \gamma^2\eta^2\norms{\hat{w}^k}^2 \\
& \overset{\myeqc{1}}{ \geq } & \norms{x^k - x^{k-1}}^2 - 2\gamma\eta\norms{\hat{w}^k}\norms{x^k - x^{k-1}} + \gamma^2\eta^2\norms{\hat{w}^k}^2 \\
& \overset{\myeqc{2}}{ \geq } & \frac{1}{2} \norms{x^k - x^{k-1}}^2 - \gamma^2\eta^2 \norms{\hat{w}^k}^2.
\end{array} 
\end{equation*}
Combining the last two expressions, we can show that
\begin{equation*} 
\begin{array}{lcl}
\Tc_{[5]} &:= & \norms{x^k - x^{k-1} + \gamma \eta(Gx^{k-1} + v^k) }^2 - \norms{x^{k+1} - x^k + \gamma \eta(Gx^{k} + v^{k+1}) }^2 \\
& = & \norms{x^k - x^{k-1} + \gamma \eta \hat{w}^k }^2 - \norms{x^{k+1} - x^k + \gamma \eta\hat{w}^{k+1} }^2 \\ 
%& \geq & \frac{1}{2}\norms{x^k - x^{k-1}}^2 -  \gamma^2 \eta^2  \norms{\hat{w}^k}^2  + (1-\gamma)(2\gamma-1)\eta^2 \norms{w^k}^2 - \gamma(2\gamma-1)\eta^2\norms{\hat{w}^k}^2 \\
%&& - {~} \gamma(1-\gamma)\eta^2 \norms{w^k - \hat{w}^k}^2 - \eta^2\norms{e^k}^2 - 2(1-\gamma)\eta^2\iprods{e^k, w^k} + 2\gamma\eta^2\iprods{e^k, \hat{w}^k} \\
& \geq & \frac{1}{2} \norms{x^k - x^{k-1}}^2 + (1-\gamma)(2\gamma-1)\eta^2 \norms{w^k}^2 - \gamma(3\gamma-1)\eta^2\norms{\hat{w}^k}^2 \\
&& - {~} \gamma(1-\gamma)\eta^2 \norms{w^k - \hat{w}^k}^2 - \eta^2\norms{e^k}^2 - 2(1-\gamma)\eta^2\iprods{e^k, w^k} + 2\gamma\eta^2\iprods{e^k, \hat{w}^k}.
\end{array} 
\end{equation*}
Multiplying $\Tc_{[5]}$ by $\mu > 0$, and adding the result to \eqref{eq:VROG4NI_key_bound_proof2}, and using $\Lc_k$ from \eqref{eq:VROG4NI_Lfuncs}, we have
\begin{equation*} 
\begin{array}{lcl}
\Lc_k - \Lc_{k+1} & = &  \norms{x^k + \gamma\eta(Gx^{k-1} + v^k) - x^{\star}}^2 - \norms{x^{k+1} + \gamma \eta(Gx^k + v^{k+1}) - x^{\star}}^2 \\
&& + {~} \mu \norms{x^k - x^{k-1} + \gamma \eta(Gx^{k-1} + v^k) }^2 - \mu \norms{x^{k+1} - x^k + \gamma \eta(Gx^k + v^{k+1}) }^2 \\
%%%
%& \geq &  2(1-\gamma)\eta \iprods{w^k, x^k - x^{\star}}  +  (1-\gamma)(2\gamma - 1) \eta^2 \norms{w^k}^2 \\
%&& + {~}  \gamma(1-\gamma)\eta^2\norms{\hat{w}^k}^2 -  \gamma(1-\gamma)\eta^2\norms{w^k - \hat{w}^k}^2 \\
%&& +  {~}  2\eta\iprods{e^k, x^k - x^{\star}}  - 2(1-\gamma)\eta^2 \iprods{e^k, w^k} + 2\gamma\eta^2 \iprods{e^k, \hat{w}^k} - \eta^2\norms{e^k}^2 \\
%%%%
%&& + {~} \frac{\mu}{2} \norms{x^k - x^{k-1}}^2 + \mu (1-\gamma)(2\gamma-1)\eta^2 \norms{w^k}^2 - \mu \gamma(3\gamma-1)\eta^2\norms{\hat{w}^k}^2 \\
%&& - {~} \mu \gamma(1-\gamma)\eta^2 \norms{w^k - \hat{w}^k}^2 - \mu \eta^2\norms{e^k}^2 - 2\mu (1-\gamma)\eta^2\iprods{e^k, w^k} + 2\mu \gamma\eta^2\iprods{e^k, \hat{w}^k} \\
%% 
& \geq & 2(1-\gamma)\eta \iprods{w^k, x^k - x^{\star}}  +  \frac{\mu}{2}\norms{x^k - x^{k-1}}^2 + (1+\mu)(1-\gamma)(2\gamma - 1)\eta^2 \norms{w^k}^2 \\
&& + {~}  \gamma[ (1-\gamma) - \mu(3\gamma-1)]\eta^2\norms{\hat{w}^k}^2 -  (1+\mu)\gamma(1-\gamma)\eta^2\norms{w^k - \hat{w}^k}^2 \\
&& +  {~}  2\eta\iprods{e^k, x^k - x^{\star}}  - 2(1+\mu)(1-\gamma)\eta^2 \iprods{e^k, w^k} \\
&& + {~} 2(1+\mu)\gamma\eta^2\iprods{e^k, \hat{w}^k} - (1+\mu)\eta^2\norms{e^k}^2.
\end{array} 
\end{equation*}
Taking the conditional expectation $\Expsk{k}{\cdot}$ both sides of this expression, and noting that
\begin{equation*}
\begin{array}{lcl}
\Expsk{k}{ \iprods{e^k, x^k - x^{\star}} } & = & \iprods{\Expsk{k}{e^k}, x^k - x^{\star}} = 0, \\
\Expsk{k}{ \iprods{e^k, w^k} } & = & \iprods{\Expsk{k}{e^k}, w^k} = 0, \\
\Expsk{k}{ \iprods{e^k, \hat{w}^k} } & = & \iprods{\Expsk{k}{e^k}, \hat{w}^k} = 0,
\end{array}
\end{equation*}
we obtain
\begin{equation*} 
\begin{array}{lcl}
\Lc_k - \Expsk{k}{ \Lc_{k+1} }
& \geq & 2(1-\gamma)\eta \iprods{w^k, x^k - x^{\star}}  +  \frac{\mu}{2}\norms{x^k - x^{k-1}}^2 + (1+\mu)(1-\gamma)(2\gamma - 1)\eta^2 \norms{w^k}^2 \\
&& + {~}  \gamma[ (1-\gamma) - \mu(3\gamma-1)]\eta^2\norms{\hat{w}^k}^2 -  (1+\mu)\gamma(1-\gamma)\eta^2\norms{w^k - \hat{w}^k}^2 \\
&& - {~} (1+\mu)\eta^2 \Expsk{k}{ \norms{e^k}^2 }.
\end{array} 
\end{equation*}
Finally, by the $L$-Lipschitz continuity of $G$ from \eqref{eq:L_smooth} of Assumption~\ref{as:A1}, we have $\norms{w^k - \hat{w}^k}^2 = \norms{Gx^k - Gx^{k-1}}^2 \leq L^2\norms{x^k - x^{k-1}}^2$ as shown in \eqref{eq:L_smooth2}.
Using this inequality into the last estimate, we can show that
\begin{equation*} 
\begin{array}{lcl}
\Lc_k - \Expsk{k}{ \Lc_{k+1} }
& \geq & 2(1-\gamma)\eta\iprods{w^k, x^k - x^{\star}}  + (1+\mu)(1-\gamma)(2\gamma - 1)\eta^2 \norms{w^k}^2 \\
&& + {~}  \gamma[ 1 - \gamma  - \mu(3\gamma-1)]\eta^2 \norms{\hat{w}^k}^2 - (1+\mu)\eta^2 \Expsk{k}{ \norms{e^k}^2 } \\
&& + {~} \frac{1}{2}\big[ \mu -  2(1+\mu)\gamma(1-\gamma)L^2\eta^2 \big] \norms{x^k - x^{k-1} }^2,
\end{array} 
\end{equation*}
which proves \eqref{eq:VROG4NI_key_est1} by recalling $w^k := Gx^k + v^k$ and $\hat{w}^k := Gx^{k-1} + v^k$.

Taking the full expectation of \eqref{eq:VROG4NI_key_est1} and using $\iprods{Gx^k + v^k, x^k - x^{\star}} \geq -\kappa \norms{Gx^k + v^k}^2$ from Assumption~\ref{as:A2} and $\Expsk{k}{ \norms{e^k}^2 } \leq  \Delta_k$ from \eqref{eq:ub_SG_estimator}, we can bound it as
\begin{equation}\label{eq:VROG4NI_key_bound_proof5} 
\hspace{-2ex}
\begin{array}{lcl}
\Expk{ \Lc_k } - \Expk{ \Lc_{k+1} } & \geq &  {~}   \frac{1}{2}\big[ \mu -  2(1+\mu)\gamma(1-\gamma)L^2\eta^2 \big] \Expk{ \norms{x^k - x^{k-1}}^2 } - (1 + \mu)\eta^2 \Delta_k \\
&& + {~} \gamma[ 1 - \gamma  - \mu(3\gamma-1)]\eta^2 \Expk{ \norms{Gx^{k-1} + v^k }^2 }  \\
&& + {~} (1-\gamma)\eta \big[  (1+\mu)(2\gamma - 1)\eta - 2\kappa \big]  \Expk{ \norms{Gx^k + v^k}^2 }.
\end{array}
\hspace{-2ex}
\end{equation}
By the third line of  \eqref{eq:ub_SG_estimator} in Definition~\ref{de:ub_SG_estimator} and utilizing again \eqref{eq:L_smooth}, we have
\begin{equation*} 
\begin{array}{lcl}
\Delta_{k} & \leq &  (1 - \rho ) \Delta_{k-1}   + C L^2 \Expk{ \norms{x^{k} - x^{k-1}}^2 }  +  \hat{C} L^2  \Expk{ \norms{x^{k-1} - x^{k-2}}^2 }.
\end{array}
\end{equation*}
Rearranging this inequality, we get
\begin{equation*} 
\begin{array}{lcl}
\Delta_{k} & \leq &  \big( \frac{1-\rho}{\rho} \big) \big( \Delta_{k-1} - \Delta_{k} \big)  + \frac{\hat{C} L^2}{ \rho } \big[ \Expk{ \norms{x^{k-1} - x^{k-2}}^2  } -  \Expk{ \norms{x^k - x^{k-1}}^2 } \big] \\
&& + {~}  \frac{(C + \hat{C})L^2}{ \rho }  \Expk{ \norms{x^{k} - x^{k-1}}^2 }.
\end{array}
\end{equation*}
Substituting this inequality into \eqref{eq:VROG4NI_key_bound_proof5}, we can show that
\begin{equation*} 
\begin{array}{lcl}
\Expk{ \Lc_k } - \Expk{ \Lc_{k+1} } & \geq &  \frac{1}{2} \Big[ \mu -  2(1+\mu)\gamma(1-\gamma)L^2\eta^2 - \frac{2L^2\eta^2(1+\mu)(C + \hat{C})}{ \rho } \Big] \Expk{ \norms{x^k - x^{k-1}}^2 } \\
&& + {~} \gamma[ 1 - \gamma  - \mu(3\gamma-1)]\eta^2 \Expk{ \norms{Gx^{k-1} + v^k }^2 }  \\
&& + {~} (1-\gamma)\eta \big[  (1+\mu)(2\gamma - 1)\eta - 2\kappa \big]  \Expk{ \norms{Gx^k + v^k}^2 } \\
&& - {~} \frac{L^2\eta^2\hat{C}(1+\mu)}{ \rho } \Big[ \Expk{ \norms{x^{k-1} - x^{k-2}}^2 }  - \Expk{ \norms{x^k - x^{k-1}}^2 } \Big] \\
&& - {~} \frac{\eta^2(1+\mu)(1-\rho)}{\rho} \big( \Delta_{k-1} - \Delta_{k} \big).
\end{array}
\end{equation*}
Rearranging this inequality and using $\Ec_k$ from \eqref{eq:VROG4NE_Lfuncs}, we obtain \eqref{eq:VROG4NI_key_est2}.
\end{proof}
%%% End of Proof.

Now, we are ready to prove our second main result, Theorem~\ref{th:VROG4NI_convergence} in the main text.

%%% Proof of Theorem 4.1.
\begin{proof}[\textbf{Proof of Theorem~\ref{th:VROG4NI_convergence}}]
Since we fix $\gamma \in \big(\frac{1}{2}, 1\big)$ and $\mu := \frac{1-\gamma}{3\gamma-1}$, we have $\mu > 0$ and $1 + \mu = \frac{2\gamma}{3\gamma-1}$.
Let us denote by $M := 4\gamma^2 + \frac{ 4\gamma}{ 1-\gamma } \cdot \frac{C + \hat{C}}{\rho} $ as in Theorem~\ref{th:VROG4NI_convergence}.
Then, \eqref{eq:VROG4NI_key_est2} reduces to
\begin{equation}\label{eq:VROG4NI_key_est2_proof}
\begin{array}{lcl}
\Expk{ \Ec_k } - \Expk{ \Ec_{k+1} }  & \geq &  \frac{(1 - \gamma)(1 -  M \cdot L^2\eta^2 )}{2(3\gamma-1)} \Expk{ \norms{x^k - x^{k-1}}^2 } \\
&& + {~} 2(1-\gamma)\eta \big[  \frac{\gamma(2\gamma - 1)\eta}{3\gamma-1} - \kappa \big]  \Expk{ \norms{Gx^k + v^k}^2 }.
\end{array}
\end{equation}
Let us choose $\eta > 0$ such that $\frac{\gamma(2\gamma - 1)\eta}{3\gamma-1} - \kappa > 0$ and $1 -  M \cdot L^2\eta^2 \geq 0$.
These two conditions lead to $\frac{(3\gamma-1)\kappa}{\gamma(2\gamma-1)} < \eta \leq \frac{1}{L\sqrt{M}}$ as stated in Theorem~\ref{th:VROG4NI_convergence}. 
However, this condition holds if $L^2\kappa^2 < \frac{\gamma^2(2\gamma-1)^2}{M(3\gamma-1)^2}$.
This condition is equivalent to $L\kappa \leq \delta$ as our condition in Theorem~\ref{th:VROG4NI_convergence}, where $\delta := \frac{\gamma(2\gamma-1)}{(3\gamma-1)\sqrt{M}}$.

Averaging  \eqref{eq:VROG4NI_key_est2_proof} from $k=0$ to $K$ and noting that $\Expk{\Ec_k} \geq 0$ for all $k\geq 0$, we get
\begin{equation*} 
\begin{array}{lcl}
\frac{1}{K+1}\sum_{k=0}^{K} \Expk{ \norms{Gx^{k} + v^k}^2 }  & \leq &  \frac{ (3\gamma - 1) \cdot \mbb{E}[ \Ec_0 ] }{2(1-\gamma)[ \gamma(2\gamma-1)\eta - (3\gamma-1)\kappa] \eta (K + 1)}, \vspace{1ex} \\
\frac{(1 - M L^2\eta^2) }{K+1}\sum_{k=0}^{K} \Expk{ \norms{x^k - x^{k-1}}^2 } & \leq & \frac{ 2(3\gamma-1) \cdot \mbb{E}[ \Ec_0 ] }{(1-\gamma)(K + 1)}.
\end{array}
\end{equation*}
Finally, since $x^{-1} = x^{-2} = x^0$, we have $\Delta_{-1} = \Delta_0$.
However, since $\widetilde{S}^0_{\gamma}  =  (1-\gamma)Gx^0 = S_{\gamma}^0$, we get $\Delta_0 = \norms{\widetilde{S}^0_{\gamma} - S_{\gamma}^0}^2 = 0$. 
Using these relations, $\rho \in [0, 1]$ and $\gamma < 1$, we can show that
\begin{equation*}
\begin{array}{lcl}
\Expk{ \Ec_0 } & = & \Expk{ \norms{x^0 + \gamma \eta (Gx^{0} + v^0) - x^{\star}}^2 } +  \frac{ \eta^2(1+\mu)(1-\rho)}{\rho} \Delta_0 \\
& \leq &  2 \Expk{ \norms{x^0 - x^{\star}}^2 } + 2\gamma^2\eta^2 \Expk{ \norms{Gx^0 + v^0}^2}  +  \frac{2\gamma \eta^2}{(3\gamma-1)\rho} \Delta_0 \\
& = &  2 \Expk{ \norms{x^0 - x^{\star}}^2 } + 2\gamma^2\eta^2 \Expk{ \norms{Gx^0 + v^0}^2}.
\end{array}
\end{equation*}
Substituting this upper bound into the above two estimates, we get two lines of \eqref{eq:VROG4NI_convergence_bound1}.
\end{proof}
%%% End of Proof.

Finally, we prove Corollaries~\ref{co:VROG2_SVRG_convergence1_v2} and \ref{co:VROG2_SAGA_convergence1_v3} in the main text.
Unlike Corollaries \ref{co:SVRG_convergence1_v2} and \ref{co:SVRG_convergence1_v3} where we fix $\gamma := \frac{3}{4}$, here we state these corollaries for any value of $\gamma \in \left(\frac{1}{2}, 1\right)$.

%%% Proof of Corollary 4.2.
\begin{proof}[\textbf{Proof of Corollary~\ref{co:VROG2_SVRG_convergence1_v2}}]
For the SVRG estimator \eqref{eq:SVRG_estimator}, we have $\rho := \frac{\mbf{p}}{2} \in (0, 1]$, $C := \frac{4-6\mbf{p} + 3\mbf{p}^2 }{b\mbf{p}}$, $\hat{C} := \frac{2\gamma^2(2 - 3\mbf{p} + \mbf{p}^2)}{b\mbf{p}}$, and $\Delta_0 = 0$ due to \eqref{eq:Deltak_quantity_full} and $x^0 = x^{-1} = w^0$.
In this case, we have $\Lambda := \frac{C+\hat{C}}{\rho} = \frac{4(1+\gamma^2)(2 - 3\mbf{p}) + 2(3+2\gamma^2)\mbf{p}^2}{b\mbf{p}^2} \leq \frac{8(1+\gamma^2)}{b\mbf{p}^2}$, and thus $M$ in Theorem~\ref{th:VROG4NE_convergence} reduces to $M :=  4\gamma^2 + \frac{4\gamma}{1-\gamma}\Lambda \leq 4\gamma^2 + \frac{32(1+\gamma^2)}{b\mbf{p}^2}$.
%Hence, our bound \eqref{eq:svrg_bound_for_NI} is the same as  the first line of \eqref{eq:VROG4NI_convergence_bound1}.

Suppose that $b\mbf{p}^2 \leq 1$.
Since $\Lambda \leq \frac{ 8(1+\gamma^2) }{b\mbf{p}^2} $ and $M =  4\gamma^2 + \frac{4\gamma}{1-\gamma}\Lambda \leq 4\gamma^2 + \frac{32\gamma(1+\gamma^2)}{(1-\gamma) b\mbf{p}^2} \leq \frac{4\gamma(8 + \gamma + 7\gamma^2)}{(1-\gamma)b\mbf{p}^2}$.
If we choose $\eta := \frac{1}{L\sqrt{M}}$, then we have $\eta \geq \frac{\sqrt{1-\gamma}\sqrt{b}\mbf{p}}{2L\sqrt{\gamma(8 + \gamma + 7\gamma^2)}} = \frac{\sigma \sqrt{b}\mbf{p}}{L}$ with $\sigma := \frac{\sqrt{1-\gamma}}{2\sqrt{8 + \gamma + 7\gamma^2}}$, then it satisfies $\frac{(3\gamma-1)\kappa}{\gamma(2\gamma-1)} < \eta \leq \frac{1}{L\sqrt{M}}$ in Theorem~\ref{th:VROG4NI_convergence}, provided that $L\kappa \leq \delta$.
Note that using $\eta \geq \frac{\sigma \sqrt{b}\mbf{p}}{ L }$ in \eqref{eq:VROG4NI_convergence_bound1} of Theorem~\ref{th:VROG4NI_convergence} we obtain the bound  \eqref{eq:svrg_bound_for_NI}.

Now, from the first line of \eqref{eq:VROG4NI_convergence_bound1}, to guarantee $\frac{1}{K+1}\sum_{k=0}^K\Exp{\norms{Gx^k + v^k}^2} \leq \epsilon^2$, we need to impose $\frac{\Theta \hat{R}_0^2}{\eta^2(K+1)} \leq \epsilon^2$, where $\hat{R}_0^2 := \norms{x^0 - x^{\star}}^2 + \gamma^2\eta^2\norms{Gx^0 + v^0}^2$.
Since $\eta \geq \frac{\sigma \sqrt{b}\mbf{p}}{ L }$, the last condition holds if we choose $K :=  \left\lfloor \Gamma \cdot \frac{L^2 \hat{R}_0^2}{ b\mbf{p}^2\epsilon^2} \right\rfloor$, where $\Gamma := \frac{\Theta}{\sigma^2}$.


Finally, at each iteration $k$, \eqref{eq:VROG4NI} requires $3$ mini-batches of size $b$, and occasionally compute the full $Gw^k$, leading to the cost of $n\mbf{p} + 3b$ per iteration.
Thus the total complexity is 
\begin{equation*}
\begin{array}{lcl}
\Tc_{c} & := & K(n\mbf{p} + 3b) = \frac{ \Gamma L^2 \hat{R}_0^2(n\mbf{p} + 3b)}{ b\mbf{p}^2\epsilon^2} = \frac{\Gamma L^2 \hat{R}_0^2}{\epsilon^2}\big(\frac{n}{b\mbf{p}} + \frac{3}{ \mbf{p}^2} \big).
\end{array}
\end{equation*}
If we choose $b := \lfloor n^{2/3} \rfloor$ and $\mbf{p} := n^{-1/3}$, then $b\mbf{p}^2 = 1$ and $\Tc_{c} = \frac{4\Gamma n^{2/3} L^2 \hat{R}_0^2}{\epsilon^2}$.
For the SVRG estimator \eqref{eq:SVRG_estimator}, one needs to compute $Gw^0$, which requires $n$ evaluations of $G_i$.
Hence, the total  evaluations of $G_i$ is $\Tc_{G_i} = n + \big\lfloor \frac{4\Gamma n^{2/3} L^2 \hat{R}_0^2}{\epsilon^2} \big\rfloor$.
Moreover, at each iteration, we need one evaluation of $J_{\gamma\eta T}$.
Therefore, the total evaluations of $J_{\gamma\eta T}$ is $\Tc_{T} := K = \big\lfloor \Gamma \cdot \frac{L^2 \hat{R}_0^2}{ b\mbf{p}^2\epsilon^2} \big\rfloor = \big\lfloor \Gamma \cdot \frac{L^2 \hat{R}_0^2}{\epsilon^2} \big\rfloor$.
\end{proof}
%%% End of Proof.

%%% Proof of Corollary 3.3.
\begin{proof}[\textbf{Proof of Corollary~\ref{co:VROG2_SAGA_convergence1_v3}}]
Since we use the SAGA estimator \eqref{eq:SAGA_estimator}, we have  $\rho := \frac{b}{2n} \in (0, 1]$, $C :=  \frac{[ 2(n - b)(2n+b) + b^2] }{ n b^2}$, and $\hat{C} := \frac{2(n - b)(2n+b) \gamma^2}{ nb^2}$.
In this case, since $b \geq 1$, we get $\Lambda := \frac{C + \hat{C}}{\rho} = \frac{2}{b} + \frac{4(1+\gamma^2)(n-b)(2n+b)}{b^3} \leq 2 + \frac{8(1+\gamma^2)n^2}{b^3}$.
Hence, $M$ in Theorem~\ref{th:VROG4NE_convergence} reduces to 
\begin{equation*}
\begin{array}{lcl}
M := 4\gamma^2 + \frac{4\gamma}{1-\gamma} \cdot \Lambda \leq  \frac{4\gamma(2 + \gamma - \gamma^2)}{1-\gamma} +  \frac{32\gamma(1+\gamma^2)n^2}{(1- \gamma) b^3} 
\end{array}
\end{equation*}
Suppose that $1 \leq b \leq n^{2/3}$.
Then, we can show that $M \leq \big[  \frac{4\gamma(2 + \gamma - \gamma^2)}{1-\gamma} +  \frac{32\gamma(1+\gamma^2)}{1- \gamma}  \big] \frac{n^2}{b^3} = \frac{4\gamma(10 + \gamma + 7\gamma^2)}{(1-\gamma)b^3} =  \frac{n^2}{ \sigma^2b^3}$, where $\sigma :=  \frac{\sqrt{1-\gamma}}{2\sqrt{\gamma(10 + \gamma + 7\gamma^2)}}$.
Hence, if we choose $\eta := \frac{1}{L\sqrt{M}}$, then we get $\eta \geq \frac{\sigma b^{3/2}}{nL}$.
Note that using $\eta \geq \frac{\sigma b^{3/2}}{ nL }$ in \eqref{eq:VROG4NI_convergence_bound1} of Theorem~\ref{th:VROG4NI_convergence} we obtain the bound  \eqref{eq:saga_bound_for_NI}.

For $\eta := \frac{1}{L\sqrt{M}} \geq \frac{\sigma b^{3/2}}{nL}$, from the first line of \eqref{eq:VROG4NI_convergence_bound1}, to guarantee $\frac{1}{K+1}\sum_{k=0}^K\Exp{\norms{Gx^k + v^k}^2} \leq \epsilon^2$, we need to impose $\frac{\Theta \hat{R}_0^2}{\eta^2(K+1)}  \leq \epsilon^2$, where $\hat{R}_0^2 := \norms{x^0 - x^{\star}}^2 + \gamma^2\eta^2\norms{Gx^0 + v^0}^2$.
Since $\eta \geq  \frac{\sigma b^{3/2} }{nL}$, the last condition holds if we choose $K :=  \big\lfloor \Gamma \cdot \frac{L^2 \hat{R}_0^2 n^2}{ b^3 \epsilon^2} \big\rfloor$, where $\Gamma := \frac{\Theta}{\sigma^2}$.

Finally,  at each iteration $k$, \eqref{eq:VROG4NI} requires $3$ mini-batches of size $b$, leading to the cost of $3b$ per iteration.
Thus the total complexity is 
\begin{equation*}
\begin{array}{lcl}
\Tc_{c} & := &  3bK = \big\lfloor \frac{ 3\Gamma L^2 \hat{R}_0^2 n^2 }{ b^2 \epsilon^2} \big\rfloor.
\end{array}
\end{equation*}
If we choose $b := \lfloor n^{2/3} \rfloor$, then $\Tc_{c} = \big\lfloor \frac{3\Gamma  L^2 \hat{R}_0^2 n^{2/3} }{\epsilon^2} \big\rfloor$.
For the SAGA estimator \eqref{eq:SAGA_estimator}, one needs to compute $Gw^0$, which requires $n$ evaluations of $G_i$.
We conclude that \eqref{eq:VROG4NI} requires $\Tc_{G_i} := n + \big\lfloor \frac{3\Gamma L^2 \hat{R}_0^2  n^{2/3}  }{\epsilon^2} \big\rfloor$ evaluations of $G_i$.
Moreover, since each iteration, it requires one evaluation of $J_{\gamma\eta T}$, we need $\Tc_T := K =  \big\lfloor \Gamma \cdot \frac{L^2 \hat{R}_0^2}{ \epsilon^2} \big\rfloor$ evaluations of $J_{\gamma\eta T}$.
\end{proof}
%%% End of Proof.

\begin{remark}\label{re:lower_bounds_of_eta}
For the SVRG estimator, if we choose $\gamma = \frac{3}{4}$, then we have $\sigma := 0.0702$.
% and $\Gamma := \frac{2030\eta}{0.75\eta - 2.5\kappa}$.
%Moreover, we also have $\Theta := \frac{10\eta}{0.75\eta - 2.5\kappa}$.
Hence, we have $\eta \geq \frac{0.0702 \sqrt{b}\mbf{p}}{L}$.
However, if we choose $\gamma := 0.55$, then $\eta \geq \frac{0.1027\sqrt{b}\mbf{p}}{L}$.
If we choose $b = \lfloor n^{2/3}\rfloor$ and $\mbf{p} = n^{-1/3}$, then the latter lower bound becomes $\eta \geq \frac{0.1027}{L}$.

For the SAGA estimator, if we choose $\gamma = \frac{3}{4}$, then we have $\sigma := 0.0753$.
Hence, we get $\eta \geq \frac{0.0753 b^{3/2}}{nL}$.
However, if we set $\gamma := 0.55$, then $\eta \geq \frac{0.1271 b^{3/2}}{nL}$.
If we choose $b = \lfloor n^{2/3}\rfloor$, then the latter lower bound becomes $\eta \geq \frac{0.1271}{L}$.
% and $\Gamma := \frac{1762.5\eta}{0.75\eta - 2.5\kappa}$ as stated.
%Moreover, we also have $\Theta := \frac{10\eta}{0.75\eta - 2.5\kappa}$.
%Using this $\Theta$ in \eqref{eq:VROG4NI_convergence_bound1} of Theorem~\ref{th:VROG4NI_convergence} to complete our proof.

Note that these lower bounds of $\eta$ can be further improved by refining the related parameters in Lemma~\ref{le:VROG2_key_bounds}, and carefully choosing $\mu$ in the proof of Theorem~\ref{th:VROG4NI_convergence}.
\end{remark}

%%%%%%%%%%%%%%%%%%%%%%%%%%%%%%%%%%%%%%%%%%%%%%%%%%%%%%%%%%%%
%%% D. Details of Experiments and Additional Experiments
%%%%%%%%%%%%%%%%%%%%%%%%%%%%%%%%%%%%%%%%%%%%%%%%%%%%%%%%%%%%
\beforesec
\section{Details of Experiments and Additional Experiments}\label{apdx:sec:num_experiments}
\aftersec
Due to space limit, we do not provide the details of experiments in Section~\ref{sec:NumExps}.
In this Supp. Doc., we provide the details of our implementation and experiments.
We also add more examples to illustrate our algorithms and compare them with existing methods.
All algorithms are implemented in Python, and all the experiments are run on a MacBookPro. 2.8GHz Quad-Core Intel Core I7, 16Gb Memory.

\beforesubsec
\subsection{Synthetic WGAN Example}\label{apdx:subsec:detail_numexps}
\aftersubsec
We modify the synthetic example in \citep{daskalakis2018training} built up on WGAN from \citep{arjovsky2017wasserstein} as our first example. 
Suppose that the generator is a simple additive model $G_{\theta}(z) =  \theta + z$ with the noise input $z$ generated from a normal distribution $\Nc(0, \Id)$, and the discriminator is also a linear function $D_{\beta}(w) = \iprods{K\beta, w}$ for a given matrix $K$, where $\theta \in \R^{p_1}$ and $\beta \in \R^{p_2}$, and $K \in \R^{p_1\times p_2}$ is a given matrix.
The goal of the generator is to find a true distribution $\theta = \theta^{*}$, leading to the following loss:
\begin{equation*}
\Lc(\theta, \beta) := \mathbb{E}_{u \sim \Nc(\theta^{*}, \Id)}\big[ \iprods{K\beta, w} \big] - \mathbb{E}_{z \sim \Nc(0, \Id)}\big[ \iprods{K\beta, \theta + z} \big].
\end{equation*}
Suppose that we have $n$ samples for both $w$ and $z$ leading to the following bilinear minimax problem:
\begin{equation}\label{eq:syn_WGAN_sup}
{\inf_{\theta \in\R^{p_1}} \sup_{\beta\in \R^{p_2}}} \big\{ \Lc(\theta, \beta) := f(\theta) + \frac{1}{n}\sum_{i=1}^n \big[ \iprods{K\beta, w_i - z_i - \theta} \big] - g(\beta) \big\}.
\end{equation}
Here, we add two convex functions $f(\theta)$ and $g(\beta)$ to possibly handle constraints or regularizers associated with $\theta$ and $\beta$, respectively.

If we define $x := [\theta,\beta] \in \R^{p_1+p_2}$, $Gx = [\nabla_{\theta}\Lc(\theta,\beta), -\nabla_{\beta}\Lc(\theta,\beta)] := -[K\beta, \frac{1}{n}\sum_{i=1}^nK^{\top}(w_i - z_i - \theta)]$, and $T := [\partial{f}(\theta), \partial{g}(\beta)]$, then the optimality condition of this minimax problem becomes $0 \in Gx + Tx$, which is a special case of \eqref{eq:NI} with $Gx$ being linear.
The model \eqref{eq:syn_WGAN_sup} is different from the one in \citep{daskalakis2018training} at two points:
\begin{compactitem}
\item It involves a linear operator $K$, making it more general than \citep{daskalakis2018training}.
\item It has two additional terms $f$ and $g$, making it broader to also cover constraints or non-smooth regularizers. 
\end{compactitem}
In our experiments below, we consider two cases:
\begin{compactitem}
\item \textbf{Case 1 (Unconstrained setting).} We assume that $\theta \in \R^{p_1}$ and $\beta \in \R^{p_2}$.
\item \textbf{Case 2 (Constrained setting).} Assume that $\theta$ and $\beta$ stays in an $\ell_{\infty}$-ball of radius $r > 0$, leading to $f(\theta) := \delta_{[-r,r]^{p_1}}(\theta)$ and $g(\beta) := \delta_{[-r, r]^{p_2}}(\beta)$, the indicator of the $\ell_{\infty}$-balls.
\end{compactitem}

\beforesubsubsec
\subsubsection{The unconstrained case}\label{apdx:subsec:synWGAN_unconstr_case}
\aftersubsubsec
\textbf{$\mathrm{(a)}$~Algorithms.}
We implement three variants of \eqref{eq:VROG4NE} to solve \eqref{eq:syn_WGAN_sup}.
\begin{compactitem}
\item The first variant is using a double-loop SVRG strategy (called \texttt{VFR-svrg}), where the full operator $Gw^s$ at a snapshot point $w^s$ is computed at the beginning of each epoch $s$.
Then, we perform $\lfloor n/b \rfloor$ iterations $k$ to update $x^k$ using \eqref{eq:VROG4NE}, where $b$ is the mini-batch size.
Finally, we set the next snapshot point $w^{s+1} := x^{k+1}$ after finishing the inner loop.
\item The second variant is called a loopless one, \texttt{LVFR-svrg}, where we implement exactly the same scheme \eqref{eq:VROG4NE} as in this paper and using the Loopless-SVRG estimator.
\item The third variant is \texttt{VFR-saga}, where we use the SAGA estimator in \eqref{eq:VROG4NE}.
\end{compactitem}
We also compare our methods with the deterministic optimistic gradient (\texttt{OG}) in \citep{daskalakis2018training}, the variance-reduced FRBS (\texttt{VFRBS}) in \citep{alacaoglu2022beyond}, and the variance-reduced extragradient (\texttt{VEG}) in \citep{alacaoglu2021stochastic}.

\textbf{$\mathrm{(b)}$~Input data.}
For \eqref{eq:NE}, we generate a vector $\theta^{*}$ from the standard normal distribution as our true mean in $\R^{p_1}$.
Then, we generate i.i.d. samples $w_i$ and $z_i$ from normal distribution $\Nc(\theta^{*}, \Id)$ and $\Nc(0, \Id)$, respectively for $i=1,2,\cdots, n$ in $\R^{p_1}$ and $\R^{p_2}$, respectively.
We perform two expertiments: \textbf{Experiment 1} with $n=5000$ and $p_1 = p_2 = 100$, and \textbf{Experiment 2} with $n=10000$ and $p_1 = p_2 = 200$.
For each experiment, we run 10 times up to 100 epochs, corresponding to $10$ problem instances, using the same setting, but different input data $(w_i, z_i)$, and then compute the mean of the relative operator norm $\norms{Gx^k} / \norms{Gx^0}$.
This mean is then plotted.

\textbf{$\mathrm{(c)}$~Parameters.}
For the optimistic gradient algorithm (\texttt{OG}), we choose its learning rate $\eta := \frac{1}{L}$, where $L$ is the Lipschitz constant of $G$, though its theoretical learning rate is much smaller.
For our methods in \eqref{eq:VROG4NE}, if $n=5000$, and we choose $b := \lfloor 0.5 n^{2/3} \rfloor = 146$, and the probability $\mathbf{p} := \frac{2}{n^{1/3}} = 0.1170$, then $\eta := \frac{1}{L\sqrt{M}} =  \frac{0.1905}{ L}$.
However, due to the under estimation of $M$, we instead use a larger learning rate $\eta := \frac{1}{2L}$ for all three variants, and choose a mini-batch of size $b := \lfloor 0.5 n^{2/3} \rfloor$, and a probability $\mathbf{p} := \frac{1}{n^{1/3}}$ for the loopless SVRG variant.

For the forward-reflected-backward splitting method with variance reduction (\texttt{VFRBS}) in \citep{alacaoglu2022beyond}, we choose its learning rate $\eta := \frac{0.95(1 - \sqrt{1 - \mathbf{p}})}{2L}$ as suggested by \citep{alacaoglu2022beyond}.
However, we still choose the probability $\mbf{p} = \frac{1}{n^{1/3}}$ and the mini-batch size $b = \lfloor  0.5 n^{2/3}\rfloor$ as our methods.
These values are much larger the ones suggested in \citep{alacaoglu2022beyond}, typically $\mbf{p} = \BigOs{1/n}$.

For the variance reduction extragradient method (\texttt{VEG}) in \citep{alacaoglu2021stochastic}, we choose its learning rate $\eta := \frac{0.95\sqrt{1 - \alpha}}{L}$ for $\alpha := 1 - \mathbf{p}$ from the paper.
However, again, we also choose $\mathbf{p} :=  \frac{1}{n^{1/3}}$ and  $b = \lfloor  0.5 n^{2/3}\rfloor$ in this method, which is the same as ours, though their theoretical results suggest smaller values of $\mathbf{p}$ (e.g., $\mathbf{p} = \frac{1}{n}$).
Note that if $n = 5000$, then the batch size $b := 150$ and the probability $\mathbf{p} := 0.062$, but if $n=10000$, then $b = 239$ and $\mathbf{p} = 0.0479$. 

\textbf{$\mathrm{(d)}$~Experiments for $K =\Id$.}
We perform two experiments: \textbf{Experiment 1} with $(n, p) = (5000, 200)$ and \textbf{Experiment 2} with $(n, p) = (10000, 400)$ as discussed above.
We run each experiment with $10$ problem instances and compute the mean of the relative residual norm $\norms{Gx^k}/\norms{Gx^0}$.
The results of this test are plotted in Figure~\ref{fig:wgan_exam1}.

\begin{figure}[ht]
\vspace{-1ex}
\centering
\includegraphics[width=0.9\textwidth]{figures/unconstr_wgan_exam1}
\vspace{-1.5ex}
\caption{
Performance of $6$ algorithms to solve \eqref{eq:syn_WGAN_sup} on $2$ experiments when $K = \Id$.
}
\label{fig:wgan_exam1}
\vspace{-1ex}
\end{figure}

For these particular experiments, our methods highly outperform \texttt{OG}, \texttt{VFRBS}, and \texttt{VEG}.
It shows that \texttt{VFR-svrg} is the best overall, while \texttt{LVFR-saga} and \texttt{VFR-svrg} have a similar performance in both experiments.
Both the competitors: \texttt{VFRBS} and \texttt{VEG} do not perform well in this test and they are much slower than ours and also \texttt{OG}.
This is perhaps due to a small learning rate of \texttt{VFRBS} although we choose the same mini-batch size $b$ and the same probability $\mbf{p}$ as ours.

\textbf{$\mathrm{(e)}$~Experiments for $K\neq \Id$.}
Now, we test these 6 algorithms for the case $K\neq \Id$ in our extended model \eqref{eq:syn_WGAN_sup}, where $K$ is generated randomly from the standard normal distribution.
Then, we normalize $K$ as $K/\norms{K}$ to get a unit Lipschitz constant $L=1$.

Again, we use the same configuration as in Figure~\ref{fig:wgan_exam1} and also run our experiments on $10$ problems and report the mean results.
We perform two experiments: \textbf{Experiment 1} with $n=5000$ and $p_1=p_2 = p =100$, and \textbf{Experiment 2} with $n=10000$ and $p_1 = p_2 = p=200$.
The results are reported in Figure~\ref{fig:wgan_exam_1k}.

\begin{figure}[ht]
\vspace{-0ex}
\centering
\includegraphics[width=0.9\textwidth]{figures/unconstr_wgan_exam2}
\vspace{-1ex}
\caption{
Performance of 6 algorithms to solve \eqref{eq:syn_WGAN_sup} on $2$ experiments when  $K \neq \Id$.
}
\label{fig:wgan_exam_1k}
\vspace{-0ex}
\end{figure}

We still observe that our algorithms work well and outperform their competitors.
However, after 100 epochs, these methods can only reach a $10^{-2}$ accuracy level for an approximate solution.

%%%%%%%%%%%%%%%%%%%%%%%%%%%%%%%%%
%%% E.1.3. The unconstrained case -- Varying b and p.
%%%%%%%%%%%%%%%%%%%%%%%%%%%%%%%%%
\beforesubsubsec
\subsubsection{The unconstrained case -- Varying $b$ and $\mbf{p}$}\label{apdx:subsubsec:other_experiment_of_WGAN}
\aftersubsubsec
We can certainly tune the parameters to make our competitors (\texttt{VFRBS}) and (\texttt{VEG}) work better.
However, such parameter configurations are far from satisfying the conditions of their theoretical results.
For example, if we set $\mathbf{p} = \frac{20}{\sqrt{n}}$, then both \texttt{VFRBS} and \texttt{VEG} work better.
In particular, if $n=5000$, then we get $\mathbf{p} = \frac{20}{\sqrt{n}} = 0.28$, which is several times larger than its suggested value $\mbf{p} =\frac{1}{n} = 2\times 10^{-4}$.

Let us further experiment  other choices of parameters (i.e. the mini-batch size $b$ and the probability  $\mathbf{p}$ of flipping a coin) to observe the performance of these algorithms.

\textbf{$\mathrm{(a)}$~Larger $b$.}
Figure~\ref{fig:wgan_ne_mb10_p033} reveals the performance of these algorithms when we increase the mini-batch size $b$ to a larger value $b = \lfloor0.1n\rfloor$, while keeping the probability $\mathbf{p} = \frac{1}{n^{1/3}}$ unchanged.

\begin{figure}[ht]
%\vspace{-1ex}
\centering
\includegraphics[width=0.9\textwidth]{figures/unconstr_wgan_exam1_p13_b01n}
\vspace{-1ex}
\caption{
Performance of 6 algorithms for a large $b = \lfloor0.1n\rfloor$ and a unchanged $\mathbf{p} = \frac{1}{n^{1/3}}$.
}
\label{fig:wgan_ne_mb10_p033}
\vspace{-1.5ex}
\end{figure}

Note that for $n=5000$, we have $b = 500$ and $\mathbf{p} = 0.058$, and for $n=10000$, we have $b = 1000$ and $\mathbf{p} = 0.046$.
With these large mini-batches, our algorithms still outperform other methods, while \texttt{VFRBS} and \texttt{VEG} are significantly slowed down. 
The double-loop variant of \eqref{eq:VROG4NE} with SVRG performs best, while \texttt{LVFR-svrg} and \texttt{VFR-saga} have a similar performance.

\textbf{$\mathrm{(b)}$~Medium $b$ and larger $\mbf{p}$.}
Next, we set $b$ to a medium size of $b = \lfloor 0.05n \rfloor$ (corresponding to $b=250$ for $n=5000$ and $b=500$ for $n=10000$) and increase $\mathbf{p} = \frac{1}{n^{1/4}}$ (corresponding to $\mathbf{p} = 0.119$ for $n=5000$ and $\mathbf{p} = 0.1$ for $n=10000$).
Then, the results are shown in Figure~\ref{fig:wgan_ne_mb05_p025}.

\begin{figure}[hpt!]
\vspace{-2ex}
\centering
\includegraphics[width=0.9\textwidth]{figures/unconstr_wgan_exam1_p14_b005n}
\vspace{-1.5ex}
\caption{
Performance of 6 algorithms for a medium $b = \lfloor0.05n\rfloor$ and larger $\mathbf{p} = \frac{1}{n^{1/4}}$.
}
\label{fig:wgan_ne_mb05_p025}
\vspace{-2ex}
\end{figure}

Then, we observe that \texttt{LVFR-svrg} and \texttt{VFR-saga} superiorly  outperform the others.
The performance of the double-loop \texttt{VFR-svrg} is still similar to the previous tests since it is not affected by $\mbf{p}$.
In addition, \texttt{VEG} is now comparable with \texttt{OG}, but \texttt{VFRBS} remains the slowest one.


\textbf{$\mathrm{(c)}$~Large $b$ and small $\mbf{p}$.}
To see the effect of $\mathbf{p}$ on our competitors: \texttt{VFRBS} and \texttt{VEG}, as suggested by their theory, we decrease $\mathbf{p}$ to $\mathbf{p} = \frac{1}{n^{1/2}}$ (corresponding to $\mathbf{p} = 0.014$ for $n=5000$ and $\mathbf{p} = 0.01$ for $n=10000$) and still set $b = \lfloor 0.1n \rfloor$, and the results are plotted in Figure~\ref{fig:wgan_ne_mb10_p05}.

\begin{figure}[hpt!]
\vspace{-2ex}
\centering
\includegraphics[width=0.9\textwidth]{figures/unconstr_wgan_exam1_p12_b01n}
\vspace{-1.5ex}
\caption{
Performance of 6 algorithms for a large $b = \lfloor 0.1n \rfloor$ and a small $\mathbf{p} = \frac{1}{n^{1/2}}$.
}
\label{fig:wgan_ne_mb10_p05}
\vspace{-2ex}
\end{figure}

As we can observed from Figure~\ref{fig:wgan_ne_mb10_p05}, our methods highly outperform \texttt{VFRBS} and \texttt{VEG}, suggesting that these competitors require a larger probability to select the snap-shot point $w^k$ for full-batch evaluation.
This is certainly not suggested in their theoretical results.

%%%% (d)~ Inclusions.
\beforesubsubsec
\subsubsection{The constrained case}\label{apdx:subsubsec:NI}
\aftersubsubsec
Next, we choose $f(\theta) = \delta_{[-r, r]^{p_1}}(\theta)$ and $g(\beta) := \delta_{[-r, r]^{p_2}}(\beta)$ as the indicators of the $\ell_{\infty}$-balls of radius $r = 5$, respectively.
In this case, we implement three variants of  \eqref{eq:VROG4NI}: the double-loop (\texttt{VFR-svrg}), the loopless  (\texttt{LVFR-svrg}), and the SAGA (\texttt{VFR-saga}) variants to solve \eqref{eq:NI} and compare against $3$ algorithms as in the unconstrained case.
Using the same data generating procedure as in the unconstrained case, we obtain the results as shown in Figure~\ref{fig:constr_wgan_exam1} when $K = \Id$.

\begin{figure}[ht]
\vspace{-0ex}
\centering
\includegraphics[width=0.9\textwidth]{figures/constr_wgan_exam1}
\vspace{-1.5ex}
\caption{
Comparison of $6$ algorithms to solve constrained instances of \eqref{eq:syn_WGAN_sup} on $2$ experiments when $K = \Id$  (The average of $10$ runs).
}
\label{fig:constr_wgan_exam1}
\vspace{-2ex}
\end{figure}
We see that the two SVRG variants of our \eqref{eq:VROG4NI}:  \texttt{VFR-svrg} and \texttt{LVFR-svrg},  as well as our \texttt{VFR-saga} variant remain working well compared to other methods.
They  superiorly outperform the  three competitors. 

Finally, we test our methods and their competitors for the case $K \neq \Id$ as we done in Figure~\ref{fig:wgan_exam_1k}.
Our results are plotted in Figure~\ref{fig:constr_wgan_exam2}, where we observe a similar behavior as in Figure~\ref{fig:wgan_exam_1k}.

\begin{figure}[ht]
\vspace{-0ex}
\centering
\includegraphics[width=0.9\textwidth]{figures/constr_wgan_exam2}
\vspace{-1.5ex}
\caption{
Comparison of $6$ algorithms to solve constrained instances of \eqref{eq:syn_WGAN_sup} on $2$ experiments when $K \neq \Id$  (The average of $10$ runs).
}
\label{fig:constr_wgan_exam2}
\vspace{-2ex}
\end{figure}

%%%% +++++++++++++++++++++++++++++++++++++++++++++++++++++++++++++++++
%% E.2. Details of Example 2 in Section 5.
\beforesubsec
\subsection{Nonconvex-Nonconcave Quadratic Minimax Problems}\label{apdx:subsec:bilinear_opt}
\aftersubsec
We recall the nonconvex-nonconcave quadratic minimax optimization problem  \eqref{eq:minimax_exam1} in this subsection:
\begin{equation}\label{eq:minimax_exam2_supp}
%\begin{array}{lcl}
\min_{u \in \R^{p_1}}\max_{v\in\R^{p_2}} \Big\{ \Lc(u, v) := \varphi(u) + \frac{1}{n}\sum_{i=1}^n\big[u^TA_iu +  u^TL_iv - v^TB_iv  + b_i^{\top}u - c_i^{\top}v\big] - \psi(v) \Big\},
%\end{array}
\end{equation}
where $A_i \in \R^{p_1\times p_1}$ and $B_i \in \R^{p_2\times p_2}$ are symmetric matrices, $L_i \in \R^{p_1\times p_2}$, $b_i \in \R^{p_1}$, $c_i \in \R^{p_2}$, and $\varphi = \delta_{\Delta_{p_1}}$ and $\psi =\delta_{\Delta_{p_2}}$ are the indicator of standard simplexes in $\R^{p_1}$ and $\R^{p_2}$, respectively.

Let us first define $x := [u, v] \in \R^p$ as the concatenation of the primal and dual variables $u$ and $v$, where $p := p_1+p_2$.
Next, we define
\begin{equation*}
G_ix = \mathbf{G}_ix + \mbf{g}_i := \begin{bmatrix}A_i & L_i \\ -L_i & B_i\end{bmatrix}\begin{bmatrix}u \\ v\end{bmatrix} + \begin{bmatrix}b_i \\ c_i\end{bmatrix} = \begin{bmatrix}A_iu + L_iv + b_i \\ -L_iu + B_iv + c_i\end{bmatrix}, \quad\text{and} \quad 
T := \begin{bmatrix}\partial{\varphi} \\ \partial{\psi}\end{bmatrix}.
\end{equation*}
Then, we denote $\mathbf{G}_i := \begin{bmatrix}A_i & L_i \\ -L_i & B_i\end{bmatrix}$, and $\mathbf{g}_i :=  \begin{bmatrix}b_i \\ c_i\end{bmatrix}$.
Clearly, $G_i(\cdot)$ is an affine mapping from $\R^p$ to $\R^p$, but $\mathbf{G}_i$ is nonsymmetric.
Let $Gx := \frac{1}{n}\sum_{i=1}^nG_ix = \big(\frac{1}{n}\sum_{i=1}^n\mathbf{G}_i\big)x + \frac{1}{n}\sum_{i=1}^n\mathbf{g}_i = \mathbf{G}x + \mathbf{g}$, where $\mathbf{G} := \frac{1}{n}\sum_{i=1}^n\mathbf{G}_i$ and $\mathbf{g} := \frac{1}{n}\sum_{i=1}^n\mathbf{g}_i$.
Then, the optimality condition of \eqref{eq:minimax_exam2_supp} becomes $0 \in Gx + Tx$, which is exactly in the form \eqref{eq:NI}.
Clearly, if $A_i$ and/or $B_i$ are not positive semidefinite, then \eqref{eq:minimax_exam2_supp} possibly covers nonconvex-nonconcave minimax optimization instances. 

\beforesubsubsec
\subsubsection{The unconstrained case}\label{subsec:exam2_unconstrained}
\aftersubsubsec
We consider the case $\varphi = 0$ and $\psi = 0$, leading to an unconstrained setting of \eqref{eq:minimax_exam2_supp}, i.e. $T = 0$ as considered in  \eqref{eq:minimax_exam1} of the main text.
Hence, the optimality condition of \eqref{eq:minimax_exam2_supp} reduces to $Gx = 0$, which is of the form \eqref{eq:NE}.

\textbf{$\mathrm{(a)}$~How to generate data?}
To run our experiments, we generate synthetic data as follows.
First, we fix the dimensions $p_1$ and $p_2$ and the number of components $n$.
We generate $A_i = Q_iD_iQ_i^T$ for a given orthonormal matrix $Q_i$ and a diagonal matrix $D_i = \mathrm{diag}(D_{i}^1, \cdots, D_{i}^{p_1})$, where its elements are generated from standard normal distribution and clipped its negative entries as $\max\sets{D_{i}^j, \varepsilon}$ for $j=1,\cdots, p_1$ and $\varepsilon := -0.1$.
This choice of $A_i$ guarantees that $A_i$ is symmetric, but possibly not positive semidefinite.
The matrix $B_i$ is also generated by the same way.
The pay-off matrix $L_i$ is an $p_1\times p_2$ matrix, which is also generated from the standard normal distribution  for all $i \in [n]$.
The vectors $b_i$ and $c_i$ are generated from the standard normal distribution for $i\in [n]$.
With this data generating procedure, $\mathbf{G}_i$ is not symmetric and possibly not positive semidefinite.

%%% (b) Algorithms
\textbf{$\mathrm{(b)}$~Algorithms.}
We again test 6 algorithms: two variants (double-loop SVRG -- \texttt{VFR-svrg}) and (loopless SVRG -- \texttt{LVFR-svrg}) of \eqref{eq:VROG4NE}, our \eqref{eq:VROG4NE} with SAGA estimator (\texttt{VFR-saga}), \texttt{VFRBS} from \citep{alacaoglu2022beyond}, \texttt{VEG} from \citep{alacaoglu2021stochastic}, and \texttt{OG} (the standard optimistic gradient method), e.g., from \citep{daskalakis2018training}.

%%% (c) Details of Example 2 in Section 5.
\textbf{$\mathrm{(c)}$~The details of Subsection~\ref{subsec:example1} in Section~\ref{sec:NumExps}.}
First, we provide the details of \textbf{Subsection}~\ref{subsec:example1} in Section~\ref{sec:NumExps}.
The purpose of this example is to verify our theoretical results stated in Corollaries~\ref{co:SVRG_convergence1_v2} and \ref{co:SVRG_convergence1_v3}.

For the SVRG estimator, let us first choose $\gamma := 0.75$, $b := \lfloor n^{2/3}\rfloor$, and $\mbf{p} := \frac{1}{n^{1/3}}$ as suggested by Corollary~\ref{co:SVRG_convergence1_v2}.
Then, we can directly compute $\eta := \frac{1}{L\sqrt{M}}$, where  $\Lambda := \frac{6.25(2-3\mbf{p}) + 4.125 \mbf{p}^2}{b\mbf{p}^2}$ and $M = 2.375 + \frac{11}{3}\Lambda$.
Clearly, if $n = 5000$, then  $\eta =  \frac{0.146153}{L}$.
Alternatively, if $n= 10000$, then $\eta = \frac{0.148934}{L}$.
These learning rates are used in our experiments plotted in Figure~\ref{fig:unconstr_minimax_exam1}.

Similarly, for the SAGA estimator, we also choose $\gamma := 0.75$ and $b := \lfloor n^{2/3}\rfloor$.
In this case, by Corollary~\ref{co:SVRG_convergence1_v3}, we can also directly compute $\eta := \frac{1}{L\sqrt{M}}$.
If $n = 5000$, then $\eta = \frac{0.146153}{L}$.
Alternatively, if $n = 10000$, then $\eta = \frac{0.145693}{L}$.
These learning rates are used in \texttt{VFR-saga}.

Note that since the theoretical value of $\mbf{p}$ in \texttt{VFRBS} and \texttt{VEG} is too small, we instead choose  $\mbf{p} := \frac{1}{n^{1/3}}$ and also $b := \lfloor n^{2/3}\rfloor$ as in our methods.
Then, we compute the learning rate $\eta$ of these methods based on the formula given in \citep{alacaoglu2022beyond} for \texttt{VFRBS} and \citep{alacaoglu2021stochastic} for  \texttt{VEG}, respectively. 

%%% (d)~Results and discussion.
\textbf{$\mathrm{(d)}$~Results for a different set of parameters.}
Unlike \textbf{Subsection~\ref{subsec:example1}} in the main text, we choose the parameters for these algorithms as in Subsection~\ref{apdx:subsec:detail_numexps}.
The $6$ algorithms are run on 2 experiments.
The first experiment is with $n=5000$ and $p_1 = p_2 =50$, while the second one is with $n=10000$ and $p_1 = p_2 = 100$.
These experiments are run 10 times, corresponding to $10$ problem instances, and the average results are reported in Figure~\ref{fig:bilinear_minimax_experiment1} in terms of the relative operator norm $\norms{Gx^k}/\norms{Gx^0}$ against the number of epochs.

\begin{figure}[ht]
%\vspace{-1ex}
\centering
\includegraphics[width=0.9\textwidth]{figures/unconstr_minimax_exam1_b}
\vspace{-1ex}
\caption{
The performance of $6$ algorithms to solve the unconstrained case of \eqref{eq:minimax_exam2_supp} on $2$ experiments  (The average of $10$ runs).
}
\label{fig:bilinear_minimax_experiment1}
\vspace{-1.5ex}
\end{figure}

Clearly, under this configuration, both SVRG variants of our methods work well and significantly outperform other competitors. 
The loopless SVRG variant (\texttt{VFR-svrg}) of \eqref{eq:VROG4NE} seems to work best, while our \texttt{VFR-saga} has a similar performance as \texttt{VEG}.
We also see that \texttt{VFRBS} has a similar performance as \texttt{OG}.

To improve the performance of these competitors, especially, \texttt{VFRBS} and \texttt{VEG}, one can tune their parameters as in Subsection~\ref{apdx:subsec:detail_numexps}, where the probability $\mathbf{p}$ of updating the snapshot point $w^k$ is increased.
However, with such a choice of $\mathbf{p}$, its value is often greater or equal to $0.5$, making these methods to be closed to deterministic variants.
Hence, their theoretical complexity bounds are no longer improved over the deterministic counterparts.

%%%% (e)~Small mini-batch experiments. 
%\textbf{$\mathrm{(e)}$~Small mini-batch experiments.}
%To see the effect of mini-batch sizes on the performance of the six algorithms, we choose $b = \lfloor 0.01n\rfloor$, corresponding to $b = 50$ if $n=5000$ and $b = 100$ if $n=10000$.
%In this experiment, we still choose $\mathbf{p} = \frac{1}{n^{1/3}}$ for all methods, but we set the learning rate for our methods to be $\eta = \frac{1}{\sqrt{b}L}$, while keeping the same learning rates  for other competitors as before.
%Our results are shown in Figure~\ref{fig:bilinear_minimax_small_batch}.
%
%\begin{figure}[ht]
%\vspace{-1ex}
%\centering
%\includegraphics[width=0.9\textwidth]{figures/unconstr_minimax_exam2_mb001}
%\vspace{-1ex}
%\caption{
%The performance of $6$ algorithms to solve \eqref{eq:minimax_exam2_supp} using a small mini-batch.
%}
%\label{fig:bilinear_minimax_small_batch}
%\vspace{-2ex}
%\end{figure}
%
%We observe that our loopless SVRG variant of \eqref{eq:VROG4NE} still performs best, while the double-loop SVRG and SAGA variants significantly slow down.
%This is due to the effect of some parameters which we set the same for all algorithms, while they require a specific attention.
%The \texttt{VEG} scheme performs quite well in both experiments, and outperforms our double-loop SVRG and SAGA variants in the first experiment.
%In the second experiment, it is comparable to our double-loop SVRG variant.

%%%%%%%%%%%%%%%%%%%%%%%%%%%%%%%%%%%%%%%%
%%% E.2.2. The constrained case
%%%%%%%%%%%%%%%%%%%%%%%%%%%%%%%%%%%%%%%%
\beforesubsubsec
\subsubsection{The constrained case}\label{subsec:exam2_constrained}
\aftersubsubsec
We conduct two more experiments for the constrained case of \eqref{eq:minimax_exam2_supp} as in the main text when $u \in \Delta_{p_1}$ and $v \in \Delta_{p_2}$, where $\Delta_p := \sets{u \in \R^p_{+}: \sum_{i=1}^pu_i = 1}$ is the standard simplex in $\R^p$. 

We run 6 algorithms for solving the constrained case of \eqref{eq:minimax_exam2_supp}  using the same parameters as Subsection~\ref{subsec:example1}, but with larger problems.
We report the relative norm of the FBS residual $\norms{\Gc_{\eta}x^k}/\norms{\Gc_{\eta}x^0}$ against the number of epochs. 
The results are revealed in Figure~\ref{fig:bilinear_minimax_experiment2} for two datasets $(p, n) = (500, 5000)$ and $(p, n) = (300, 10000)$. 
\begin{figure}[ht]
\vspace{-0ex}
\centering
\includegraphics[width=\textwidth]{figures/constr_minimax_exam2_c}
\vspace{-2ex}
\caption{
The performance of $6$ algorithms to solve the constrained case of \eqref{eq:minimax_exam2_supp} on $2$ experiments  (The average of $10$ runs).
}
\label{fig:bilinear_minimax_experiment2}
\vspace{-2ex}
\end{figure}

With these two additional experiments, both SVRG variants of our method \eqref{eq:VROG4NI} again work well and significantly outperform other competitors. 
The loopless SVRG variant (\texttt{VFR-svrg}) of \eqref{eq:VROG4NI} tends to work best, while our \texttt{VFR-saga} has a relatively similar performance as \texttt{VEG}.
We also see that \texttt{VFRBS}  has a similar performance trend as \texttt{OG}, but works better.



%%%%%%%%%%%%%%%%%%%%%%%%%%%%%%%%%%%%%%%%%%%%%%%%%%%%
%%% E5. Regularized Logistic Regression with Ambiguous Features}
%%%%%%%%%%%%%%%%%%%%%%%%%%%%%%%%%%%%%%%%%%%%%%%%%%%%
\beforesubsec
\subsection{The $\ell_1$-Regularized Logistic Regression with Ambiguous Features}\label{apdx:subsec:logistic_example}
\aftersubsec
This Supp. Doc. provides the details of \textbf{Subsection}~\ref{subsec:example2} in Section~\ref{sec:NumExps} in the main text.

$\mathrm{(a)}$~\textbf{Model.}
We consider a standard regularized logistic regression model associated with a given dataset $\sets{(\hat{X}_i, y_i)}_{i=1}^N$, where $\hat{X}_i$ is an i.i.d. sample of a feature vector and $y_i \in \{0, 1\}$ is the associated label of $\hat{X}_i$.
Unfortunately, $\hat{X}_i$ is ambiguous, i.e. it belongs to one of $m$ possible examples  $\{X_{ij} \}_{j=1}^m$.
Since we do not know $\hat{X}_i$ to evaluate the loss, we consider the worst-case loss $f_i(w) := \max_{1\leq j \leq m}\ell( \iprods{X_{ij}, w}, y_i)$ computed from $m$ examples, where $\ell(\tau, s) := \log(1 + \exp(\tau)) - s\tau$ is the standard logistic loss.

Using the fact that $\max_{1\leq j\leq m}\ell_j(\cdot) = \max_{z \in\Delta_m}\sum_{j=1}^m z_j \ell_j(\cdot)$, where $\Delta_m$ is  the standard simplex in $\R^m$, we can model this regularized logistic regression into the following minimax problem:
\begin{equation}\label{eq:logistic_reg_exam_apdx}
\begin{array}{lcl}
{\displaystyle\min_{w \in\R^d}} {\displaystyle\max_{z \in \R^m }} \Big\{ \mathcal{L}(w, z) := \frac{1}{N} \sum_{i=1}^{N} \sum_{j=1}^m z_j  \ell ( \iprods{X_{ij}, w}, y_i ) + \tau R(w) - \delta_{\Delta_m}(z) \Big\},
\end{array}
\end{equation}
where  $\ell(\tau, s) := \log(1 + \exp(\tau)) - s\tau$ is the standard logistic loss,  $R(w) := \norms{w}_1$ is an $\ell_1$-norm regularizer, $\tau > 0$ is a regularization parameter, and $\delta_{\Delta_m}$ is the indicator of $\Delta_m$ that handles the constraint $z \in \Delta_m$.
This problem is exactly the one stated in \eqref{eq:logistic_reg_exam} of the main text.

First, let us denote $x := [w; z] \in \R^p$ as the concatenation of $w$ and $z$ with $p = d+m$, and
\begin{equation*}
\begin{array}{lcl}
G_ix & := &  \begin{bmatrix} \sum_{j=1}^mz_j \ell'(\iprods{X_{ij},w}, y_i) X_{ij} \\ -\ell(\iprods{X_{i1}, w}, y_i)\\ 
\cdots\\
-\ell(\iprods{X_{im}, w}, y_i) \end{bmatrix} \quad \text{and}  \quad Tx :=  \begin{bmatrix} \tau \partial{R}(w)\\ \partial{\delta_{\Delta_m}}(z)\end{bmatrix}, 
\end{array}
\end{equation*}
where $\ell'(\tau, s) = \frac{\exp(\tau)}{1+\exp(\tau)} - s$.
Then, the optimality condition of \eqref{eq:logistic_reg_exam_apdx} can be written as \eqref{eq:NI}: $0 \in Gx + Tx$, where $Gx := \frac{1}{n} \sum_{i=1}^nG_ix$.

$\mathrm{(b)}$~\textbf{Input data.}
We test our algorithms and their competitors on two real datasets: \texttt{a9a} (134 features and 3561 samples) and \texttt{w8a} (311 features and 45546 samples)  downloaded from \texttt{LIBSVM} \citep{CC01a}.
For a given nominal dataset $\sets{(\hat{X}_i, y_i)}_{i=1}^n$, we first normalize the feature vector $\hat{X}_i$ such that its column norm is one,  and then add a column of  all ones to address the bias term.
To generate ambiguous features, we take the nominal feature vector $\hat{X}_i$ and add a random noise generated from a normal distribution of zero mean and variance of $\sigma = 0.5$.
In our test, we choose $\tau := 10^{-3}$ and $m := 10$ for all the experiments.



$\mathrm{(c)}$~\textbf{Algorithms.}
As before, we implement $3$ variants of our method \eqref{eq:VROG4NI}: \texttt{VFR-svrg}, \texttt{LVFR-svrg}, and \texttt{VFR-saga} to solve \eqref{eq:logistic_reg_exam_apdx}.
We also compare them with \texttt{OG}, \texttt{VFRBS}, and \texttt{VEG}.
We choose $x^0 := 0.5\cdot\texttt{ones}(p)$ in all experiments. 
We run all the algorithms for $100$ epochs and report the relative FBS residual norm $\norms{\Gc_{\eta}x^k}/\norms{\Gc_{\eta}x^0}$ against the epochs.


$\mathrm{(d)}$~\textbf{Parameters.}
Since it is very difficult to estimate the Lipschitz constant $L$ of $G$, we are unable to set a correct learning rate $\eta$ in the underlying algorithms.
We instead compute an estimation $\hat{L} := \norms{\hat{X}}$, and then set $\eta := \frac{\omega}{L}$, by tuning $\omega$ for each algorithm.
More specifically, after tuning, we obtain the following configuration.
\begin{compactitem}
\item For the three variants of \eqref{eq:VROG4NI}: \texttt{VFR-svrg}, \texttt{LVFR-svrg}, and \texttt{VFR-saga},  we set $\eta = \frac{25}{\hat{L}}$ for \texttt{a9a} and $\eta = \frac{50}{\hat{L}}$ for \texttt{w8a}.
\item For \texttt{OG}, we set $\eta = \frac{50}{\hat{L}}$ for \texttt{a9a} and $\eta = \frac{100}{\hat{L}}$ for \texttt{w8a}.
\item For \texttt{VFRBS}, we choose $\eta = \frac{47.5(1 - \sqrt{1 - \mbf{p}})}{2\hat{L}}$ for \texttt{a9a} and $\eta = \frac{95(1 - \sqrt{1 - \mbf{p}})}{2\hat{L}}$ for \texttt{w8a}.
\item For \texttt{VEG}, we select $\eta = \frac{47.5\sqrt{1-\alpha}}{\hat{L}}$ for \texttt{a9a} and $\eta = \frac{95\sqrt{1-\alpha}}{\hat{L}}$ for \texttt{w8a} with $\alpha := 1 - \mbf{p}$.
\end{compactitem}
We still choose the mini-batch size $b$ and the probability $\mbf{p}$ of updating the snapshot point $w^k$ in SVRG variants as $b = \lfloor 0.5 n^{2/3}\rfloor$ and $\mbf{p} = n^{-1/3}$, respectively  for all the algorithms.

We conduct two more experiments using the well-known MNIST dataset ($n=70000$ and $p = 780$) where we want to classify the even and odd numbers into two different classes, respectively.
We use the same parameter selection as in the experiment with the \texttt{a9a} dataset.
In the first experiment, we choose $m = 10$ and the variance of noise $\sigma = 0.5$.
In the second experiment, we choose $m=20$ and the variance of noise $\sigma = 0.25$.
We only run 5 algorithms and leave out the \texttt{VFRBS} method since we have not managed to find the parameters that make it work stably.
The result of this experiment is shown in Figure~\ref{fig:logistic_regression_exam2_b}.

\begin{figure}[ht]
\vspace{-1ex}
\begin{center}
\includegraphics[width=\textwidth]{figures/logistic_reg_exam2_b}
\vspace{-1.5ex}
\caption{Comparison of $6$ algorithms to solve \eqref{eq:logistic_reg_exam} on the real dataset: \texttt{MNIST}.}
\label{fig:logistic_regression_exam2_b}
\end{center}
\vspace{-2ex}
\end{figure}

We can see from Figure~\ref{fig:logistic_regression_exam2_b} again that three variants \texttt{VFR-svrg}, \texttt{LVFR-svrg}, and \texttt{VFR-saga} have similar performance and are much better than their two competitors.
Here, \texttt{VEG} is better than \texttt{OG}, but both methods are slower than ours.

%%%%%%%%%%%%%%%%%%%%%%%%%%%%%%%%%%%%%%%%%%%%%%%%%%%%%%%
%%% End of the supplementary document
%%%%%%%%%%%%%%%%%%%%%%%%%%%%%%%%%%%%%%%%%%%%%%%%%%%%%%%


%\section{You \emph{can} have an appendix here.}
%
%You can have as much text here as you want. The main body must be at most $8$ pages long.
%For the final version, one more page can be added.
%If you want, you can use an appendix like this one.  
%
%The $\mathtt{\backslash onecolumn}$ command above can be kept in place if you prefer a one-column appendix, or can be removed if you prefer a two-column appendix.  Apart from this possible change, the style (font size, spacing, margins, page numbering, etc.) should be kept the same as the main body.
%%%%%%%%%%%%%%%%%%%%%%%%%%%%%%%%%%%%%%%%%%%%%%%%%%%%%%%%%%%%%%%%%%%%%%%%%%%%%%%
%%%%%%%%%%%%%%%%%%%%%%%%%%%%%%%%%%%%%%%%%%%%%%%%%%%%%%%%%%%%%%%%%%%%%%%%%%%%%%%

%\bibliography{/Users/quoctd/Dropbox/E-Books/tran_bibtex_new.bib}
\bibliographystyle{icml2025}
\begin{thebibliography}{57}
\providecommand{\natexlab}[1]{#1}
\providecommand{\url}[1]{\texttt{#1}}
\expandafter\ifx\csname urlstyle\endcsname\relax
  \providecommand{\doi}[1]{doi: #1}\else
  \providecommand{\doi}{doi: \begingroup \urlstyle{rm}\Url}\fi

\bibitem[Adly \& Attouch(2021)Adly and Attouch]{adly2021first}
Adly, S. and Attouch, H.
\newblock First-order inertial algorithms involving dry friction damping.
\newblock \emph{Math. Program.}, pp.\  1--41, 2021.

\bibitem[Alacaoglu \& Malitsky(2022)Alacaoglu and
  Malitsky]{alacaoglu2021stochastic}
Alacaoglu, A. and Malitsky, Y.
\newblock Stochastic variance reduction for variational inequality methods.
\newblock In \emph{Conference on Learning Theory}, pp.\  778--816. PMLR, 2022.

\bibitem[Alacaoglu et~al.(2023)Alacaoglu, B{\"o}hm, and
  Malitsky]{alacaoglu2022beyond}
Alacaoglu, A., B{\"o}hm, A., and Malitsky, Y.
\newblock Beyond the golden ratio for variational inequality algorithms.
\newblock \emph{J. Mach. Learn. Res}, 24\penalty0 (172):\penalty0 1--33, 2023.

\bibitem[Arjovsky et~al.(2017)Arjovsky, Chintala, and
  Bottou]{arjovsky2017wasserstein}
Arjovsky, M., Chintala, S., and Bottou, L.
\newblock Wasserstein generative adversarial networks.
\newblock In \emph{International Conference on Machine Learning}, pp.\
  214--223, 2017.

\bibitem[Attouch \& Cabot(2020)Attouch and Cabot]{attouch2020convergence}
Attouch, H. and Cabot, A.
\newblock Convergence of a relaxed inertial proximal algorithm for maximally
  monotone operators.
\newblock \emph{Math. Program.}, 184\penalty0 (1):\penalty0 243--287, 2020.

\bibitem[Bauschke et~al.(2020)Bauschke, Moursi, and
  Wang]{bauschke2020generalized}
Bauschke, H.~H., Moursi, W.~M., and Wang, X.
\newblock Generalized monotone operators and their averaged resolvents.
\newblock \emph{Math. Program.}, pp.\  1--20, 2020.

\bibitem[Beznosikov et~al.(2023)Beznosikov, Gorbunov, Berard, and
  Loizou]{beznosikov2023stochastic}
Beznosikov, A., Gorbunov, E., Berard, H., and Loizou, N.
\newblock Stochastic gradient descent-ascent: Unified theory and new efficient
  methods.
\newblock In \emph{International Conference on Artificial Intelligence and
  Statistics}, pp.\  172--235. PMLR, 2023.

\bibitem[B{\"o}hm(2022)]{bohm2022solving}
B{\"o}hm, A.
\newblock Solving nonconvex-nonconcave min-max problems exhibiting weak {M}inty
  solutions.
\newblock \emph{Transactions on Machine Learning Research}, 2022.

\bibitem[Bo{\c{t}} et~al.(2021)Bo{\c{t}}, Mertikopoulos, Staudigl, and
  Vuong]{boct2021minibatch}
Bo{\c{t}}, R.~I., Mertikopoulos, P., Staudigl, M., and Vuong, P.~T.
\newblock Minibatch forward-backward-forward methods for solving stochastic
  variational inequalities.
\newblock \emph{Stochastic Systems}, 11\penalty0 (2):\penalty0 112--139, 2021.

\bibitem[Cai et~al.(2022)Cai, Song, Guzm\'{a}n, and
  Diakonikolas]{cai2022stochastic}
Cai, X., Song, C., Guzm\'{a}n, C., and Diakonikolas, J.
\newblock A stochastic {H}alpern iteration with variance reduction for
  stochastic monotone inclusion problems.
\newblock In \emph{Proceedings of the 12th International Conference on Learning
  Representations (ICLR 2022)}, 2022.
\newblock URL \url{https://openreview.net/forum?id=BRZos-8TpCf}.

\bibitem[Cai et~al.(2024)Cai, Alacaoglu, and Diakonikolas]{cai2023variance}
Cai, X., Alacaoglu, A., and Diakonikolas, J.
\newblock Variance reduced halpern iteration for finite-sum monotone
  inclusions.
\newblock In \emph{The 12th International Conference on Learning
  Representations (ICLR)}, pp.\  1--33, 2024.

\bibitem[Chakrabarti et~al.(2024)Chakrabarti, Diakonikolas, and
  Kroer]{chakrabarti2024block}
Chakrabarti, D., Diakonikolas, J., and Kroer, C.
\newblock Block-coordinate methods and restarting for solving extensive-form
  games.
\newblock \emph{Advances in Neural Information Processing Systems}, 36, 2024.

\bibitem[Chang \& Lin(2011)Chang and Lin]{CC01a}
Chang, C.-C. and Lin, C.-J.
\newblock {LIBSVM}: {A} library for {S}upport {V}ector {M}achines.
\newblock \emph{ACM Transactions on Intelligent Systems and Technology},
  2:\penalty0 27:1--27:27, 2011.

\bibitem[Combettes \& Eckstein(2018)Combettes and
  Eckstein]{combettes2018asynchronous}
Combettes, P.~L. and Eckstein, J.
\newblock Asynchronous block-iterative primal-dual decomposition methods for
  monotone inclusions.
\newblock \emph{Math. Program.}, 168\penalty0 (1):\penalty0 645--672, 2018.

\bibitem[Combettes \& Pesquet(2015)Combettes and
  Pesquet]{combettes2015stochastic}
Combettes, P.~L. and Pesquet, J.-C.
\newblock Stochastic quasi-{F}ej{\'e}r block-coordinate fixed point iterations
  with random sweeping.
\newblock \emph{SIAM J. Optim.}, 25\penalty0 (2):\penalty0 1221--1248, 2015.

\bibitem[Cui \& Shanbhag(2021)Cui and Shanbhag]{cui2021analysis}
Cui, S. and Shanbhag, U.
\newblock On the analysis of variance-reduced and randomized projection
  variants of single projection schemes for monotone stochastic variational
  inequality problems.
\newblock \emph{Set-Valued and Variational Analysis}, 29\penalty0 (2):\penalty0
  453--499, 2021.

\bibitem[Cutkosky \& Orabona(2019)Cutkosky and Orabona]{Cutkosky2019}
Cutkosky, A. and Orabona, F.
\newblock Momentum-based variance reduction in non-convex {SGD}.
\newblock In \emph{Advances in Neural Information Processing Systems}, pp.\
  15210--15219, 2019.

\bibitem[Daskalakis et~al.(2018)Daskalakis, Ilyas, Syrgkanis, and
  Zeng]{daskalakis2018training}
Daskalakis, C., Ilyas, A., Syrgkanis, V., and Zeng, H.
\newblock Training {GANs} with {O}ptimism.
\newblock In \emph{International Conference on Learning Representations (ICLR
  2018)}, 2018.

\bibitem[Defazio et~al.(2014)Defazio, Bach, and Lacoste-Julien]{Defazio2014}
Defazio, A., Bach, F., and Lacoste-Julien, S.
\newblock {SAGA}: {A} fast incremental gradient method with support for
  non-strongly convex composite objectives.
\newblock In \emph{Advances in Neural Information Processing Systems (NIPS)},
  pp.\  1646--1654, 2014.

\bibitem[Diakonikolas et~al.(2021)Diakonikolas, Daskalakis, and
  Jordan]{diakonikolas2021efficient}
Diakonikolas, J., Daskalakis, C., and Jordan, M.
\newblock Efficient methods for structured nonconvex-nonconcave min-max
  optimization.
\newblock In \emph{International Conference on Artificial Intelligence and
  Statistics}, pp.\  2746--2754. PMLR, 2021.

\bibitem[Driggs et~al.(2022)Driggs, Liang, and Sch{\"o}nlieb]{driggs2019bias}
Driggs, D., Liang, J., and Sch{\"o}nlieb, C.-B.
\newblock On biased stochastic gradient estimation.
\newblock \emph{Journal of Machine Learning Research}, 23\penalty0
  (24):\penalty0 1--43, 2022.

\bibitem[Gorbunov et~al.(2022{\natexlab{a}})Gorbunov, Berard, Gidel, and
  Loizou]{gorbunov2022stochastic}
Gorbunov, E., Berard, H., Gidel, G., and Loizou, N.
\newblock Stochastic extragradient: General analysis and improved rates.
\newblock In \emph{International Conference on Artificial Intelligence and
  Statistics}, pp.\  7865--7901. PMLR, 2022{\natexlab{a}}.

\bibitem[Gorbunov et~al.(2022{\natexlab{b}})Gorbunov, Loizou, and
  Gidel]{gorbunov2022extragradient}
Gorbunov, E., Loizou, N., and Gidel, G.
\newblock Extragradient method: $\mathcal{O} (1/k)$ last-iterate convergence
  for monotone variational inequalities and connections with cocoercivity.
\newblock In \emph{International Conference on Artificial Intelligence and
  Statistics}, pp.\  366--402. PMLR, 2022{\natexlab{b}}.

\bibitem[Gower et~al.(2021)Gower, Richt{\'a}rik, and Bach]{gower2021stochastic}
Gower, R.~M., Richt{\'a}rik, P., and Bach, F.
\newblock Stochastic quasi-gradient methods: {V}ariance reduction via
  {J}acobian sketching.
\newblock \emph{Math. Program.}, 188\penalty0 (1):\penalty0 135--192, 2021.

\bibitem[Hamedani et~al.(2018)Hamedani, Jalilzadeh, Aybat, and
  Shanbhag]{hamedani2018iteration}
Hamedani, E.~Y., Jalilzadeh, A., Aybat, N.~S., and Shanbhag, U.~V.
\newblock Iteration complexity of randomized primal-dual methods for
  convex-concave saddle point problems.
\newblock \emph{arXiv preprint arXiv:1806.04118}, 2018.

\bibitem[Hanzely et~al.(2018)Hanzely, Mishchenko, and
  Richt{\'a}rik]{hanzely2018sega}
Hanzely, F., Mishchenko, K., and Richt{\'a}rik, P.
\newblock {SEGA}: {V}ariance reduction via gradient sketching.
\newblock In \emph{Advances in Neural Information Processing Systems}, pp.\
  2082--2093, 2018.

\bibitem[Horv{\'a}th et~al.(2023)Horv{\'a}th, Kovalev, Mishchenko,
  Richt{\'a}rik, and Stich]{horvath2023stochastic}
Horv{\'a}th, S., Kovalev, D., Mishchenko, K., Richt{\'a}rik, P., and Stich, S.
\newblock Stochastic distributed learning with gradient quantization and
  double-variance reduction.
\newblock \emph{Optimization Methods and Software}, 38\penalty0 (1):\penalty0
  91--106, 2023.

\bibitem[Hsieh et~al.(2019)Hsieh, Iutzeler, Malick, and
  Mertikopoulos]{hsieh2019convergence}
Hsieh, Y.-G., Iutzeler, F., Malick, J., and Mertikopoulos, P.
\newblock On the convergence of single-call stochastic extra-gradient methods.
\newblock In \emph{Advances in Neural Information Processing Systems}, pp.\
  6938--6948, 2019.

\bibitem[Iusem et~al.(2017)Iusem, Jofr{\'e}, Oliveira, and
  Thompson]{iusem2017extragradient}
Iusem, A.~N., Jofr{\'e}, A., Oliveira, R.~I., and Thompson, P.
\newblock Extragradient method with variance reduction for stochastic
  variational inequalities.
\newblock \emph{SIAM J. Optim.}, 27\penalty0 (2):\penalty0 686--724, 2017.

\bibitem[Johnson \& Zhang(2013)Johnson and Zhang]{SVRG}
Johnson, R. and Zhang, T.
\newblock Accelerating stochastic gradient descent using predictive variance
  reduction.
\newblock In \emph{NIPS}, pp.\  315--323, 2013.

\bibitem[Juditsky et~al.(2011)Juditsky, Nemirovski, and
  Tauvel]{juditsky2011solving}
Juditsky, A., Nemirovski, A., and Tauvel, C.
\newblock Solving variational inequalities with stochastic mirror-prox
  algorithm.
\newblock \emph{Stochastic Systems}, 1\penalty0 (1):\penalty0 17--58, 2011.

\bibitem[Kannan \& Shanbhag(2019)Kannan and Shanbhag]{kannan2019optimal}
Kannan, A. and Shanbhag, U.~V.
\newblock Optimal stochastic extragradient schemes for pseudomonotone
  stochastic variational inequality problems and their variants.
\newblock \emph{Comput. Optim. Appl.}, 74\penalty0 (3):\penalty0 779--820,
  2019.

\bibitem[Konnov(2001)]{Konnov2001}
Konnov, I.
\newblock \emph{Combined relaxation methods for variational inequalities.}
\newblock Springer-Verlag, 2001.

\bibitem[Kotsalis et~al.(2022)Kotsalis, Lan, and Li]{kotsalis2022simple}
Kotsalis, G., Lan, G., and Li, T.
\newblock Simple and optimal methods for stochastic variational inequalities,
  i: operator extrapolation.
\newblock \emph{SIAM J. Optim.}, 32\penalty0 (3):\penalty0 2041--2073, 2022.

\bibitem[Li et~al.(2022)Li, Yu, Loizou, Gidel, Ma, Roux, and
  Jordan]{li2022convergence}
Li, C.~J., Yu, Y., Loizou, N., Gidel, G., Ma, Y., Roux, N.~L., and Jordan, M.
\newblock On the convergence of stochastic extragradient for bilinear games
  using restarted iteration averaging.
\newblock In \emph{International Conference on Artificial Intelligence and
  Statistics}, pp.\  9793--9826. PMLR, 2022.

\bibitem[Loizou et~al.(2021)Loizou, Berard, Gidel, Mitliagkas, and
  Lacoste-Julien]{loizou2021stochastic}
Loizou, N., Berard, H., Gidel, G., Mitliagkas, I., and Lacoste-Julien, S.
\newblock Stochastic gradient descent-ascent and consensus optimization for
  smooth games: {C}onvergence analysis under expected co-coercivity.
\newblock \emph{Advances in Neural Information Processing Systems},
  34:\penalty0 19095--19108, 2021.

\bibitem[Luo \& Tran-Dinh(2022)Luo and Tran-Dinh]{luo2022last}
Luo, Y. and Tran-Dinh, Q.
\newblock Extragradient-type methods for co-monotone root-finding problems.
\newblock \emph{(UNC-STOR Technical Report)}, 2022.

\bibitem[Malitsky \& Tam(2020)Malitsky and Tam]{malitsky2020forward}
Malitsky, Y. and Tam, M.~K.
\newblock A forward-backward splitting method for monotone inclusions without
  cocoercivity.
\newblock \emph{SIAM J. Optim.}, 30\penalty0 (2):\penalty0 1451--1472, 2020.

\bibitem[Mishchenko et~al.(2020)Mishchenko, Kovalev, Shulgin, Richt{\'a}rik,
  and Malitsky]{mishchenko2020revisiting}
Mishchenko, K., Kovalev, D., Shulgin, E., Richt{\'a}rik, P., and Malitsky, Y.
\newblock Revisiting stochastic extragradient.
\newblock In \emph{International Conference on Artificial Intelligence and
  Statistics}, pp.\  4573--4582. PMLR, 2020.

\bibitem[Noor(2003)]{noor2003extragradient}
Noor, M.~A.
\newblock Extragradient methods for pseudomonotone variational inequalities.
\newblock \emph{J. Optim. Theory Appl.}, 117\penalty0 (3):\penalty0 475--488,
  2003.

\bibitem[Noor \& Al-Said(1999)Noor and Al-Said]{noor1999wiener}
Noor, M.~A. and Al-Said, E.
\newblock {W}iener--{H}opf equations technique for quasimonotone variational
  inequalities.
\newblock \emph{J. Optim. Theory Appl.}, 103:\penalty0 705--714, 1999.

\bibitem[Peng et~al.(2016)Peng, Xu, Yan, and Yin]{peng2016arock}
Peng, Z., Xu, Y., Yan, M., and Yin, W.
\newblock {AR}ock: an algorithmic framework for asynchronous parallel
  coordinate updates.
\newblock \emph{SIAM J. Scientific Comput.}, 38\penalty0 (5):\penalty0
  2851--2879, 2016.

\bibitem[Pethick et~al.(2023)Pethick, Fercoq, Latafat, Patrinos, and
  Cevher]{pethick2023solving}
Pethick, T., Fercoq, O., Latafat, P., Patrinos, P., and Cevher, V.
\newblock Solving stochastic weak {M}inty variational inequalities without
  increasing batch size.
\newblock In \emph{Proceedings of International Conference on Learning
  Representations (ICLR)}, pp.\  1--34, 2023.

\bibitem[Pham et~al.(2020)Pham, Nguyen, Phan, and Tran-Dinh]{Pham2019}
Pham, H.~N., Nguyen, M.~L., Phan, T.~D., and Tran-Dinh, Q.
\newblock {ProxSARAH}: {A}n efficient algorithmic framework for stochastic
  composite nonconvex optimization.
\newblock \emph{J. Mach. Learn. Res.}, 21:\penalty0 1--48, 2020.

\bibitem[Robbins \& Monro(1951)Robbins and Monro]{RM1951}
Robbins, H. and Monro, S.
\newblock A stochastic approximation method.
\newblock \emph{The Annals of Mathematical Statistics}, 22\penalty0
  (3):\penalty0 400--407, 1951.

\bibitem[Rockafellar \& Wets(1997)Rockafellar and Wets]{Rockafellar1997}
Rockafellar, R. and Wets, R. J.-B.
\newblock \emph{{V}ariational {A}nalysis}.
\newblock Springer-Verlag, 1997.

\bibitem[Song \& Diakonikolas(2023)Song and Diakonikolas]{song2023cyclic}
Song, C. and Diakonikolas, J.
\newblock Cyclic coordinate dual averaging with extrapolation.
\newblock \emph{SIAM J. Optim.}, 33\penalty0 (4):\penalty0 2935--2961, 2023.

\bibitem[Song et~al.(2020)Song, Zhou, Zhou, Jiang, and Ma]{song2020optimistic}
Song, C., Zhou, Z., Zhou, Y., Jiang, Y., and Ma, Y.
\newblock Optimistic dual extrapolation for coherent non-monotone variational
  inequalities.
\newblock \emph{Advances in Neural Information Processing Systems},
  33:\penalty0 14303--14314, 2020.

\bibitem[Tran-Dinh(2023)]{tran2023sublinear}
Tran-Dinh, Q.
\newblock {S}ublinear {C}onvergence {R}ates of {E}xtragradient-{T}ype
  {M}ethods: {A} {S}urvey on {C}lassical and {R}ecent {D}evelopments.
\newblock \emph{arXiv preprint arXiv:2303.17192}, 2023.

\bibitem[Tran-Dinh \& Luo(2025)Tran-Dinh and Luo]{tran2022accelerated}
Tran-Dinh, Q. and Luo, Y.
\newblock {R}andomized {B}lock-{C}oordinate {O}ptimistic {G}radient
  {A}lgorithms for {R}oot-{F}inding {P}roblems.
\newblock \emph{Math. Oper. Res.}, in press, 2025.

\bibitem[Tran-Dinh \& Nguyen-Trung(2025{\natexlab{a}})Tran-Dinh and
  Nguyen-Trung]{TranDinh2025a}
Tran-Dinh, Q. and Nguyen-Trung, N.
\newblock {V}ariance-{R}educed {A}ccelerated {F}ixed-{P}oint-{B}ased
  {M}ethodsfor {G}eneralized {E}quations: {B}etter {C}onvergence {G}uarantees.
\newblock \emph{Tech. Report -- UNC-STOR}, 2025{\natexlab{a}}.

\bibitem[Tran-Dinh \& Nguyen-Trung(2025{\natexlab{b}})Tran-Dinh and
  Nguyen-Trung]{tran2025accelerated}
Tran-Dinh, Q. and Nguyen-Trung, N.
\newblock Accelerated extragradient-type methods--part 2: Generalization and
  sublinear convergence rates under co-hypomonotonicity.
\newblock \emph{arXiv preprint arXiv:2501.04585}, 2025{\natexlab{b}}.

\bibitem[Tran-Dinh et~al.(2022)Tran-Dinh, Pham, Phan, and
  Nguyen]{Tran-Dinh2019a}
Tran-Dinh, Q., Pham, N.~H., Phan, D.~T., and Nguyen, L.~M.
\newblock A hybrid stochastic optimization framework for stochastic composite
  nonconvex optimization.
\newblock \emph{Math. Program.}, 191:\penalty0 1005--1071, 2022.

\bibitem[Tu(2018)]{vuong2018weak}
Tu, V.~P.
\newblock On the weak convergence of the extragradient method for solving
  pseudo-monotone variational inequalities.
\newblock \emph{J. Optim. Theory Appl.}, 176\penalty0 (2):\penalty0 399--409,
  2018.

\bibitem[Yang et~al.(2020)Yang, Kiyavash, and He]{yang2020global}
Yang, J., Kiyavash, N., and He, N.
\newblock Global convergence and variance-reduced optimization for a class of
  nonconvex-nonconcave minimax problems.
\newblock \emph{arXiv preprint arXiv:2002.09621}, 2020.

\bibitem[Yousefian et~al.(2018)Yousefian, Nedi{\'c}, and
  Shanbhag]{yousefian2018stochastic}
Yousefian, F., Nedi{\'c}, A., and Shanbhag, U.~V.
\newblock On stochastic mirror-prox algorithms for stochastic cartesian
  variational inequalities: {R}andomized block coordinate and optimal averaging
  schemes.
\newblock \emph{Set-Valued and Variational Analysis}, 26:\penalty0 789--819,
  2018.

\bibitem[Zhou et~al.(2018)Zhou, Xu, and Gu]{zhou2018stochastic}
Zhou, D., Xu, P., and Gu, Q.
\newblock Stochastic nested variance reduction for nonconvex optimization.
\newblock In \emph{Proceedings of the 32nd International Conference on Neural
  Information Processing Systems}, pp.\  3925--3936. Curran Associates Inc.,
  2018.

\end{thebibliography}

%\bibliography{/Users/quoc/Dropbox/E-Books/tran_bibtex_new.bib}

\end{document}


% This document was modified from the file originally made available by
% Pat Langley and Andrea Danyluk for ICML-2K. This version was created
% by Iain Murray in 2018, and modified by Alexandre Bouchard in
% 2019 and 2021 and by Csaba Szepesvari, Gang Niu and Sivan Sabato in 2022.
% Modified again in 2023 and 2024 by Sivan Sabato and Jonathan Scarlett.
% Previous contributors include Dan Roy, Lise Getoor and Tobias
% Scheffer, which was slightly modified from the 2010 version by
% Thorsten Joachims & Johannes Fuernkranz, slightly modified from the
% 2009 version by Kiri Wagstaff and Sam Roweis's 2008 version, which is
% slightly modified from Prasad Tadepalli's 2007 version which is a
% lightly changed version of the previous year's version by Andrew
% Moore, which was in turn edited from those of Kristian Kersting and
% Codrina Lauth. Alex Smola contributed to the algorithmic style files.
